\documentclass[11pt,a4paper]{article}
\usepackage[margin=18mm,top=19mm,bottom=19mm,headheight=15pt,headsep=6mm]{geometry}
\usepackage{fontspec}
\usepackage{xeCJK}
\IfFontExistsTF{Times New Roman}{\setmainfont{Times New Roman}}{\setmainfont{TeX Gyre Termes}}
\IfFontExistsTF{Arial}{\setsansfont{Arial}}{\setsansfont{TeX Gyre Heros}}
\IfFontExistsTF{Consolas}{\setmonofont[Scale=0.94]{Consolas}}{\setmonofont[Scale=0.94]{Latin Modern Mono}}
\IfFontExistsTF{SimSun}{\setCJKmainfont[AutoFakeBold=1.8]{SimSun}}{\setCJKmainfont{FandolSong-Regular}}
\IfFontExistsTF{Microsoft YaHei}{\setCJKsansfont{Microsoft YaHei}}{\setCJKsansfont{FandolHei-Regular}}
\IfFontExistsTF{SimSun}{\setCJKmonofont{SimSun}}{\setCJKmonofont{FandolSong-Regular}}
\usepackage{amsmath,amssymb}
\usepackage{booktabs,array,longtable}
\usepackage{fvextra}
\usepackage{graphicx}
\usepackage{tikz,pgfplots}
\usetikzlibrary{arrows.meta,positioning,calc,fit,backgrounds,patterns,decorations.pathreplacing,matrix}
\pgfplotsset{compat=1.18}
\usepackage{caption}
\usepackage{titlesec}
\usepackage{enumitem}
\usepackage{needspace}
\usepackage{placeins}
\usepackage{setspace}
\usepackage{fancyhdr}
\usepackage{xurl}
\usepackage[unicode,hidelinks,bookmarksnumbered=true,hypertexnames=false,pdfencoding=auto]{hyperref}
\usepackage{bookmark}
\hypersetup{pdftitle={字节级 BPE 分析与多语言抽样设计},pdfsubject={BPE 与多语言分词器训练语料比例},pdfauthor={}}
\urlstyle{same}
\setlength{\parindent}{0pt}
\setlength{\parskip}{0.48em}
\setstretch{1.20}
\setlength{\emergencystretch}{2em}
\raggedbottom
\setcounter{secnumdepth}{3}
\setcounter{tocdepth}{2}
\renewcommand{\contentsname}{目录}
\renewcommand{\figurename}{图}
\renewcommand{\tablename}{表}
\titleformat{\section}{\sffamily\large\bfseries}{\thesection}{0.65em}{}
\titleformat{\subsection}{\sffamily\normalsize\bfseries}{\thesubsection}{0.65em}{}
\titleformat{\subsubsection}{\sffamily\normalsize\bfseries}{\thesubsubsection}{0.65em}{}
\titlespacing*{\section}{0pt}{1.5em}{0.7em}
\titlespacing*{\subsection}{0pt}{1.1em}{0.45em}
\titlespacing*{\subsubsection}{0pt}{0.9em}{0.35em}
\setlist[itemize]{leftmargin=1.7em,itemsep=0.3em,topsep=0.4em,parsep=0pt}
\captionsetup{font=small,labelfont=bf,labelsep=quad,skip=7pt,justification=justified,singlelinecheck=false}
\pagestyle{fancy}
\fancyhf{}
\fancyhead[L]{\small Tokenization Analysis}
\fancyhead[R]{\small 原理、源码与实验记录}
\fancyfoot[C]{\small\thepage}
\renewcommand{\headrulewidth}{0.3pt}
\renewcommand{\footrulewidth}{0pt}
\setlength{\tabcolsep}{4pt}
\setlength{\LTpre}{0.45em}
\setlength{\LTpost}{0.65em}
\newlength{\BpeDocTableWidth}
\newcommand{\BpeDocCodeCell}[1]{{\ttfamily\fontsize{7.8}{10.3}\selectfont #1}}
\newcommand{\SourceLineNumber}[1]{{\ttfamily\fontsize{8.2}{10.3}\selectfont #1}}
\newcommand{\BpeDocCodeInline}[1]{{\ttfamily #1}}
\newcommand{\BpeDocSection}[2]{\section{\texorpdfstring{#1}{#2}}}
\newcommand{\BpeDocSubsection}[2]{\subsection{\texorpdfstring{#1}{#2}}}
\newcommand{\BpeDocSubsubsection}[2]{\subsubsection{\texorpdfstring{#1}{#2}}}
\begin{document}


\begin{center}
{\sffamily\fontsize{19}{26}\selectfont\bfseries 字节级 BPE 分析与多语言抽样设计\par}
\vspace{3mm}
{\large Byte-Level BPE Analysis and Multilingual Sampling Design\par}
\vspace{2mm}
{\small Technical Research Report\par}
\end{center}
\renewcommand{\abstractname}{摘要}
\begin{abstract}
本文通过源码与运行记录，跟踪 MinBPE 如何学习和应用合并规则。训练 abababab 时，第一轮得到四个连续的 ab 词元。第二轮统计到三个相邻位置对，从左往右合并时每个词元只能用一次，因此执行两次替换。另一个例子说明，编码新文本时按训练得到的规则优先级合并。在选定的八个普通文本输入上，MinBPE 包装器与 cl100k\_base 逐个编号一致，解码均还原原文。验证程序的十七组行为检查全部通过。德语／西语部分在假设库存上完成了四项配额计算检查，并给出在真实语料上比较不同抽样比例的后续实验方案。
\end{abstract}
\noindent\textbf{关键词：}字节级 BPE；编码一致性；UTF-8；多语言抽样；可复现性。
\par\vspace{2mm}
\pagenumbering{arabic}


% Content source: reports/content/research_overview.json

\Needspace{9\baselineskip}
\BpeDocSection{研究问题与分析框架}{研究问题与分析框架}


分词器将文本映射为离散编号，其合并规则、边界约束和训练语料分布共同决定输入序列的表示方式。准确分析这些环节，需要区分算法定义、具体实现和实证检查的证据范围。本项目以固定版本的 MinBPE 与 \BpeDocCodeInline{cl100k\char95\relax\allowbreak{}base} 为对象，结合受控合并轨迹、源码审计和差分编码核对，建立能够逐项复查的机制证据；在此基础上，提出德语／西语分词器语料抽样的控制实验方案。


分析遵循三个层次：首先用规模有限、结果可完全追踪的输入检查 BPE 的局部操作；其次在不同文字与格式类型上比较已有编码的 \BpeDocCodeInline{token} ID 和完整往返结果；最后把文献中的语言抽样思想转化为明确的计量单位、数据预算和语言内评价条件。三个层次对应不同的推断目标，不能仅凭前一层的检查通过，就宣布后一层的语言质量或模型能力得到改善。


\Needspace{9\baselineskip}
\BpeDocSubsection{研究问题与可检验命题}{研究问题与可检验命题}


\begingroup\fontsize{9.3}{12.7}\selectfont\setstretch{1}
\setlength{\BpeDocTableWidth}{\dimexpr\linewidth-6\tabcolsep\relax}
\renewcommand{\arraystretch}{1.15}
\par\addvspace{0.45em}
\noindent\begin{minipage}{\linewidth}
\begin{tabular}{>{\raggedright\arraybackslash}p{0.280000\BpeDocTableWidth}>{\raggedright\arraybackslash}p{0.380000\BpeDocTableWidth}>{\raggedright\arraybackslash}p{0.340000\BpeDocTableWidth}}
\toprule
\textbf{问题} & \textbf{操作性检查} & \textbf{证据能够回答的范围} \\

\midrule
RQ1：BPE 训练与编码中的选择依据有何不同？ & 跟踪相邻对频率、实际替换次数和已学合并优先级；构造两种选择规则产生不同编号的输入。 & 定位重叠计数与不重叠替换的差别，并检查新文本编码是否遵循既有规则。 \\[3pt]
RQ2：具体编码包装器在什么范围内复现参考实现？ & 在固定编码、固定依赖和同一输入上比较完整 \BpeDocCodeInline{token} ID 列表、字节拼接和文本往返。 & 确认所覆盖普通文本案例的一致性，识别正则边界、UTF-8 片段与字节编号映射的作用。 \\[3pt]
RQ3：如何使多语言语料比例成为可评价的设计变量？ & 固定词表、字符预算、文本清洗与数据划分；先计算候选配额，再规定独立验证与测试指标。 & 给出可执行的比例选择程序；真实语料上的最优候选仍需按该程序测量。 \\[3pt]
\bottomrule
\end{tabular}
\end{minipage}\par\addvspace{0.65em}
\endgroup


对 RQ1，检验重点是可区分的行为命题：统计阶段允许相邻窗口共享位置，替换阶段按从左到右的不重叠配对消耗位置；训练时选择当前最高频相邻对，而编码时选择已学规则中优先级最高的可用合并。对 RQ2，判据采用精确编号相等与完整文本还原，而不是仅比较 \BpeDocCodeInline{token} 总数。对 RQ3，判据首先要求比较条件一致，再讨论语言内序列长度的变化。


\Needspace{9\baselineskip}
\BpeDocSubsection{分析设计与证据来源}{分析设计与证据来源}


算法分析固定在 MinBPE 提交 \BpeDocCodeInline{1acefe89\allowbreak{}412b2024\allowbreak{}5db5a22d\allowbreak{}2a02001e\allowbreak{}547dc602}。上游文件保留原始实现、许可与哈希，项目中的跟踪和核对代码调用这些实现并保存中间状态。因此，可以将报告中的每个关键步骤追溯到具体函数、执行行和结构化结果，而不必依赖图形外观反推算法行为。


\begin{itemize}

\item 受控机制检查：训练微型词表，记录初始字节、每轮全部相邻对计数、所选规则、实际替换次数、词表内容与最终编号。将跟踪结果与 \BpeDocCodeInline{BasicTokenizer} 的训练、编码及解码输出核对。

\item 差分编码检查：分别初始化 MinBPE \BpeDocCodeInline{GPT4Tokenizer} 与 tiktoken 的 \BpeDocCodeInline{cl100k\char95\relax\allowbreak{}base}，对英文、德语、西班牙语、中文、数字、换行、Emoji 和空输入执行同一组核对。输入范围及逐 \BpeDocCodeInline{token} 字节内容保存在结果文件中。

\item 抽样设计检查：独立计算平滑语言份额与整数配额，并检查总预算约束、概率归一化和零库存边界。语料库存为明确标注的假设值；该部分的执行证据是配额运算，而非语料训练或性能评估。

\end{itemize}



\begingroup\fontsize{9.3}{12.7}\selectfont\setstretch{1}
\setlength{\BpeDocTableWidth}{\dimexpr\linewidth-6\tabcolsep\relax}
\renewcommand{\arraystretch}{1.15}
\par\addvspace{0.45em}
\noindent\begin{minipage}{\linewidth}
\begin{tabular}{>{\raggedright\arraybackslash}p{0.230000\BpeDocTableWidth}>{\raggedright\arraybackslash}p{0.400000\BpeDocTableWidth}>{\raggedright\arraybackslash}p{0.370000\BpeDocTableWidth}}
\toprule
\textbf{记录单位} & \textbf{具体含义} & \textbf{解释约束} \\

\midrule
相邻对频率／替换次数 & 前者计数所有相邻位置窗口，后者记录一次实际合并消耗的非重叠配对。 & 两个数值描述不同操作，存在重复元素时不能直接互换。 \\[3pt]
\BpeDocCodeInline{token} ID／UTF-8 字节 & ID 是编码词表中的索引，字节是该索引对应的原始内容。 & ID 数值本身不表示字符语义；一个 \BpeDocCodeInline{token} 的字节可以不是完整字符。 \\[3pt]
Unicode 码点／\BpeDocCodeInline{token} 数 & 前者作为一致清洗规则下的文本预算单位，后者作为候选编码的输出长度。 & 语言比例以输入文本计量，避免因候选分词器改变而同时改变输入预算。 \\[3pt]
检查项／输入案例 & 检查项可包含多条断言；输入案例是一次具体的文本及其编码记录。 & 17 项行为检查与 8 类普通文本是不同计数，不视为独立数据集数量。 \\[3pt]
\bottomrule
\end{tabular}
\end{minipage}\par\addvspace{0.65em}
\endgroup


直接证据位于 \BpeDocCodeInline{results/\allowbreak{}bpe\char95\relax\allowbreak{}walkthrough.\allowbreak{}json}、\BpeDocCodeInline{results/\allowbreak{}gpt4\char95\relax\allowbreak{}check.\allowbreak{}json}、\BpeDocCodeInline{results/\allowbreak{}verification.\allowbreak{}json} 和 \BpeDocCodeInline{results/\allowbreak{}language\char95\relax\allowbreak{}mix.\allowbreak{}json}。图表数据、图源与文件哈希形成另一条核对链。正式结果章节呈现具体输入、编号和检查结论；源码审计附录补充实现细节与行号。


\Needspace{9\baselineskip}
\BpeDocSubsection{主要观察及机制解释}{主要观察及机制解释}


输入 abababab 在第一轮后变成四个 ab 词元，按从左到右编号为 1 至 4。第二轮统计时，第 1、2 个，第 2、3 个和第 3、4 个都构成相同的相邻对，共三个候选位置。替换时先合并第 1、2 个，第 2 个已经用过，接着合并第 3、4 个，得到两个 abab 词元。频次 3 计的是候选位置，替换次数 2 计的是本轮实际执行的合并。


合并优先级反例进一步把训练与推断区别转化为可检查的输出差异。同一新输入上，若错误地优先合并当前更高频的相邻对，会得到不同的编号序列。已有规则的先后顺序因而是编码规范的一部分，而不只是训练过程的记录。仅在新文本上重新统计频率，不能替代已训练的合并表。


编码对照中，八个给定输入的完整编号与参考编码逐项一致，拼接各 \BpeDocCodeInline{token} 的字节后也能还原原文。结果章节列出全部输入，包括空字符串。检查还展示了两个实现细节：原始字节值与单字节 \BpeDocCodeInline{token} ID 存在重排关系；某些 \BpeDocCodeInline{token} 只含一个 UTF-8 字符的部分字节。


对多语言抽样，配额计算确认了平滑指数能够在库存比例与均衡份额之间生成可复算的候选设置。但配额正确只解决了输入分配问题，不能确定候选分词器在独立文本上的切分效率。后续评估需要固定语料质量、领域、词表和文本预算，再分别比较德语与西语相对同一基线的变化。相关论文用于支持这些控制条件与诊断思路，其原有任务范围和适用边界在多语言设计章节单独讨论。


\Needspace{9\baselineskip}
\BpeDocSubsection{复现与推断边界}{复现与推断边界}


现有算法轨迹和编码对照是确定性计算，复现对象由输入、固定源码、依赖版本和参考编码共同定义。保存中间序列与原始字节，使差异可以定位到具体步骤；保留完整结果而不只保存汇总数字，使后续扩展能够重复既有检查并记录新增覆盖范围。AI 辅助开发和审查记录与上游归属单独保留，以说明材料的产生过程。


当前证据支持所记录配置上的机制解释、有限输入的一致性和候选配额运算。输入覆盖仍然有限；普通文本一致性不自动包含所有特殊 \BpeDocCodeInline{token} 配置，微型词表行为也不能直接外推为大型分词器的质量收益。真实德语／西语语料比例实验与下游语言模型评估属于尚待执行的研究设计。报告将这些层次分别呈现，使已执行验证与后续实证目标保持清楚对应。


\FloatBarrier\par\bigskip

% Content source: reports/content/execution_evidence.json

\Needspace{9\baselineskip}
\BpeDocSection{编码验证结果与机制证据}{编码验证结果与机制证据}


以下结果来自仓库保留的实际运行记录，可用 \BpeDocCodeInline{code/\allowbreak{}verify.\allowbreak{}py} 与 results 中的原始编号、字节和检查项核对。受控输入用于分析合并机制；普通文本用于比较两个编码实现。所有结果均针对所记录的配置与输入范围。


\begingroup\fontsize{9.3}{12.7}\selectfont\setstretch{1}
\setlength{\BpeDocTableWidth}{\dimexpr\linewidth-4\tabcolsep\relax}
\renewcommand{\arraystretch}{1.15}
\par\addvspace{0.45em}
\noindent\begin{minipage}{\linewidth}
\begin{tabular}{>{\raggedright\arraybackslash}p{0.240000\BpeDocTableWidth}>{\raggedright\arraybackslash}p{0.760000\BpeDocTableWidth}}
\toprule
\textbf{环境项} & \textbf{实际记录} \\

\midrule
Python & 3.12.14 \\[3pt]
regex / tiktoken & 2026.9.29 / 0.14.0 \\[3pt]
MinBPE & \BpeDocCodeInline{1acefe89\allowbreak{}412b2024\allowbreak{}5db5a22d\allowbreak{}2a02001e\allowbreak{}547dc602} \\[3pt]
随机性 & BPE 示例与对照均为确定性计算，无随机种子 \\[3pt]
\bottomrule
\end{tabular}
\end{minipage}\par\addvspace{0.65em}
\endgroup


\Needspace{9\baselineskip}
\BpeDocSubsection{三轮合并与新文本编码}{三轮合并与新文本编码}


\begin{figure}[htbp]
\centering
% Source: results/bpe_walkthrough.json, training.initial_ids and training.steps.
% Teaching example: abababab, vocab_size=259. This is not cl100k_base output.
% Input fragment. Requires tikz + arrows.meta + calc; font setup belongs to main.tex.
% Fixed design box: 15.5 cm x 7.2 cm. Text content is not scaled.
\begingroup
\definecolor{tokenizerBlue}{HTML}{285D77}
\definecolor{tokenizerGrey}{HTML}{66747B}
\definecolor{tokenizerGrid}{HTML}{E1E6E9}
\begin{tikzpicture}[
  x=1cm,y=1cm,
  font=\small,
  token/.style={
    draw=tokenizerBlue,line width=.5pt,fill=tokenizerBlue!3,
    minimum height=.82cm,inner xsep=1pt,inner ysep=2pt,
    outer sep=0pt,align=center,text=tokenizerBlue
  },
  link/.style={draw=tokenizerGrey!80,line width=.45pt},
  finish/.style={link,-{Stealth[length=1.5mm,width=1.05mm]}},
  side/.style={font=\footnotesize,text=tokenizerGrey},
  id/.style={font=\footnotesize,text=tokenizerGrey}
]
  \path[use as bounding box] (-1.55,-.82) rectangle (13.95,6.38);

  % Original-byte boundaries are shared by all four stages.
  \foreach \j in {0,...,8} {
    \draw[tokenizerGrid,line width=.3pt,densely dotted]
      ({1.5*\j},.41) -- ({1.5*\j},5.80);
  }
  \node[side,anchor=west] at (-1.48,6.08) {字节位置};
  \node[side,anchor=west] at (12.43,6.08) {token 数};
  \foreach \j in {1,...,8} {
    \node[side] at ({1.5*\j-.75},6.08) {$\j$};
  }

  % Widths retain the same byte-position scale: 1, 2, 4, and 8 bytes.
  \foreach \i/\letter/\tokenid in {
    0/a/97,1/b/98,2/a/97,3/b/98,4/a/97,5/b/98,6/a/97,7/b/98
  } {
    \node[token,minimum width=1.42cm] (s0-\i)
      at ({1.5*\i+.75},5.35)
      {\texttt{\letter}\\[-1pt]
       {\footnotesize\color{tokenizerGrey}\texttt{\tokenid}}};
  }
  \foreach \i in {0,...,3} {
    \node[token,minimum width=2.92cm] (s1-\i)
      at ({3*\i+1.5},3.85)
      {\texttt{ab}\\[-1pt]
       {\footnotesize\color{tokenizerGrey}\texttt{256}}};
  }
  \foreach \i in {0,1} {
    \node[token,minimum width=5.92cm] (s2-\i)
      at ({6*\i+3},2.35)
      {\texttt{abab}\\[-1pt]
       {\footnotesize\color{tokenizerGrey}\texttt{257}}};
  }
  \node[token,minimum width=11.92cm] (s3-0) at (6,.85)
    {\texttt{abababab}\\[-1pt]
     {\footnotesize\color{tokenizerGrey}\texttt{258}}};

  % A single arrowhead per parent makes each two-to-one relation explicit.
  \foreach \i in {0,...,3} {
    \pgfmathtruncatemacro{\leftchild}{2*\i}
    \pgfmathtruncatemacro{\rightchild}{2*\i+1}
    \coordinate (join1-\i) at ($(s1-\i.north)+(0,.34)$);
    \draw[link] (s0-\leftchild.south) -- (join1-\i);
    \draw[link] (s0-\rightchild.south) -- (join1-\i);
    \draw[finish] (join1-\i) -- (s1-\i.north);
  }
  \foreach \i in {0,1} {
    \pgfmathtruncatemacro{\leftchild}{2*\i}
    \pgfmathtruncatemacro{\rightchild}{2*\i+1}
    \coordinate (join2-\i) at ($(s2-\i.north)+(0,.34)$);
    \draw[link] (s1-\leftchild.south) -- (join2-\i);
    \draw[link] (s1-\rightchild.south) -- (join2-\i);
    \draw[finish] (join2-\i) -- (s2-\i.north);
  }
  \coordinate (join3) at ($(s3-0.north)+(0,.34)$);
  \draw[link] (s2-0.south) -- (join3);
  \draw[link] (s2-1.south) -- (join3);
  \draw[finish] (join3) -- (s3-0.north);

  \foreach \stage/\yy/\count in {0/5.35/8,1/3.85/4,2/2.35/2,3/.85/1} {
    \node[side,anchor=west] at (-1.32,\yy) {$t=\stage$};
    \node[side,anchor=west] at (12.43,\yy) {$\count$};
  }
  \node[side] at (6,-.34)
    {$|\mathcal V|:\quad 256\ \longrightarrow\ 257\ \longrightarrow\ 258\ \longrightarrow\ 259$};
\end{tikzpicture}
\endgroup

\caption{受控输入 abababab 的真实 BPE 合并轨迹。横向位置与块宽表示覆盖的原始字节；每轮把已有相邻 \BpeDocCodeInline{token} 合成新 \BpeDocCodeInline{token}。图中 \BpeDocCodeInline{token} 数是当前文本的长度，词表大小则是可用编号的总数，两者不同。数据来源：\BpeDocCodeInline{results/\allowbreak{}bpe\char95\relax\allowbreak{}walkthrough.\allowbreak{}json}。}
\label{fig:bpe-merges}
\end{figure}


\begingroup\fontsize{9.3}{12.7}\selectfont\setstretch{1}
\setlength{\BpeDocTableWidth}{\dimexpr\linewidth-10\tabcolsep\relax}
\renewcommand{\arraystretch}{1.15}
\par\addvspace{0.45em}
\noindent\begin{minipage}{\linewidth}
\begin{tabular}{>{\raggedright\arraybackslash}p{0.100000\BpeDocTableWidth}>{\raggedright\arraybackslash}p{0.250000\BpeDocTableWidth}>{\raggedright\arraybackslash}p{0.130000\BpeDocTableWidth}>{\raggedright\arraybackslash}p{0.130000\BpeDocTableWidth}>{\raggedright\arraybackslash}p{0.390000\BpeDocTableWidth}}
\toprule
\textbf{轮次} & \textbf{选中相邻对} & \textbf{统计次数} & \textbf{实际替换} & \textbf{合并后的 ID} \\

\midrule
1 & [97, 98] & 4 & 4 & [256, 256, 256, 256] \\[3pt]
2 & [256, 256] & 3 & 2 & [257, 257] \\[3pt]
3 & [257, 257] & 1 & 1 & [258] \\[3pt]
\bottomrule
\end{tabular}
\end{minipage}\par\addvspace{0.65em}
\endgroup


第二轮开始时四个词元均为 ab（编号 256）。三个候选窗口分别覆盖位置 (1,2)、(2,3)、(3,4)。按从左到右的规则，替换 (1,2) 后直接到位置 3，再替换 (3,4)。三轮结束后，训练文本 abababab 编为 [258]；新文本 ababab 编为 [257, 256]，分别对应 abab 与 ab。两段文本均能解码还原。


\begin{figure}[htbp]
\centering
% Source: results/bpe_walkthrough.json, training.steps[1].
% Input [256,256,256,256]; overlapping pair count=3; actual replacements=2.
% Input fragment. Requires tikz + arrows.meta + calc.
% Fixed design box: 15.5 cm x 6.5 cm. Text content is not scaled.
\begingroup
\definecolor{tokenizerBlue}{HTML}{285D77}
\definecolor{tokenizerOchre}{HTML}{B36826}
\definecolor{tokenizerGrey}{HTML}{66747B}
\definecolor{tokenizerGrid}{HTML}{E1E6E9}
\begin{tikzpicture}[
  x=1cm,y=1cm,font=\small,
  token/.style={
    draw=tokenizerBlue,line width=.5pt,fill=tokenizerBlue!3,
    minimum height=.82cm,minimum width=1.52cm,
    inner xsep=1pt,inner ysep=2pt,outer sep=0pt,
    align=center,text=tokenizerBlue
  },
  result/.style={token,draw=tokenizerOchre,fill=tokenizerOchre!4,
    text=tokenizerOchre,minimum width=3.12cm},
  window/.style={draw=tokenizerBlue,line width=.5pt},
  link/.style={draw=tokenizerOchre!85,line width=.45pt},
  finish/.style={link,-{Stealth[length=1.5mm,width=1.05mm]}},
  note/.style={font=\footnotesize,text=tokenizerGrey}
]
  \path[use as bounding box] (0,0) rectangle (15.5,6.5);
  \node[anchor=west] at (.10,6.12) {\textbf{(a)}\quad 相邻窗口};
  \node[anchor=west] at (8.45,6.12) {\textbf{(b)}\quad 不重叠替换};
  \draw[tokenizerGrid,line width=.35pt] (7.65,.65) -- (7.65,5.98);

  % Both panels repeat the same four-token input, at the same scale.
  \foreach \i in {0,...,3} {
    \pgfmathtruncatemacro{\position}{\i+1}
    \node[note] at ({.90+1.6*\i},5.57) {$\position$};
    \node[note] at ({9.25+1.6*\i},5.57) {$\position$};
    \node[token] (a-\i) at ({.90+1.6*\i},4.80)
      {\texttt{ab}\\[-1pt]{\footnotesize\color{tokenizerGrey}\texttt{256}}};
    \node[token] (b-\i) at ({9.25+1.6*\i},4.80)
      {\texttt{ab}\\[-1pt]{\footnotesize\color{tokenizerGrey}\texttt{256}}};
  }

  % All three adjacent windows count, including the overlapping middle one.
  \foreach \left/\right/\yy/\label in {
    .14/3.26/3.75/{(1,2)},1.74/4.86/2.85/{(2,3)},3.34/6.46/1.95/{(3,4)}
  } {
    \draw[window] (\left,{\yy+.15}) -- (\left,\yy)
      -- (\right,\yy) -- (\right,{\yy+.15});
    \node[note,anchor=north] at ({(\left+\right)/2},{\yy-.14}) {$\label$};
  }

  % Left-to-right merging consumes positions 1,2 and then 3,4 exactly once.
  \node[result] (out-0) at (10.05,2.85)
    {\texttt{abab}\\[-1pt]{\footnotesize\color{tokenizerGrey}\texttt{257}}};
  \node[result] (out-1) at (13.25,2.85)
    {\texttt{abab}\\[-1pt]{\footnotesize\color{tokenizerGrey}\texttt{257}}};
  \foreach \i in {0,1} {
    \pgfmathtruncatemacro{\leftchild}{2*\i}
    \pgfmathtruncatemacro{\rightchild}{2*\i+1}
    \coordinate (join-\i) at ($(out-\i.north)+(0,.56)$);
    \draw[link] (b-\leftchild.south) -- (join-\i);
    \draw[link] (b-\rightchild.south) -- (join-\i);
    \draw[finish] (join-\i) -- (out-\i.north);
  }
  \node[note] at (10.05,2.08) {$(1,2)$};
  \node[note] at (13.25,2.08) {$(3,4)$};
  \node[font=\small,text=tokenizerBlue] at (3.30,.70) {$3$ 次计数};
  \node[font=\small,text=tokenizerOchre] at (11.65,.70) {$2$ 次替换};
\end{tikzpicture}
\endgroup

\caption{同一序列的两个视角：\BpeDocCodeInline{get\char95\relax\allowbreak{}stats} 检查所有相邻窗口，共计 3 个；\BpeDocCodeInline{merge} 从左到右只替换不重叠的相邻对，共 2 次。因此日志中的出现次数不能直接当成实际替换次数。数据来源：上述实测记录的第 2 轮。}
\label{fig:overlap}
\end{figure}


\Needspace{9\baselineskip}
\BpeDocSubsection{一个可以证伪错误理解的反例}{一个可以证伪错误理解的反例}


用 ab1ab2ab3ab4bcbcbc 训练两次合并后，规则为 ab\(\rightarrow\)256、bc\(\rightarrow\)257。对新文本 abcbcbcbc，ab 只出现 1 次，bc 出现 4 次；正确编码仍先执行编号更小的 ab 规则。


\begin{Verbatim}[fontsize=\small,breaklines=true,breakanywhere=true,breakautoindent=true,tabsize=4,breaksymbolleft={},breaksymbolright={}]
正确：[256, 99, 257, 257, 257]
错误地先按新文本频率合并 bc：[97, 257, 257, 257, 257]
\end{Verbatim}


这说明 \BpeDocCodeInline{encode} 中虽然也调用 \BpeDocCodeInline{get\char95\relax\allowbreak{}stats}，返回的出现次数却不是本轮选取依据；该函数还用于枚举当前有哪些相邻对。选择规则由 \BpeDocCodeInline{self.\allowbreak{}merges} 中的已有优先级决定。


\Needspace{9\baselineskip}
\BpeDocSubsection{与 \BpeDocCodeInline{cl100k\char95\relax\allowbreak{}base} 的普通文本对照}{与 cl100k\_base 的普通文本对照}


\begingroup\fontsize{9.3}{12.7}\selectfont\setstretch{1}
\setlength{\BpeDocTableWidth}{\dimexpr\linewidth-8\tabcolsep\relax}
\renewcommand{\arraystretch}{1.15}
\par\addvspace{0.45em}
\noindent\begin{minipage}{\linewidth}
\begin{tabular}{>{\raggedright\arraybackslash}p{0.130000\BpeDocTableWidth}>{\raggedright\arraybackslash}p{0.520000\BpeDocTableWidth}>{\raggedright\arraybackslash}p{0.130000\BpeDocTableWidth}>{\raggedright\arraybackslash}p{0.220000\BpeDocTableWidth}}
\toprule
\textbf{类别} & \textbf{输入（换行以转义表示）} & \textbf{\BpeDocCodeInline{token} 数} & \textbf{编码／还原} \\

\midrule
英文 & 'Hello, world!' & 4 & 一致 / 成功 \\[3pt]
德语 & 'Füße, Straße und Häuser.' & 10 & 一致 / 成功 \\[3pt]
西班牙语 & '¡Hola! ¿Cómo estás? El niño come piña.' & 15 & 一致 / 成功 \\[3pt]
中文 & '你好，世界！' & 7 & 一致 / 成功 \\[3pt]
数字 & '1234567890' & 4 & 一致 / 成功 \\[3pt]
换行 & 'line one\textbackslash{}n\textbackslash{}nline two\textbackslash{}n' & 6 & 一致 / 成功 \\[3pt]
Emoji & 'U+1F642U+1F680' & 5 & 一致 / 成功 \\[3pt]
空字符串 & '' & 0 & 一致 / 成功 \\[3pt]
\bottomrule
\end{tabular}
\end{minipage}\par\addvspace{0.65em}
\endgroup


PDF 中 Emoji 用 Unicode 码点名称显示；真实字符、完整 ID 与字节数据保存在 \BpeDocCodeInline{results/\allowbreak{}gpt4\char95\relax\allowbreak{}check.\allowbreak{}json} 和对应中文 TXT 中。


\begingroup\fontsize{9.3}{12.7}\selectfont\setstretch{1}
\setlength{\BpeDocTableWidth}{\dimexpr\linewidth-4\tabcolsep\relax}
\renewcommand{\arraystretch}{1.15}
\par\addvspace{0.45em}
\noindent\begin{minipage}{\linewidth}
\begin{tabular}{>{\raggedright\arraybackslash}p{0.270000\BpeDocTableWidth}>{\raggedright\arraybackslash}p{0.730000\BpeDocTableWidth}}
\toprule
\textbf{具体输入} & \textbf{真实 \BpeDocCodeInline{token} ID} \\

\midrule
Hello, world! & [9906, 11, 1917, 0] \\[3pt]
你好，世界！ & [57668, 53901, 3922, 3574, 244, 98220, 6447] \\[3pt]
1234567890 & [4513, 10961, 16474, 15] \\[3pt]
\bottomrule
\end{tabular}
\end{minipage}\par\addvspace{0.65em}
\endgroup


中文例子中的“世”对应的 UTF-8 字节被拆在两个 \BpeDocCodeInline{token} 中：ID 3574 对应 e4 b8，ID 244 对应 96。这两个 \BpeDocCodeInline{token} 单独都不是完整字符；拼回字节后 e4 b8 96 才是“世”。因此，不应把每个 \BpeDocCodeInline{token} 单独解码后出现的替换符当成分词错误。


数字 1234567890 先被正则分成 123、456、789、0，再在各块内部编码。字节编号也与 \BpeDocCodeInline{BasicTokenizer} 不同：原始字节 97（a）对应 \BpeDocCodeInline{cl100k\char95\relax\allowbreak{}base} 单字节 \BpeDocCodeInline{token} ID 64；空格字节 32 对应 ID 220。这就是 \BpeDocCodeInline{GPT4Tokenizer} 需要 \BpeDocCodeInline{byte\char95\relax\allowbreak{}shuffle} 的原因。


\Needspace{9\baselineskip}
\BpeDocSubsection{检查范围与运行方法}{检查范围与运行方法}


\BpeDocCodeInline{code/\allowbreak{}verify.\allowbreak{}py} 中的 17 组检查全部通过，覆盖重叠合并、尾部剩余、空输入、Unicode 完整往返、规则优先级、普通文本编码对照、编号重排和 \BpeDocCodeInline{GPT4Tokenizer} 不可训练等行为。有限样例不能证明所有可能文本均正确，也未测试特殊 \BpeDocCodeInline{token} 的完整往返。


\begin{Verbatim}[fontsize=\small,breaklines=true,breakanywhere=true,breakautoindent=true,tabsize=4,breaksymbolleft={},breaksymbolright={}]
python -m pip install -r requirements.txt
python code/verify.py
python code/language_mix.py
\end{Verbatim}


当前电脑也可以执行项目根目录的 \BpeDocCodeInline{run\char95\relax\allowbreak{}examples.\allowbreak{}ps1}。无需 GPU 或模型 API；详细中文步骤见“使用说明.txt”。官方源文件、MIT 许可证、提交号及文件哈希都随项目保留。


\FloatBarrier\par\bigskip

% Content source: reports/content/multilingual_sampling.json

\Needspace{9\baselineskip}
\BpeDocSection{多语言分词器的语料抽样与评估设计}{多语言分词器的语料抽样与评估设计}


我的方案是把语言比例看成需要验证的设计变量：先确定两种语言的使用需求，清洗并统计可用语料，再在相同词表大小和相同训练字符预算下比较若干比例。50:50 可以作为两种语言同等重要时的基线，但不能在实验之前称为最佳比例。本文完成的是文献核对和实验方案设计；尚未进行真实德语／西语语料的配比实验，也没有得到一个经实验证实的最优比例。


\Needspace{9\baselineskip}
\BpeDocSubsection{先明确：调的是分词器比例，不是大模型比例}{先明确：调的是分词器比例，不是大模型比例}


BPE 训练是在决定“哪些相邻片段值得合并、有限的词表要记住哪些常见片段”。如果某种语言在训练材料中占比更高，它的片段往往更容易获得合并机会。因此，语言比例会影响两种语言的切分长度。词表也存在跨语言共享片段，所以不能简单认定“德语占 70\% 的语料，就必然占 70\% 的词表”。


\begingroup\fontsize{9.3}{12.7}\selectfont\setstretch{1}
\setlength{\BpeDocTableWidth}{\dimexpr\linewidth-6\tabcolsep\relax}
\renewcommand{\arraystretch}{1.15}
\par\addvspace{0.45em}
\noindent\begin{minipage}{\linewidth}
\begin{tabular}{>{\raggedright\arraybackslash}p{0.230000\BpeDocTableWidth}>{\raggedright\arraybackslash}p{0.385000\BpeDocTableWidth}>{\raggedright\arraybackslash}p{0.385000\BpeDocTableWidth}}
\toprule
\textbf{要决定的比例} & \textbf{它影响什么} & \textbf{本项目如何处理} \\

\midrule
\BpeDocCodeInline{Tokenizer} 训练语料比例 & BPE 合并规则与共享词表 & 本项目主要研究对象 \\[3pt]
LLM 预训练语料比例 & 模型学习各语言的机会及能力 & 若做下游对照，应固定它，避免混淆 \\[3pt]
实际使用需求比例 & 系统服务哪些语言与任务 & 用于定义评价权重，不直接等同于前两者 \\[3pt]
\bottomrule
\end{tabular}
\end{minipage}\par\addvspace{0.65em}
\endgroup


例如，用户请求中有 70\% 德语，不代表分词器必须用 70\% 德语训练。德语是否已经能被高效切分、西语是否被切得过碎、语料是否覆盖真实业务，都需要一起考虑。反过来，分词器给两种语言相同训练份额，也不能保证大模型获得相同语言能力。


\Needspace{9\baselineskip}
\BpeDocSubsection{三篇论文提供了什么证据}{三篇论文提供了什么证据}


[R1] Zhang 等（2022）直接研究 tokenizer 训练中的语言不平衡。他们固定模型训练设置，改变分词器语料比例；在其机器翻译实验中，更均衡的比例往往较好，但性能对比例失衡也比直觉预期更稳健。他们强调，tokenizer 抽样与模型训练抽样需要分别研究。主要实验采用 SentencePiece Unigram，并以 BPE 做消融。研究是英语中心、最多涉及八种语言，未直接测试德语／西语组合，因此不能把结论写成“德西 1:1 已被证明最优”。


\begin{quote}\footnotesize\setlength{\parskip}{0pt} [R1] Zhang et al. (2022). How Robust is Neural Machine Translation to Language Imbalance in Multilingual \BpeDocCodeInline{Tokenizer} Training? 重点：§3.1、§3.4、§4、§6。
\par\nobreak
\url{https://arxiv.org/html/2204.14268v2}
\end{quote}


[R2] Rust 等（2021）比较多语言与单语言模型，发现语言在词表中是否得到充分表示，与下游表现有关。他们使用 fertility（每个词平均切成多少个子词）和被拆分词的比例描述切分程度。这支持“逐语言检查是否过度切分”，但不支持“\BpeDocCodeInline{token} 越少，模型必然越好”，也没有直接给出德语／西语应采用的采样比例。


\begin{quote}\footnotesize\setlength{\parskip}{0pt} [R2] Rust et al. (2021). How Good is Your \BpeDocCodeInline{Tokenizer}? On the Monolingual Performance of Multilingual Language Models. 重点：§4.2、§4.3 与相关性分析。
\par\nobreak
\url{https://aclanthology.org/2021.acl-long.243/}
\end{quote}


[R3] Chung 等（2023）的 UniMax 尽量均匀分配语言预算，同时限制小语种语料重复次数。论文还为不同采样策略分别训练词表，并以字符数衡量各语言语料规模。可借鉴的是“均衡、避免过度重复、先定义计量单位”。它研究的是多语言预训练整体设置，其性能提升不能全部归因于 tokenizer 比例，也不能直接当作 byte-level BPE 的最佳配比结论。


\begin{quote}\footnotesize\setlength{\parskip}{0pt} [R3] Chung et al. (2023). UniMax: Fairer and more Effective Language Sampling for Large-Scale Multilingual Pretraining. 重点：§3、§3.1、§4.2。
\par\nobreak
\url{https://arxiv.org/html/2304.09151v1}
\end{quote}


\begingroup\fontsize{9.3}{12.7}\selectfont\setstretch{1}
\setlength{\BpeDocTableWidth}{\dimexpr\linewidth-6\tabcolsep\relax}
\renewcommand{\arraystretch}{1.15}
\par\addvspace{0.45em}
\noindent\begin{minipage}{\linewidth}
\begin{tabular}{>{\raggedright\arraybackslash}p{0.230000\BpeDocTableWidth}>{\raggedright\arraybackslash}p{0.385000\BpeDocTableWidth}>{\raggedright\arraybackslash}p{0.385000\BpeDocTableWidth}}
\toprule
\textbf{证据层次} & \textbf{本文采用的内容} & \textbf{不能越过的边界} \\

\midrule
文献直接支持 & 在原论文的语言、模型和任务范围内，分词质量与表现有关；抽样阶段需要区分 & 不推广成所有 LLM、所有语言对的必然规律 \\[3pt]
可迁移、尚待德西验证 & 参考平滑抽样、字符计量、重复限制与逐语言切分诊断 & 应重新在真实德语／西语文本上检验 \\[3pt]
本报告自己的建议 & 候选配比、5 百万字符、8,192 词表、三种子与 5\% 门槛 & 是实验设计选择，不是论文给出的最优配置 \\[3pt]
\bottomrule
\end{tabular}
\end{minipage}\par\addvspace{0.65em}
\endgroup


\Needspace{9\baselineskip}
\BpeDocSubsection{收集语料时，比例前面还有三个条件}{收集语料时，比例前面还有三个条件}


\begin{itemize}

\item 内容质量：先检查语言标签，去掉乱码、模板、广告、重复或近重复文档。误标为西语的英文网页，不能算成有效西语训练材料。

\item 领域覆盖：两种语言分别覆盖目标任务需要的领域，例如通用文本、技术说明与对话。若一边全是新闻、另一边全是代码，比较出的差异同时包含语言与领域差异。

\item 计量单位：本文将比例定义为清洗后 Unicode 字符数占比，并固定正规化方式；不能混用文档数、字节数、词数和 \BpeDocCodeInline{token} 数。

\end{itemize}



具体可统一采用 NFC 正规化，并按 Unicode 码点计数，空格和标点也计入；训练、验证与测试使用完全一致的规则。这是一个可复现的工程口径，不表示字符数能完整衡量语义信息。文档长短不同，所以“各取 100 篇”不等于材料数量相同；UTF-8 中带重音字符可能占多个字节，所以字节占比也不等于字符占比。不要用正在比较的 tokenizer 输出 \BpeDocCodeInline{token} 数反过来定义输入份额，否则比较口径会随候选方案改变。


德语与西语共享不少字母和标点，但仍应保留并检查各自常见的特殊字符、词形变化与长词。建议人工查看如德语复合词，以及西语带重音或变位的词，确认常见内容没有异常碎裂。这些是诊断样例，不应仅靠几个手选单词决定整个比例。


\Needspace{9\baselineskip}
\BpeDocSubsection{用平滑抽样生成候选比例}{用平滑抽样生成候选比例}


设 \(n_l\) 是语言 \(l\) 清洗后可用的字符数。可先用下面的指数平滑规则产生候选份额 \(q_l\)。\(\alpha\) 是需要尝试的参数，不是从论文复制后就永远适用的常数；这里仅考虑 \(n_l\) 大于零的语言。


\begin{equation}\label{eq:language-sampling}
q_l=\frac{n_l^{\alpha}}{\displaystyle\sum_j n_j^{\alpha}},\qquad 0\leq\alpha\leq1.
\end{equation}


\begin{figure}[htbp]
\centering
% Mathematical scenario, not measured performance:
% available character-count ratio n_de:n_es=9:1;
% p_i=n_i^alpha/sum_j(n_j^alpha), alpha in [0,1].
% Cross-check table: figures/language_sampling_key_points.csv.
% Input fragment. Requires pgfplots (compat=1.18 recommended).
% Requested outer design size: approximately 15.5 cm x 8 cm, including axes.
\begingroup
\definecolor{tokenizerBlue}{HTML}{285D77}
\definecolor{tokenizerOchre}{HTML}{B36826}
\definecolor{tokenizerGrey}{HTML}{66747B}
\definecolor{tokenizerGrid}{HTML}{E1E6E9}
\begin{tikzpicture}[font=\small]
\begin{axis}[
  width=15.5cm,height=8cm,
  xmin=0,xmax=1,ymin=0,ymax=1,enlargelimits=false,
  axis lines=left,
  axis line style={draw=tokenizerGrey,line width=.45pt},
  tick align=outside,
  tick style={draw=tokenizerGrey,line width=.4pt},
  major tick length=2.5pt,
  xtick={0,.25,.5,.75,1},
  xticklabels={$0$,$0.25$,$0.5$,$0.75$,$1$},
  ytick={0,.25,.5,.75,1},
  yticklabels={$0\%$,$25\%$,$50\%$,$75\%$,$100\%$},
  tick label style={font=\footnotesize,text=tokenizerGrey},
  xlabel={抽样指数 $\alpha$},ylabel={语料份额},
  label style={font=\small},
  xlabel near ticks,ylabel near ticks,
  grid=major,
  major grid style={draw=tokenizerGrid,line width=.3pt},
  legend style={
    at={(.025,.985)},anchor=north west,
    draw=none,fill=none,font=\footnotesize,
    legend columns=2,inner sep=1pt,
    /tikz/every even column/.append style={column sep=16pt}
  },
  legend cell align=left,
  clip=false,
  samples=101,domain=0:1
]
  \addplot[tokenizerBlue,line width=.6pt,no marks,forget plot]
    {9^x/(9^x+1)};
  \addplot[tokenizerOchre,line width=.6pt,dashed,no marks,forget plot]
    {1/(9^x+1)};

  % Different marks also distinguish the curves in grayscale.
  % At alpha=0 both markers share the same (0,.5) location, without jitter.
  \addplot[only marks,mark=square,mark size=2.5pt,
    mark options={draw=tokenizerOchre,fill=white,line width=.6pt},forget plot]
    coordinates {(0,.5) (.5,.25) (1,.1)};
  \addplot[only marks,mark=*,mark size=1.8pt,
    mark options={draw=tokenizerBlue,fill=white,line width=.6pt},forget plot]
    coordinates {(0,.5) (.5,.75) (1,.9)};

  \addlegendimage{tokenizerBlue,line width=.6pt,mark=*,mark size=1.8pt,
    mark options={draw=tokenizerBlue,fill=white,line width=.6pt}}
  \addlegendentry{德语}
  \addlegendimage{tokenizerOchre,line width=.6pt,dashed,mark=square,
    mark size=2.3pt,mark options={draw=tokenizerOchre,fill=white,line width=.6pt}}
  \addlegendentry{西班牙语}

  % Keep value labels close to their points, without placing text on a curve.
  \draw[tokenizerGrey,line width=.35pt]
    (axis cs:.009,.5) -- (axis cs:.12,.5);
  \node[anchor=west,font=\footnotesize,text=tokenizerGrey,inner sep=1pt]
    at (axis cs:.14,.5) {$50\%$};
  \node[anchor=south,font=\footnotesize,text=tokenizerBlue,yshift=4pt,inner sep=1pt]
    at (axis cs:.5,.75) {$75\%$};
  \node[anchor=north,font=\footnotesize,text=tokenizerOchre,yshift=-4pt,inner sep=1pt]
    at (axis cs:.5,.25) {$25\%$};
  \node[anchor=south east,font=\footnotesize,text=tokenizerBlue,
    xshift=-4pt,yshift=3pt,inner sep=1pt]
    at (axis cs:1,.9) {$90\%$};
  \node[anchor=north east,font=\footnotesize,text=tokenizerOchre,
    xshift=-4pt,yshift=-3pt,inner sep=1pt]
    at (axis cs:1,.1) {$10\%$};
\end{axis}
\end{tikzpicture}
\endgroup

\caption{假设清洗后德语与西语的库存字符数为 9:1，指数平滑参数 \(\alpha\) 如何改变训练份额。曲线由公式精确计算，标记点与 \BpeDocCodeInline{language\char95\relax\allowbreak{}mix.\allowbreak{}py} 的配额结果一致。这是数学情景，不是语料采集结果或模型性能曲线，也不能据此确定最佳 \(\alpha\)。完整点数据见 \BpeDocCodeInline{figures/\allowbreak{}language\char95\relax\allowbreak{}sampling.\allowbreak{}csv}。}
\label{fig:language-sampling}
\end{figure}


\begingroup\fontsize{9.3}{12.7}\selectfont\setstretch{1}
\setlength{\BpeDocTableWidth}{\dimexpr\linewidth-6\tabcolsep\relax}
\renewcommand{\arraystretch}{1.15}
\par\addvspace{0.45em}
\noindent\begin{minipage}{\linewidth}
\begin{tabular}{>{\raggedright\arraybackslash}p{0.230000\BpeDocTableWidth}>{\raggedright\arraybackslash}p{0.385000\BpeDocTableWidth}>{\raggedright\arraybackslash}p{0.385000\BpeDocTableWidth}}
\toprule
\textbf{\(\alpha\)} & \textbf{直观含义} & \textbf{假设原始德语:西语=9:1 时} \\

\midrule
1 & 保持清洗后的原始份额 & 9:1，即 90\%:10\% \\[3pt]
0.5 & 减小大语种的优势 & \(\sqrt{9}\):\(\sqrt{1}\)=3:1，即 75\%:25\% \\[3pt]
0 & 在语料非空的语言间均匀分配 & 1:1，即 50\%:50\% \\[3pt]
\bottomrule
\end{tabular}
\end{minipage}\par\addvspace{0.65em}
\endgroup


附带的 \BpeDocCodeInline{code/\allowbreak{}language\char95\relax\allowbreak{}mix.\allowbreak{}py} 已实际计算这个数学例子：假设德语有 900 万字符、西语有 100 万字符，训练预算为 100 万字符，\(\alpha\)=0、0.5、1 对应的德语／西语配额分别为 50 万／50 万、75 万／25 万、90 万／10 万。配额总和、份额总和、平方根平滑与零库存边界共 4 项检查通过。输入数量是假设值；程序没有训练或评估德语／西语分词器，完整输出在 \BpeDocCodeInline{results/\allowbreak{}language\char95\relax\allowbreak{}mix.\allowbreak{}json}。


表中的 9:1 只是数学例子，不是本次收集到的真实语料统计。指数越小，分配越均匀。这个规则帮助系统地提出候选，但不能表达所有业务偏好；还可以直接比较 30:70、50:50、70:30 等比例。若预计主要服务德语，应在报告中事先说明评价权重，同时保留西语的最低质量要求，而不是在看到结果后挑一个对自己有利的权重。


提高某语言的相对份额不一定需要重复文本：可以减少另一语言的抽样量。本文的 tokenizer 先导实验优先不放回抽样，即同一文档最多使用一次。若某个候选比例需要的唯一语料超过库存，应补充语料、对所有候选统一降低总预算，或明确标记该候选不可行；不能静默复制少量文本来凑比例。若另做带重复上限的方案，应记录上限、实际重复量及与无重复方案的差别。这是借鉴 UniMax 的控制思想，而不是声称复现其全部算法。


\Needspace{9\baselineskip}
\BpeDocSubsection{我提出的可复现实验方案（尚未执行）}{我提出的可复现实验方案（尚未执行）}


先做只训练 tokenizer 的低成本筛选，再决定是否值得训练下游模型。下面数值是可执行的先导配置，可按唯一语料量与计算资源统一缩小；它们不是经验最优参数。更换预算或词表大小后，应重跑全部候选。实现时使用能够承受该规模的 BPE 训练器；机制分析用 MinBPE 的速度不能代表生产训练器。


\begingroup\fontsize{9.3}{12.7}\selectfont\setstretch{1}
\setlength{\BpeDocTableWidth}{\dimexpr\linewidth-4\tabcolsep\relax}
\renewcommand{\arraystretch}{1.15}
\par\addvspace{0.45em}
\noindent\begin{minipage}{\linewidth}
\begin{tabular}{>{\raggedright\arraybackslash}p{0.290000\BpeDocTableWidth}>{\raggedright\arraybackslash}p{0.710000\BpeDocTableWidth}}
\toprule
\textbf{项目} & \textbf{预先固定的设置} \\

\midrule
数据划分 & 按完整文档或来源先去重，再划分训练／验证／测试；三者不共享重复或近重复内容 \\[3pt]
\BpeDocCodeInline{Tokenizer} 家族 & 同一种 byte-level BPE、同一预切分规则、NFC 规则与特殊 \BpeDocCodeInline{token} 策略 \\[3pt]
训练预算 & 每个候选共 5,000,000 个字符；不足时统一降低所有候选的预算 \\[3pt]
词表大小 & 8,192 个普通 \BpeDocCodeInline{token}，包含 256 个基础字节；特殊 \BpeDocCodeInline{token} 如需使用则另外固定 \\[3pt]
候选德语:西语份额 & 10:90、30:70、50:50、70:30、90:10；另加入实际库存的 \(\alpha\)=0.5 候选，重复项合并 \\[3pt]
验证／测试 & 两种语言分别留出 200,000 字符验证集和 200,000 字符最终测试集；不足则统一缩小并报告 \\[3pt]
抽样重复 & 每个候选用 3 个预先固定的随机种子独立抽样；记录种子、语料来源、哈希与实际份额 \\[3pt]
\bottomrule
\end{tabular}
\end{minipage}\par\addvspace{0.65em}
\endgroup


先按语言与领域分层，再按字符预算抽取完整文档，避免一个比例恰好抽到更多简单文本。完整文档长度可能让最终份额偏离目标，应设置小的预算容差并报告实际字符数与份额。如果必须截取片段，应先将文件解码成 Unicode 文本，再按完整文本边界截取；不能在原始 UTF-8 字节流中直接切到指定长度，否则可能把一个字符切坏。公平比较要求固定输入字符预算，而不是固定候选 tokenizer 各自产生的 \BpeDocCodeInline{token} 数。不同候选的输出 \BpeDocCodeInline{token} 数本来就是要测量的结果。


验证集用于比较候选、选比例；测试集在方案确定后才用于最终报告。即使所有数据都来自公开网页，也不能把训练集上的压缩效果写成泛化效果。随机种子的重复用于观察语料抽样带来的波动，不代表训练了三个大语言模型。


\Needspace{9\baselineskip}
\BpeDocSubsection{用什么指标决定保留哪个比例}{用什么指标决定保留哪个比例}


\begin{equation}\label{eq:tokens-per-character}
T_l=\frac{\#\,\mathrm{tokens}_l}{\#\,\mathrm{characters}_l}.
\end{equation}


\(T_l\) 表示平均每个字符需要多少 \BpeDocCodeInline{token}。在同一语言、同一份验证文本上，\(T_l\) 较低说明切分更紧凑。例：同一段 1,000 字符的文本被候选 A 切为 400 tokens，候选 B 切为 500 tokens，则 \(T\) 分别是 0.40 与 0.50。这只是指标计算例子，不是实际测量。具体计算时，逐个编码正规化后的非空文档，分子累加普通文本 \BpeDocCodeInline{token} 数，分母累加这些相同文本的 Unicode 字符数，包含空格与标点，不额外插入特殊 \BpeDocCodeInline{token}。不要先把文档拼接得可跨文档合并。本文使用总数之比，与逐句比率再平均不同，应在实现与报告中保持一致。


\begingroup\fontsize{9.3}{12.7}\selectfont\setstretch{1}
\setlength{\BpeDocTableWidth}{\dimexpr\linewidth-6\tabcolsep\relax}
\renewcommand{\arraystretch}{1.15}
\par\addvspace{0.45em}
\noindent\begin{minipage}{\linewidth}
\begin{tabular}{>{\raggedright\arraybackslash}p{0.230000\BpeDocTableWidth}>{\raggedright\arraybackslash}p{0.385000\BpeDocTableWidth}>{\raggedright\arraybackslash}p{0.385000\BpeDocTableWidth}}
\toprule
\textbf{指标／检查} & \textbf{作用} & \textbf{解释边界} \\

\midrule
逐语言 tokens / character & 比较同一语言在各候选下的序列长度 & 跨语言不能直接当成能力或公平分数 \\[3pt]
Fertility：每词平均子词数 & 观察词被拆分的程度；需固定词边界规则 & 德语长复合词与西语词形不同，不能要求两语数值相等 \\[3pt]
长文 \BpeDocCodeInline{token} 长度的 P95 & 观察较长输入的上下文压力 & 需固定文档集合；P95 是 95\% 文档不超过的长度 \\[3pt]
编码再解码一致性 & 确认规范化后的文本可完整还原 & 完整字节覆盖通常无 UNK，但不等于切分高效 \\[3pt]
常见词、专名、数字和混合语言样例 & 发现平均指标掩盖的局部异常 & 用于诊断，不替代独立验证集 \\[3pt]
\bottomrule
\end{tabular}
\end{minipage}\par\addvspace{0.65em}
\endgroup


Fertility 只作为辅助指标。为了让实现明确，可以先按空白和标点切出词，并在所有候选上使用同一规则，对词分别分词，计算每词 \BpeDocCodeInline{token} 数；应说明这是一种操作性定义，与整段文本编码可能不同。如果改用语言学词边界标注，则两种语言分别使用固定标准，并注明来源。最有意义的比较是“同一语言在新方案相对基线改善多少”，而不是直接比较德语和西语谁的 fertility 更低。


若两种语言同等重要，可以先以 50:50 为基线，再比较各语言 \(T\) 相对基线的变化。一个明确的选择规则是：先排除任一语言相对基线恶化超过预设容忍度的候选，再在剩余候选中选加权平均相对长度最小者。容忍度和权重属于项目要求，必须在看结果之前写下；例如先导阶段可暂设每语不恶化超过 5\%，两语各占一半权重，再做敏感性检查。这里的 5\% 是提议的工程门槛，不是论文结论。


\begin{equation}\label{eq:comparison-objective}
\begin{aligned}
R_l(q)&=\frac{T_l(q)}{T_l(0.5)},\\[4pt]
J(q)&=w_{\mathrm{de}}R_{\mathrm{de}}(q)+w_{\mathrm{es}}R_{\mathrm{es}}(q),\qquad w_{\mathrm{de}}+w_{\mathrm{es}}=1.
\end{aligned}
\end{equation}


报告每个候选的三次抽样均值和范围。若差异小于抽样波动，不应宣称某个比例明显更优；可保留几个折中候选。tokenizer 筛选完成后，若预算允许，再用相同模型架构、固定 LM 语言采样分布和预先约定的训练计算预算，对少数候选做下游验证。不同词表的每 \BpeDocCodeInline{token} 困惑度不能直接横比，可使用同一原文上的按字节或字符归一化损失，并辅以德语／西语目标任务评价。


\Needspace{9\baselineskip}
\BpeDocSubsection{最终观点与适用范围}{最终观点与适用范围}


我会从清洗后的库存比例、均匀比例和若干中间比例出发，用独立文本验证各语言的切分效率，在明确的业务权重与逐语言质量约束下选取比例；随后用下游任务检查这项选择是否真正有益。若暂时没有实验资源，我只把 50:50 作为双语同等重要时的合理起点，并保留“需要验证”的结论。


这套方案的局限是：字符预算并不等同于语义信息量；德语与西语的领域和词形差异可能影响指标；小规模词表实验不能直接外推到大规模模型；切分更紧凑只降低序列长度，不能独自证明理解与生成能力提高。文献也没有给出适用于所有德语／西语 LLM 的固定最佳比例。因此，本项目提出“有证据、有控制变量、有评价标准的选择方法”，不能用一个未经验证的数字代替论证。


\clearpage

\appendix

% Content source: reports/content/bpe_analysis.json

\Needspace{9\baselineskip}
\BpeDocSection{算法基础与源码审计补充材料}{算法基础与源码审计补充材料}


本节分析 BPE 的规则学习与编码机制，并将每个关键步骤对应到固定版本 MinBPE 的 Python 函数及代码行号。通过可手算的输入观察相邻对统计、合并优先级和 UTF-8 字节处理，再检查 \BpeDocCodeInline{cl100k\char95\relax\allowbreak{}base} 编码包装器的实现边界。


BPE 的核心做法是：先把文字拆成很小的单位，反复寻找最常一起出现的相邻两块，把它们合成一个新单位。训练结束后，这些合并规则固定下来。之后处理新文字时，分词器照着已有规则合并，不会重新学习一套词表。MinBPE 采用字节级 BPE，最小单位是 UTF-8 字节，而不是汉字或完整英文单词。


\begingroup\fontsize{9.3}{12.7}\selectfont\setstretch{1}
\setlength{\BpeDocTableWidth}{\dimexpr\linewidth-4\tabcolsep\relax}
\renewcommand{\arraystretch}{1.15}
\par\addvspace{0.45em}
\noindent\begin{minipage}{\linewidth}
\begin{tabular}{>{\raggedright\arraybackslash}p{0.290000\BpeDocTableWidth}>{\raggedright\arraybackslash}p{0.710000\BpeDocTableWidth}}
\toprule
\textbf{名词} & \textbf{先这样理解} \\

\midrule
\BpeDocCodeInline{token} & 模型处理文字的一个单位；可能是单个字节、一个字的一部分、几个字母、单词或其他片段，并不等于“单词”。 \\[3pt]
\BpeDocCodeInline{token} ID & 这个单位的整数编号。编号是查表用的标签，数值大小本身不表示词义。 \\[3pt]
字节（byte） & 计算机保存数据的基本单位之一，一个字节的数值在 0 到 255 之间。 \\[3pt]
UTF-8 & 把文字转换成字节序列的一种规则。英文 a 只需一个字节，汉字通常需要多个字节。 \\[3pt]
\BpeDocCodeInline{vocab}（词表） & 在本仓库中主要指“\BpeDocCodeInline{token} ID \(\rightarrow\) 对应字节串”的表，用于恢复文本。 \\[3pt]
\BpeDocCodeInline{merges}（合并表） & “两个相邻 \BpeDocCodeInline{token} 的编号 \(\rightarrow\) 新 \BpeDocCodeInline{token} 编号”的表。 \\[3pt]
\BpeDocCodeInline{train} / \BpeDocCodeInline{encode} / \BpeDocCodeInline{decode} & 分别是学习合并规则、将文字变成编号、将编号恢复为文字。不要把三者混为一谈。 \\[3pt]
\bottomrule
\end{tabular}
\end{minipage}\par\addvspace{0.65em}
\endgroup


\Needspace{9\baselineskip}
\BpeDocSubsection{先手算一遍：abababab 如何被压缩成一个 \BpeDocCodeInline{token}}{先手算一遍：abababab 如何被压缩成一个 token}


下面是教学手算，不是某个真实 GPT-4 输出。为便于看清过程，只用英文 a 和 b，并按 \BpeDocCodeInline{BasicTokenizer} 的实现设定 \BpeDocCodeInline{vocab\char95\relax\allowbreak{}size}=259。初始 256 个 \BpeDocCodeInline{token} 对应所有可能的单字节；a 的字节值是 97，b 的字节值是 98。目标词表大小为 259，所以要学习 259\(-\)256=3 条合并规则。


\begingroup\fontsize{9.3}{12.7}\selectfont\setstretch{1}
\setlength{\BpeDocTableWidth}{\dimexpr\linewidth-8\tabcolsep\relax}
\renewcommand{\arraystretch}{1.15}
\par\addvspace{0.45em}
\noindent\begin{minipage}{\linewidth}
\begin{tabular}{>{\raggedright\arraybackslash}p{0.140000\BpeDocTableWidth}>{\raggedright\arraybackslash}p{0.290000\BpeDocTableWidth}>{\raggedright\arraybackslash}p{0.340000\BpeDocTableWidth}>{\raggedright\arraybackslash}p{0.230000\BpeDocTableWidth}}
\toprule
\textbf{阶段} & \textbf{当前序列} & \textbf{统计与选择} & \textbf{合并结果} \\

\midrule
开始 & [97,98,97,98,97,98,97,98] & (97,98) 出现 4 次；(98,97) 出现 3 次。选择 (97,98)，新编号 256。 & [256,256,256,256] \\[3pt]
第 2 轮 & [256,256,256,256] & 相邻窗口中的 (256,256) 共 3 次；把它定为新编号 257。 & [257,257] \\[3pt]
第 3 轮 & [257,257] & (257,257) 出现 1 次；把它定为新编号 258。 & [258] \\[3pt]
\bottomrule
\end{tabular}
\end{minipage}\par\addvspace{0.65em}
\endgroup


第 2 轮最容易看错：四个 256 之间有三个相邻窗口，所以统计次数是 3；但真正替换时，一个位置不能被两个合并同时使用。程序从左向右先合并第 1、2 个，再合并第 3、4 个，因此只发生 2 次替换。相邻对的统计允许重叠，实际替换采用不重叠配对；二者不必相等。


\begingroup\fontsize{9.3}{12.7}\selectfont\setstretch{1}
\setlength{\BpeDocTableWidth}{\dimexpr\linewidth-6\tabcolsep\relax}
\renewcommand{\arraystretch}{1.15}
\par\addvspace{0.45em}
\noindent\begin{minipage}{\linewidth}
\begin{tabular}{>{\raggedright\arraybackslash}p{0.230000\BpeDocTableWidth}>{\raggedright\arraybackslash}p{0.385000\BpeDocTableWidth}>{\raggedright\arraybackslash}p{0.385000\BpeDocTableWidth}}
\toprule
\textbf{新 \BpeDocCodeInline{token} ID} & \textbf{合并规则} & \textbf{对应的实际字节内容} \\

\midrule
256 & (97,98) \(\rightarrow\) 256 & ab \\[3pt]
257 & (256,256) \(\rightarrow\) 257 & abab \\[3pt]
258 & (257,257) \(\rightarrow\) 258 & abababab \\[3pt]
\bottomrule
\end{tabular}
\end{minipage}\par\addvspace{0.65em}
\endgroup


训练后，词表仍保留原来的 256 个单字节 \BpeDocCodeInline{token}，并增加这 3 个新 \BpeDocCodeInline{token}。不是把 a、b 删除了。输入另一个字符串 ababab 时，按已有规则可先得到 [256,256,256]，再得到 [257,256]；最终分别代表 abab 和 ab。训练没有学过 (257,256) 的合并规则，所以不能继续把它合成一个 \BpeDocCodeInline{token}。


下面是附带 \BpeDocCodeInline{code/\allowbreak{}bpe\char95\relax\allowbreak{}walkthrough.\allowbreak{}py} 中的核心调用方式。完整脚本已配置 vendor 源码的导入路径，直接运行该脚本即可。这里只训练一个很小的 \BpeDocCodeInline{BasicTokenizer}；编号由该微型词表决定。手算已与实际程序结果核对，完整运行记录见后文。


\begin{Verbatim}[fontsize=\small,breaklines=true,breakanywhere=true,breakautoindent=true,tabsize=4,breaksymbolleft={},breaksymbolright={}]
from minbpe import BasicTokenizer

tokenizer = BasicTokenizer()
tokenizer.train("abababab", vocab_size=259, verbose=True)
print(tokenizer.encode("abababab"))  # 手算预期：[258]
print(tokenizer.encode("ababab"))    # 手算预期：[257, 256]
print(tokenizer.decode([257, 256]))  # 手算预期：ababab
\end{Verbatim}


\Needspace{9\baselineskip}
\BpeDocSubsubsection{再看中文：“你好”为什么从 6 个字节开始}{再看中文：“你好”为什么从 6 个字节开始}


英文例子中，一个字母恰好只占一个字节；中文不能照搬这种对应关系。Python 中“你好”有两个 Unicode 字符，却会被 UTF-8 编码成六个字节。


\begingroup\fontsize{9.3}{12.7}\selectfont\setstretch{1}
\setlength{\BpeDocTableWidth}{\dimexpr\linewidth-6\tabcolsep\relax}
\renewcommand{\arraystretch}{1.15}
\par\addvspace{0.45em}
\noindent\begin{minipage}{\linewidth}
\begin{tabular}{>{\raggedright\arraybackslash}p{0.230000\BpeDocTableWidth}>{\raggedright\arraybackslash}p{0.385000\BpeDocTableWidth}>{\raggedright\arraybackslash}p{0.385000\BpeDocTableWidth}}
\toprule
\textbf{文字} & \textbf{UTF-8 字节值} & \textbf{字节数} \\

\midrule
你 & 228, 189, 160 & 3 \\[3pt]
好 & 229, 165, 189 & 3 \\[3pt]
你好 & 228, 189, 160, 229, 165, 189 & 6 \\[3pt]
\bottomrule
\end{tabular}
\end{minipage}\par\addvspace{0.65em}
\endgroup


只用 abababab 训练的微型分词器没有学过这些中文片段的合并，所以“你好”会保留为 6 个初始 \BpeDocCodeInline{token}。解码时把这些字节完整接起来，再按 UTF-8 还原，仍能得到“你好”。有更多中文训练材料的分词器可以学到不同的合并，因此不能据此断言所有模型都把“你好”切成 6 个 \BpeDocCodeInline{token}。实际程序已核对这组字节和完整往返。


\Needspace{9\baselineskip}
\BpeDocSubsubsection{Python 语法与实现约定}{Python 语法与实现约定}


\begingroup\fontsize{9.3}{12.7}\selectfont\setstretch{1}
\setlength{\BpeDocTableWidth}{\dimexpr\linewidth-4\tabcolsep\relax}
\renewcommand{\arraystretch}{1.15}
\begin{longtable}{>{\raggedright\arraybackslash}p{0.290000\BpeDocTableWidth}>{\raggedright\arraybackslash}p{0.710000\BpeDocTableWidth}}
\toprule
\textbf{写法} & \textbf{本项目中的意思} \\

\midrule\endfirsthead
\toprule
\textbf{写法} & \textbf{本项目中的意思} \\

\midrule\endhead
\midrule\endfoot
\bottomrule\endlastfoot
\BpeDocCodeInline{self} & 当前这个分词器对象。\BpeDocCodeInline{self.\allowbreak{}merges} 表示这个对象保存的合并表。 \\[3pt]
\BpeDocCodeInline{class} A(B) & 定义一种 A 对象，它继承 B 已有的功能；A 还可以补充或改写功能。 \\[3pt]
\BpeDocCodeInline{\char95\relax\allowbreak{}\char95\relax\allowbreak{}init\char95\relax\allowbreak{}\char95\relax} / \BpeDocCodeInline{super}().\BpeDocCodeInline{\char95\relax\allowbreak{}\char95\relax\allowbreak{}init\char95\relax\allowbreak{}\char95\relax}() & 创建对象时执行的初始化方法；\BpeDocCodeInline{super}() 用来调用父类版本，先建立公共状态。 \\[3pt]
字典 d[key] 与 \BpeDocCodeInline{d.\allowbreak{}get}(key, default) & 前者直接查一个键；后者在键不存在时返回 default。这里常用默认值 0 或正无穷。 \\[3pt]
\BpeDocCodeInline{zip}(\BpeDocCodeInline{ids}, \BpeDocCodeInline{ids}[1:]) & 把原列表与“去掉第一个元素后的列表”并排配对，得到每一对相邻编号。 \\[3pt]
\BpeDocCodeInline{max}(\BpeDocCodeInline{stats}, key=\BpeDocCodeInline{stats.\allowbreak{}get}) & 遍历字典里的键，按各键对应的统计次数比较，选次数最大的那一对。不是选数值最大的 \BpeDocCodeInline{token} ID。 \\[3pt]
\BpeDocCodeInline{lambda} p: ... & 临时写一个短函数；这里规定比较相邻对 p 时应该看它的合并顺序。 \\[3pt]
\BpeDocCodeInline{min}(..., key=...) & 按 key 指定的标准找最小项。编码时用它选最早学到、优先级最高的可用合并。 \\[3pt]
\BpeDocCodeInline{bytes}([97,98]) 与 [97,98] & 前者是字节串 b'ab'；后者是整数列表。内容可以对应，但它们是不同类型。 \\[3pt]
列表推导式 [f(x) for x in xs] & 依次对 xs 中的每个元素执行 f(x)，收集为新列表。 \\[3pt]
\BpeDocCodeInline{append} 与 \BpeDocCodeInline{extend} & \BpeDocCodeInline{append} 把一个对象加到末尾；\BpeDocCodeInline{extend} 把另一个列表中的各个元素依次接到末尾。 \\[3pt]
\BpeDocCodeInline{return}、\BpeDocCodeInline{break}、\BpeDocCodeInline{continue} & 分别是返回函数结果、结束当前循环、跳过本轮并进入下一轮。 \\[3pt]
\end{longtable}
\endgroup


理解下面四个文件的分工即可：\BpeDocCodeInline{base.\allowbreak{}py} 提供公共工具；\BpeDocCodeInline{basic.\allowbreak{}py} 展示最简训练和编码；\BpeDocCodeInline{regex.\allowbreak{}py} 在它的思路上增加正则预切分与特殊 \BpeDocCodeInline{token}；\BpeDocCodeInline{gpt4.\allowbreak{}py} 读取已有的 \BpeDocCodeInline{cl100k\char95\relax\allowbreak{}base} 规则，并处理其编号细节。代码段中的空行和英文说明不执行；配注重点逐一解释每条实际执行语句。所有行号均按本报告固定的 MinBPE 提交。


\Needspace{18\baselineskip}
\BpeDocSubsection{\BpeDocCodeInline{base.\allowbreak{}py}：\BpeDocCodeInline{get\char95\relax\allowbreak{}stats} 统计相邻对，13–22 行}{base.py：get\_stats 统计相邻对，13–22 行}


\Needspace{4\baselineskip}
\begingroup\fontsize{8.9}{12.2}\selectfont\setstretch{1}
\setlength{\BpeDocTableWidth}{\dimexpr\linewidth-6\tabcolsep\relax}
\renewcommand{\arraystretch}{1.08}
\par\addvspace{0.45em}
\noindent\begin{minipage}{\linewidth}
{\small \BpeDocCodeInline{minbpe/\allowbreak{}base.\allowbreak{}py}，第 13–22 行（仅列执行语句）}\par\vspace{0.35em}
\begin{tabular}{>{\raggedleft\arraybackslash}p{0.065\BpeDocTableWidth}>{\raggedright\arraybackslash}p{0.475\BpeDocTableWidth}>{\raggedright\arraybackslash}p{0.460\BpeDocTableWidth}}
\toprule
\textbf{行号} & \textbf{原代码} & \textbf{逐行解释} \\
\midrule
\SourceLineNumber{13} & \BpeDocCodeCell{d\allowbreak{}e\allowbreak{}f\allowbreak{}\char32\relax\allowbreak{}g\allowbreak{}e\allowbreak{}t\allowbreak{}\char95\relax\allowbreak{}s\allowbreak{}t\allowbreak{}a\allowbreak{}t\allowbreak{}s\allowbreak{}(\allowbreak{}i\allowbreak{}d\allowbreak{}s\allowbreak{},\allowbreak{}\char32\relax\allowbreak{}c\allowbreak{}o\allowbreak{}u\allowbreak{}n\allowbreak{}t\allowbreak{}s\allowbreak{}=\allowbreak{}N\allowbreak{}o\allowbreak{}n\allowbreak{}e\allowbreak{})\allowbreak{}:} & 定义统计函数。\BpeDocCodeInline{ids} 是整数列表；\BpeDocCodeInline{counts} 可传入已有统计表，默认不传。 \\[3pt]
\SourceLineNumber{19} & \BpeDocCodeCell{\char32\relax\allowbreak{}\char32\relax\allowbreak{}\char32\relax\allowbreak{}\char32\relax\allowbreak{}c\allowbreak{}o\allowbreak{}u\allowbreak{}n\allowbreak{}t\allowbreak{}s\allowbreak{}\char32\relax\allowbreak{}=\allowbreak{}\char32\relax\allowbreak{}\char123\relax\allowbreak{}\char125\relax\allowbreak{}\char32\relax\allowbreak{}i\allowbreak{}f\allowbreak{}\char32\relax\allowbreak{}c\allowbreak{}o\allowbreak{}u\allowbreak{}n\allowbreak{}t\allowbreak{}s\allowbreak{}\char32\relax\allowbreak{}i\allowbreak{}s\allowbreak{}\char32\relax\allowbreak{}N\allowbreak{}o\allowbreak{}n\allowbreak{}e\allowbreak{}\char32\relax\allowbreak{}e\allowbreak{}l\allowbreak{}s\allowbreak{}e\allowbreak{}\char32\relax\allowbreak{}c\allowbreak{}o\allowbreak{}u\allowbreak{}n\allowbreak{}t\allowbreak{}s} & 没传 \BpeDocCodeInline{counts} 就新建空字典；传了就继续使用同一个字典，保留已有计数。 \\[3pt]
\SourceLineNumber{20} & \BpeDocCodeCell{\char32\relax\allowbreak{}\char32\relax\allowbreak{}\char32\relax\allowbreak{}\char32\relax\allowbreak{}f\allowbreak{}o\allowbreak{}r\allowbreak{}\char32\relax\allowbreak{}p\allowbreak{}a\allowbreak{}i\allowbreak{}r\allowbreak{}\char32\relax\allowbreak{}i\allowbreak{}n\allowbreak{}\char32\relax\allowbreak{}z\allowbreak{}i\allowbreak{}p\allowbreak{}(\allowbreak{}i\allowbreak{}d\allowbreak{}s\allowbreak{},\allowbreak{}\char32\relax\allowbreak{}i\allowbreak{}d\allowbreak{}s\allowbreak{}[\allowbreak{}1\allowbreak{}:\allowbreak{}]\allowbreak{})\allowbreak{}:\allowbreak{}\char32\relax\allowbreak{}\char35\relax\allowbreak{}\char32\relax\allowbreak{}i\allowbreak{}t\allowbreak{}e\allowbreak{}r\allowbreak{}a\allowbreak{}t\allowbreak{}e\allowbreak{}\char32\relax\allowbreak{}c\allowbreak{}o\allowbreak{}n\allowbreak{}s\allowbreak{}e\allowbreak{}c\allowbreak{}u\allowbreak{}t\allowbreak{}i\allowbreak{}v\allowbreak{}e\allowbreak{}\char32\relax\allowbreak{}e\allowbreak{}l\allowbreak{}e\allowbreak{}m\allowbreak{}e\allowbreak{}n\allowbreak{}t\allowbreak{}s} & \BpeDocCodeInline{ids}[1:] 从第 2 个元素开始；\BpeDocCodeInline{zip} 将它与原列表配成相邻对，循环逐对处理。 \\[3pt]
\SourceLineNumber{21} & \BpeDocCodeCell{\char32\relax\allowbreak{}\char32\relax\allowbreak{}\char32\relax\allowbreak{}\char32\relax\allowbreak{}\char32\relax\allowbreak{}\char32\relax\allowbreak{}\char32\relax\allowbreak{}\char32\relax\allowbreak{}c\allowbreak{}o\allowbreak{}u\allowbreak{}n\allowbreak{}t\allowbreak{}s\allowbreak{}[\allowbreak{}p\allowbreak{}a\allowbreak{}i\allowbreak{}r\allowbreak{}]\allowbreak{}\char32\relax\allowbreak{}=\allowbreak{}\char32\relax\allowbreak{}c\allowbreak{}o\allowbreak{}u\allowbreak{}n\allowbreak{}t\allowbreak{}s\allowbreak{}.\allowbreak{}g\allowbreak{}e\allowbreak{}t\allowbreak{}(\allowbreak{}p\allowbreak{}a\allowbreak{}i\allowbreak{}r\allowbreak{},\allowbreak{}\char32\relax\allowbreak{}0\allowbreak{})\allowbreak{}\char32\relax\allowbreak{}+\allowbreak{}\char32\relax\allowbreak{}1} & 读出 \BpeDocCodeInline{pair} 原有次数；不存在时按 0 处理，再加 1 并存回字典。 \\[3pt]
\SourceLineNumber{22} & \BpeDocCodeCell{\char32\relax\allowbreak{}\char32\relax\allowbreak{}\char32\relax\allowbreak{}\char32\relax\allowbreak{}r\allowbreak{}e\allowbreak{}t\allowbreak{}u\allowbreak{}r\allowbreak{}n\allowbreak{}\char32\relax\allowbreak{}c\allowbreak{}o\allowbreak{}u\allowbreak{}n\allowbreak{}t\allowbreak{}s} & 返回统计字典。传入已有字典时，原字典也已经被直接更新。 \\[3pt]
\bottomrule
\end{tabular}
\end{minipage}\par\addvspace{0.65em}
\endgroup


例如 \BpeDocCodeInline{ids}=[97,98,97]：\BpeDocCodeInline{zip} 的结果是 (97,98)、(98,97)。对 [256,256,256,256]，结果是三个相同的 (256,256)，因此计数为 3。这里没有做任何合并，也没有决定新 \BpeDocCodeInline{token} 编号。它只负责回答“每一种相邻对出现了多少次”。


\Needspace{9\baselineskip}
\BpeDocSubsubsection{\BpeDocCodeInline{base.\allowbreak{}py}：\BpeDocCodeInline{merge} 从左向右替换一类相邻对，25–41 行}{base.py：merge 从左向右替换一类相邻对，25–41 行}


\Needspace{4\baselineskip}
\begingroup\fontsize{8.9}{12.2}\selectfont\setstretch{1}
\setlength{\BpeDocTableWidth}{\dimexpr\linewidth-6\tabcolsep\relax}
\renewcommand{\arraystretch}{1.08}
{\small \BpeDocCodeInline{minbpe/\allowbreak{}base.\allowbreak{}py}，第 25–41 行（仅列执行语句）}

\begin{longtable}{>{\raggedleft\arraybackslash}p{0.065\BpeDocTableWidth}>{\raggedright\arraybackslash}p{0.475\BpeDocTableWidth}>{\raggedright\arraybackslash}p{0.460\BpeDocTableWidth}}
\toprule
\textbf{行号} & \textbf{原代码} & \textbf{逐行解释} \\
\midrule\endfirsthead
\toprule
\textbf{行号} & \textbf{原代码} & \textbf{逐行解释} \\
\midrule\endhead
\midrule\endfoot
\bottomrule\endlastfoot
\SourceLineNumber{25} & \BpeDocCodeCell{d\allowbreak{}e\allowbreak{}f\allowbreak{}\char32\relax\allowbreak{}m\allowbreak{}e\allowbreak{}r\allowbreak{}g\allowbreak{}e\allowbreak{}(\allowbreak{}i\allowbreak{}d\allowbreak{}s\allowbreak{},\allowbreak{}\char32\relax\allowbreak{}p\allowbreak{}a\allowbreak{}i\allowbreak{}r\allowbreak{},\allowbreak{}\char32\relax\allowbreak{}i\allowbreak{}d\allowbreak{}x\allowbreak{})\allowbreak{}:} & 定义替换函数：\BpeDocCodeInline{ids} 是原序列，\BpeDocCodeInline{pair} 是要合并的两个编号，\BpeDocCodeInline{idx} 是合并后的新编号。 \\[3pt]
\SourceLineNumber{31} & \BpeDocCodeCell{\char32\relax\allowbreak{}\char32\relax\allowbreak{}\char32\relax\allowbreak{}\char32\relax\allowbreak{}n\allowbreak{}e\allowbreak{}w\allowbreak{}i\allowbreak{}d\allowbreak{}s\allowbreak{}\char32\relax\allowbreak{}=\allowbreak{}\char32\relax\allowbreak{}[\allowbreak{}]} & 建立空列表，用来收集合并后的结果。 \\[3pt]
\SourceLineNumber{32} & \BpeDocCodeCell{\char32\relax\allowbreak{}\char32\relax\allowbreak{}\char32\relax\allowbreak{}\char32\relax\allowbreak{}i\allowbreak{}\char32\relax\allowbreak{}=\allowbreak{}\char32\relax\allowbreak{}0} & 从原序列的第 0 个位置开始检查。Python 下标从 0 开始。 \\[3pt]
\SourceLineNumber{33} & \BpeDocCodeCell{\char32\relax\allowbreak{}\char32\relax\allowbreak{}\char32\relax\allowbreak{}\char32\relax\allowbreak{}w\allowbreak{}h\allowbreak{}i\allowbreak{}l\allowbreak{}e\allowbreak{}\char32\relax\allowbreak{}i\allowbreak{}\char32\relax\allowbreak{}<\allowbreak{}\char32\relax\allowbreak{}l\allowbreak{}e\allowbreak{}n\allowbreak{}(\allowbreak{}i\allowbreak{}d\allowbreak{}s\allowbreak{})\allowbreak{}:} & 只要还没走到原序列末尾，就继续处理。 \\[3pt]
\SourceLineNumber{35} & \BpeDocCodeCell{\char32\relax\allowbreak{}\char32\relax\allowbreak{}\char32\relax\allowbreak{}\char32\relax\allowbreak{}\char32\relax\allowbreak{}\char32\relax\allowbreak{}\char32\relax\allowbreak{}\char32\relax\allowbreak{}i\allowbreak{}f\allowbreak{}\char32\relax\allowbreak{}i\allowbreak{}d\allowbreak{}s\allowbreak{}[\allowbreak{}i\allowbreak{}]\allowbreak{}\char32\relax\allowbreak{}=\allowbreak{}=\allowbreak{}\char32\relax\allowbreak{}p\allowbreak{}a\allowbreak{}i\allowbreak{}r\allowbreak{}[\allowbreak{}0\allowbreak{}]\allowbreak{}\char32\relax\allowbreak{}a\allowbreak{}n\allowbreak{}d\allowbreak{}\char32\relax\allowbreak{}i\allowbreak{}\char32\relax\allowbreak{}<\allowbreak{}\char32\relax\allowbreak{}l\allowbreak{}e\allowbreak{}n\allowbreak{}(\allowbreak{}i\allowbreak{}d\allowbreak{}s\allowbreak{})\allowbreak{}\char32\relax\allowbreak{}-\allowbreak{}\char32\relax\allowbreak{}1\allowbreak{}\char32\relax\allowbreak{}a\allowbreak{}n\allowbreak{}d\allowbreak{}\char32\relax\allowbreak{}i\allowbreak{}d\allowbreak{}s\allowbreak{}[\allowbreak{}i\allowbreak{}+\allowbreak{}1\allowbreak{}]\allowbreak{}\char32\relax\allowbreak{}=\allowbreak{}=\allowbreak{}\char32\relax\allowbreak{}p\allowbreak{}a\allowbreak{}i\allowbreak{}r\allowbreak{}[\allowbreak{}1\allowbreak{}]\allowbreak{}:} & 检查当前编号、右侧位置是否存在、右侧编号是否匹配。and 从左向右短路求值，边界检查通过后才访问 \BpeDocCodeInline{ids}[i+1]。 \\[3pt]
\SourceLineNumber{36} & \BpeDocCodeCell{\char32\relax\allowbreak{}\char32\relax\allowbreak{}\char32\relax\allowbreak{}\char32\relax\allowbreak{}\char32\relax\allowbreak{}\char32\relax\allowbreak{}\char32\relax\allowbreak{}\char32\relax\allowbreak{}\char32\relax\allowbreak{}\char32\relax\allowbreak{}\char32\relax\allowbreak{}\char32\relax\allowbreak{}n\allowbreak{}e\allowbreak{}w\allowbreak{}i\allowbreak{}d\allowbreak{}s\allowbreak{}.\allowbreak{}a\allowbreak{}p\allowbreak{}p\allowbreak{}e\allowbreak{}n\allowbreak{}d\allowbreak{}(\allowbreak{}i\allowbreak{}d\allowbreak{}x\allowbreak{})} & 匹配成功，把新 \BpeDocCodeInline{token} 编号追加到结果。 \\[3pt]
\SourceLineNumber{37} & \BpeDocCodeCell{\char32\relax\allowbreak{}\char32\relax\allowbreak{}\char32\relax\allowbreak{}\char32\relax\allowbreak{}\char32\relax\allowbreak{}\char32\relax\allowbreak{}\char32\relax\allowbreak{}\char32\relax\allowbreak{}\char32\relax\allowbreak{}\char32\relax\allowbreak{}\char32\relax\allowbreak{}\char32\relax\allowbreak{}i\allowbreak{}\char32\relax\allowbreak{}+\allowbreak{}=\allowbreak{}\char32\relax\allowbreak{}2} & 原来的两个位置都已使用，指针向右走 2 格，避免重叠替换。 \\[3pt]
\SourceLineNumber{38} & \BpeDocCodeCell{\char32\relax\allowbreak{}\char32\relax\allowbreak{}\char32\relax\allowbreak{}\char32\relax\allowbreak{}\char32\relax\allowbreak{}\char32\relax\allowbreak{}\char32\relax\allowbreak{}\char32\relax\allowbreak{}e\allowbreak{}l\allowbreak{}s\allowbreak{}e\allowbreak{}:} & 如果相邻对不匹配，进入此分支。 \\[3pt]
\SourceLineNumber{39} & \BpeDocCodeCell{\char32\relax\allowbreak{}\char32\relax\allowbreak{}\char32\relax\allowbreak{}\char32\relax\allowbreak{}\char32\relax\allowbreak{}\char32\relax\allowbreak{}\char32\relax\allowbreak{}\char32\relax\allowbreak{}\char32\relax\allowbreak{}\char32\relax\allowbreak{}\char32\relax\allowbreak{}\char32\relax\allowbreak{}n\allowbreak{}e\allowbreak{}w\allowbreak{}i\allowbreak{}d\allowbreak{}s\allowbreak{}.\allowbreak{}a\allowbreak{}p\allowbreak{}p\allowbreak{}e\allowbreak{}n\allowbreak{}d\allowbreak{}(\allowbreak{}i\allowbreak{}d\allowbreak{}s\allowbreak{}[\allowbreak{}i\allowbreak{}]\allowbreak{})} & 把当前这个原编号原样保留到结果中。 \\[3pt]
\SourceLineNumber{40} & \BpeDocCodeCell{\char32\relax\allowbreak{}\char32\relax\allowbreak{}\char32\relax\allowbreak{}\char32\relax\allowbreak{}\char32\relax\allowbreak{}\char32\relax\allowbreak{}\char32\relax\allowbreak{}\char32\relax\allowbreak{}\char32\relax\allowbreak{}\char32\relax\allowbreak{}\char32\relax\allowbreak{}\char32\relax\allowbreak{}i\allowbreak{}\char32\relax\allowbreak{}+\allowbreak{}=\allowbreak{}\char32\relax\allowbreak{}1} & 只处理了一个位置，因此指针向右走 1 格。 \\[3pt]
\SourceLineNumber{41} & \BpeDocCodeCell{\char32\relax\allowbreak{}\char32\relax\allowbreak{}\char32\relax\allowbreak{}\char32\relax\allowbreak{}r\allowbreak{}e\allowbreak{}t\allowbreak{}u\allowbreak{}r\allowbreak{}n\allowbreak{}\char32\relax\allowbreak{}n\allowbreak{}e\allowbreak{}w\allowbreak{}i\allowbreak{}d\allowbreak{}s} & 返回新列表。函数没有直接改写传入的原列表。 \\[3pt]
\end{longtable}
\endgroup


这个函数不统计频次，也不比较不同规则的优先级；它只执行调用者已经选好的那条规则。训练和编码都会使用它，但“如何选 \BpeDocCodeInline{pair}”由外层代码决定。函数内使用新列表，所以外层通常写 \BpeDocCodeInline{ids}=\BpeDocCodeInline{merge}(...)，让 \BpeDocCodeInline{ids} 指向新的结果。


\Needspace{18\baselineskip}
\BpeDocSubsection{\BpeDocCodeInline{basic.\allowbreak{}py}：创建最简单的分词器，15–18 行}{basic.py：创建最简单的分词器，15–18 行}


\Needspace{4\baselineskip}
\begingroup\fontsize{8.9}{12.2}\selectfont\setstretch{1}
\setlength{\BpeDocTableWidth}{\dimexpr\linewidth-6\tabcolsep\relax}
\renewcommand{\arraystretch}{1.08}
\par\addvspace{0.45em}
\noindent\begin{minipage}{\linewidth}
{\small \BpeDocCodeInline{minbpe/\allowbreak{}basic.\allowbreak{}py}，第 15–18 行（仅列执行语句）}\par\vspace{0.35em}
\begin{tabular}{>{\raggedleft\arraybackslash}p{0.065\BpeDocTableWidth}>{\raggedright\arraybackslash}p{0.475\BpeDocTableWidth}>{\raggedright\arraybackslash}p{0.460\BpeDocTableWidth}}
\toprule
\textbf{行号} & \textbf{原代码} & \textbf{逐行解释} \\
\midrule
\SourceLineNumber{15} & \BpeDocCodeCell{c\allowbreak{}l\allowbreak{}a\allowbreak{}s\allowbreak{}s\allowbreak{}\char32\relax\allowbreak{}B\allowbreak{}a\allowbreak{}s\allowbreak{}i\allowbreak{}c\allowbreak{}T\allowbreak{}o\allowbreak{}k\allowbreak{}e\allowbreak{}n\allowbreak{}i\allowbreak{}z\allowbreak{}e\allowbreak{}r\allowbreak{}(\allowbreak{}T\allowbreak{}o\allowbreak{}k\allowbreak{}e\allowbreak{}n\allowbreak{}i\allowbreak{}z\allowbreak{}e\allowbreak{}r\allowbreak{})\allowbreak{}:} & 定义 \BpeDocCodeInline{BasicTokenizer}，它继承公共 \BpeDocCodeInline{Tokenizer} 类。 \\[3pt]
\SourceLineNumber{17} & \BpeDocCodeCell{\char32\relax\allowbreak{}\char32\relax\allowbreak{}\char32\relax\allowbreak{}\char32\relax\allowbreak{}d\allowbreak{}e\allowbreak{}f\allowbreak{}\char32\relax\allowbreak{}\char95\relax\allowbreak{}\char95\relax\allowbreak{}i\allowbreak{}n\allowbreak{}i\allowbreak{}t\allowbreak{}\char95\relax\allowbreak{}\char95\relax\allowbreak{}(\allowbreak{}s\allowbreak{}e\allowbreak{}l\allowbreak{}f\allowbreak{})\allowbreak{}:} & 定义创建该对象时调用的初始化方法。 \\[3pt]
\SourceLineNumber{18} & \BpeDocCodeCell{\char32\relax\allowbreak{}\char32\relax\allowbreak{}\char32\relax\allowbreak{}\char32\relax\allowbreak{}\char32\relax\allowbreak{}\char32\relax\allowbreak{}\char32\relax\allowbreak{}\char32\relax\allowbreak{}s\allowbreak{}u\allowbreak{}p\allowbreak{}e\allowbreak{}r\allowbreak{}(\allowbreak{})\allowbreak{}.\allowbreak{}\char95\relax\allowbreak{}\char95\relax\allowbreak{}i\allowbreak{}n\allowbreak{}i\allowbreak{}t\allowbreak{}\char95\relax\allowbreak{}\char95\relax\allowbreak{}(\allowbreak{})} & 执行父类初始化，建立空合并表、初始词表等公共状态。 \\[3pt]
\bottomrule
\end{tabular}
\end{minipage}\par\addvspace{0.65em}
\endgroup


父类的初始词表含 256 个单字节 \BpeDocCodeInline{token}。\BpeDocCodeInline{BasicTokenizer} 不做正则预切分，也不专门识别特殊 \BpeDocCodeInline{token}；整个输入的字节序列直接进入 BPE。因此它适合先理解核心算法，但不能仅凭它声称已经复现 GPT-4 分词。


\Needspace{9\baselineskip}
\BpeDocSubsubsection{\BpeDocCodeInline{basic.\allowbreak{}py}：\BpeDocCodeInline{train} 学习合并规则，20–49 行}{basic.py：train 学习合并规则，20–49 行}


\Needspace{4\baselineskip}
\begingroup\fontsize{8.9}{12.2}\selectfont\setstretch{1}
\setlength{\BpeDocTableWidth}{\dimexpr\linewidth-6\tabcolsep\relax}
\renewcommand{\arraystretch}{1.08}
{\small \BpeDocCodeInline{minbpe/\allowbreak{}basic.\allowbreak{}py}，第 20–49 行（仅列执行语句）}

\begin{longtable}{>{\raggedleft\arraybackslash}p{0.065\BpeDocTableWidth}>{\raggedright\arraybackslash}p{0.475\BpeDocTableWidth}>{\raggedright\arraybackslash}p{0.460\BpeDocTableWidth}}
\toprule
\textbf{行号} & \textbf{原代码} & \textbf{逐行解释} \\
\midrule\endfirsthead
\toprule
\textbf{行号} & \textbf{原代码} & \textbf{逐行解释} \\
\midrule\endhead
\midrule\endfoot
\bottomrule\endlastfoot
\SourceLineNumber{20} & \BpeDocCodeCell{\char32\relax\allowbreak{}\char32\relax\allowbreak{}\char32\relax\allowbreak{}\char32\relax\allowbreak{}d\allowbreak{}e\allowbreak{}f\allowbreak{}\char32\relax\allowbreak{}t\allowbreak{}r\allowbreak{}a\allowbreak{}i\allowbreak{}n\allowbreak{}(\allowbreak{}s\allowbreak{}e\allowbreak{}l\allowbreak{}f\allowbreak{},\allowbreak{}\char32\relax\allowbreak{}t\allowbreak{}e\allowbreak{}x\allowbreak{}t\allowbreak{},\allowbreak{}\char32\relax\allowbreak{}v\allowbreak{}o\allowbreak{}c\allowbreak{}a\allowbreak{}b\allowbreak{}\char95\relax\allowbreak{}s\allowbreak{}i\allowbreak{}z\allowbreak{}e\allowbreak{},\allowbreak{}\char32\relax\allowbreak{}v\allowbreak{}e\allowbreak{}r\allowbreak{}b\allowbreak{}o\allowbreak{}s\allowbreak{}e\allowbreak{}=\allowbreak{}F\allowbreak{}a\allowbreak{}l\allowbreak{}s\allowbreak{}e\allowbreak{})\allowbreak{}:} & 定义训练方法。text 是训练文本；\BpeDocCodeInline{vocab\char95\relax\allowbreak{}size} 是目标词表大小；\BpeDocCodeInline{verbose} 控制是否打印过程。 \\[3pt]
\SourceLineNumber{21} & \BpeDocCodeCell{\char32\relax\allowbreak{}\char32\relax\allowbreak{}\char32\relax\allowbreak{}\char32\relax\allowbreak{}\char32\relax\allowbreak{}\char32\relax\allowbreak{}\char32\relax\allowbreak{}\char32\relax\allowbreak{}a\allowbreak{}s\allowbreak{}s\allowbreak{}e\allowbreak{}r\allowbreak{}t\allowbreak{}\char32\relax\allowbreak{}v\allowbreak{}o\allowbreak{}c\allowbreak{}a\allowbreak{}b\allowbreak{}\char95\relax\allowbreak{}s\allowbreak{}i\allowbreak{}z\allowbreak{}e\allowbreak{}\char32\relax\allowbreak{}>\allowbreak{}=\allowbreak{}\char32\relax\allowbreak{}2\allowbreak{}5\allowbreak{}6} & 要求词表至少容纳 256 个初始字节 \BpeDocCodeInline{token}，否则断言失败。 \\[3pt]
\SourceLineNumber{22} & \BpeDocCodeCell{\char32\relax\allowbreak{}\char32\relax\allowbreak{}\char32\relax\allowbreak{}\char32\relax\allowbreak{}\char32\relax\allowbreak{}\char32\relax\allowbreak{}\char32\relax\allowbreak{}\char32\relax\allowbreak{}n\allowbreak{}u\allowbreak{}m\allowbreak{}\char95\relax\allowbreak{}m\allowbreak{}e\allowbreak{}r\allowbreak{}g\allowbreak{}e\allowbreak{}s\allowbreak{}\char32\relax\allowbreak{}=\allowbreak{}\char32\relax\allowbreak{}v\allowbreak{}o\allowbreak{}c\allowbreak{}a\allowbreak{}b\allowbreak{}\char95\relax\allowbreak{}s\allowbreak{}i\allowbreak{}z\allowbreak{}e\allowbreak{}\char32\relax\allowbreak{}-\allowbreak{}\char32\relax\allowbreak{}2\allowbreak{}5\allowbreak{}6} & 每次合并新增一个 \BpeDocCodeInline{token}，所以训练轮数等于目标大小减去 256。 \\[3pt]
\SourceLineNumber{25} & \BpeDocCodeCell{\char32\relax\allowbreak{}\char32\relax\allowbreak{}\char32\relax\allowbreak{}\char32\relax\allowbreak{}\char32\relax\allowbreak{}\char32\relax\allowbreak{}\char32\relax\allowbreak{}\char32\relax\allowbreak{}t\allowbreak{}e\allowbreak{}x\allowbreak{}t\allowbreak{}\char95\relax\allowbreak{}b\allowbreak{}y\allowbreak{}t\allowbreak{}e\allowbreak{}s\allowbreak{}\char32\relax\allowbreak{}=\allowbreak{}\char32\relax\allowbreak{}t\allowbreak{}e\allowbreak{}x\allowbreak{}t\allowbreak{}.\allowbreak{}e\allowbreak{}n\allowbreak{}c\allowbreak{}o\allowbreak{}d\allowbreak{}e\allowbreak{}(\allowbreak{}"\allowbreak{}u\allowbreak{}t\allowbreak{}f\allowbreak{}-\allowbreak{}8\allowbreak{}"\allowbreak{})\allowbreak{}\char32\relax\allowbreak{}\char35\relax\allowbreak{}\char32\relax\allowbreak{}r\allowbreak{}a\allowbreak{}w\allowbreak{}\char32\relax\allowbreak{}b\allowbreak{}y\allowbreak{}t\allowbreak{}e\allowbreak{}s} & 按 UTF-8 规则把 Python 字符串转换为字节串。这里的 \BpeDocCodeInline{encode} 是字符串方法，不是分词器的 \BpeDocCodeInline{encode} 方法。 \\[3pt]
\SourceLineNumber{26} & \BpeDocCodeCell{\char32\relax\allowbreak{}\char32\relax\allowbreak{}\char32\relax\allowbreak{}\char32\relax\allowbreak{}\char32\relax\allowbreak{}\char32\relax\allowbreak{}\char32\relax\allowbreak{}\char32\relax\allowbreak{}i\allowbreak{}d\allowbreak{}s\allowbreak{}\char32\relax\allowbreak{}=\allowbreak{}\char32\relax\allowbreak{}l\allowbreak{}i\allowbreak{}s\allowbreak{}t\allowbreak{}(\allowbreak{}t\allowbreak{}e\allowbreak{}x\allowbreak{}t\allowbreak{}\char95\relax\allowbreak{}b\allowbreak{}y\allowbreak{}t\allowbreak{}e\allowbreak{}s\allowbreak{})\allowbreak{}\char32\relax\allowbreak{}\char35\relax\allowbreak{}\char32\relax\allowbreak{}l\allowbreak{}i\allowbreak{}s\allowbreak{}t\allowbreak{}\char32\relax\allowbreak{}o\allowbreak{}f\allowbreak{}\char32\relax\allowbreak{}i\allowbreak{}n\allowbreak{}t\allowbreak{}e\allowbreak{}g\allowbreak{}e\allowbreak{}r\allowbreak{}s\allowbreak{}\char32\relax\allowbreak{}i\allowbreak{}n\allowbreak{}\char32\relax\allowbreak{}r\allowbreak{}a\allowbreak{}n\allowbreak{}g\allowbreak{}e\allowbreak{}\char32\relax\allowbreak{}0\allowbreak{}.\allowbreak{}.\allowbreak{}2\allowbreak{}5\allowbreak{}5} & 把字节串转换成整数列表，初始每项都在 0–255 之间。 \\[3pt]
\SourceLineNumber{29} & \BpeDocCodeCell{\char32\relax\allowbreak{}\char32\relax\allowbreak{}\char32\relax\allowbreak{}\char32\relax\allowbreak{}\char32\relax\allowbreak{}\char32\relax\allowbreak{}\char32\relax\allowbreak{}\char32\relax\allowbreak{}m\allowbreak{}e\allowbreak{}r\allowbreak{}g\allowbreak{}e\allowbreak{}s\allowbreak{}\char32\relax\allowbreak{}=\allowbreak{}\char32\relax\allowbreak{}\char123\relax\allowbreak{}\char125\relax\allowbreak{}\char32\relax\allowbreak{}\char35\relax\allowbreak{}\char32\relax\allowbreak{}(\allowbreak{}i\allowbreak{}n\allowbreak{}t\allowbreak{},\allowbreak{}\char32\relax\allowbreak{}i\allowbreak{}n\allowbreak{}t\allowbreak{})\allowbreak{}\char32\relax\allowbreak{}-\allowbreak{}>\allowbreak{}\char32\relax\allowbreak{}i\allowbreak{}n\allowbreak{}t} & 新建合并表，记录“两个旧编号 \(\rightarrow\) 一个新编号”。 \\[3pt]
\SourceLineNumber{30} & \BpeDocCodeCell{\char32\relax\allowbreak{}\char32\relax\allowbreak{}\char32\relax\allowbreak{}\char32\relax\allowbreak{}\char32\relax\allowbreak{}\char32\relax\allowbreak{}\char32\relax\allowbreak{}\char32\relax\allowbreak{}v\allowbreak{}o\allowbreak{}c\allowbreak{}a\allowbreak{}b\allowbreak{}\char32\relax\allowbreak{}=\allowbreak{}\char32\relax\allowbreak{}\char123\relax\allowbreak{}i\allowbreak{}d\allowbreak{}x\allowbreak{}:\allowbreak{}\char32\relax\allowbreak{}b\allowbreak{}y\allowbreak{}t\allowbreak{}e\allowbreak{}s\allowbreak{}(\allowbreak{}[\allowbreak{}i\allowbreak{}d\allowbreak{}x\allowbreak{}]\allowbreak{})\allowbreak{}\char32\relax\allowbreak{}f\allowbreak{}o\allowbreak{}r\allowbreak{}\char32\relax\allowbreak{}i\allowbreak{}d\allowbreak{}x\allowbreak{}\char32\relax\allowbreak{}i\allowbreak{}n\allowbreak{}\char32\relax\allowbreak{}r\allowbreak{}a\allowbreak{}n\allowbreak{}g\allowbreak{}e\allowbreak{}(\allowbreak{}2\allowbreak{}5\allowbreak{}6\allowbreak{})\allowbreak{}\char125\relax\allowbreak{}\char32\relax\allowbreak{}\char35\relax\allowbreak{}\char32\relax\allowbreak{}i\allowbreak{}n\allowbreak{}t\allowbreak{}\char32\relax\allowbreak{}-\allowbreak{}>\allowbreak{}\char32\relax\allowbreak{}b\allowbreak{}y\allowbreak{}t\allowbreak{}e\allowbreak{}s} & 建立初始词表：每个 0–255 的编号对应一个单字节串。 \\[3pt]
\SourceLineNumber{31} & \BpeDocCodeCell{\char32\relax\allowbreak{}\char32\relax\allowbreak{}\char32\relax\allowbreak{}\char32\relax\allowbreak{}\char32\relax\allowbreak{}\char32\relax\allowbreak{}\char32\relax\allowbreak{}\char32\relax\allowbreak{}f\allowbreak{}o\allowbreak{}r\allowbreak{}\char32\relax\allowbreak{}i\allowbreak{}\char32\relax\allowbreak{}i\allowbreak{}n\allowbreak{}\char32\relax\allowbreak{}r\allowbreak{}a\allowbreak{}n\allowbreak{}g\allowbreak{}e\allowbreak{}(\allowbreak{}n\allowbreak{}u\allowbreak{}m\allowbreak{}\char95\relax\allowbreak{}m\allowbreak{}e\allowbreak{}r\allowbreak{}g\allowbreak{}e\allowbreak{}s\allowbreak{})\allowbreak{}:} & 按预定次数循环；i 从 0 开始，每轮产生一个新 \BpeDocCodeInline{token}。 \\[3pt]
\SourceLineNumber{33} & \BpeDocCodeCell{\char32\relax\allowbreak{}\char32\relax\allowbreak{}\char32\relax\allowbreak{}\char32\relax\allowbreak{}\char32\relax\allowbreak{}\char32\relax\allowbreak{}\char32\relax\allowbreak{}\char32\relax\allowbreak{}\char32\relax\allowbreak{}\char32\relax\allowbreak{}\char32\relax\allowbreak{}\char32\relax\allowbreak{}s\allowbreak{}t\allowbreak{}a\allowbreak{}t\allowbreak{}s\allowbreak{}\char32\relax\allowbreak{}=\allowbreak{}\char32\relax\allowbreak{}g\allowbreak{}e\allowbreak{}t\allowbreak{}\char95\relax\allowbreak{}s\allowbreak{}t\allowbreak{}a\allowbreak{}t\allowbreak{}s\allowbreak{}(\allowbreak{}i\allowbreak{}d\allowbreak{}s\allowbreak{})} & 对当前 \BpeDocCodeInline{ids} 统计所有相邻对的频次。每轮都重新统计，因为序列已经改变。 \\[3pt]
\SourceLineNumber{35} & \BpeDocCodeCell{\char32\relax\allowbreak{}\char32\relax\allowbreak{}\char32\relax\allowbreak{}\char32\relax\allowbreak{}\char32\relax\allowbreak{}\char32\relax\allowbreak{}\char32\relax\allowbreak{}\char32\relax\allowbreak{}\char32\relax\allowbreak{}\char32\relax\allowbreak{}\char32\relax\allowbreak{}\char32\relax\allowbreak{}p\allowbreak{}a\allowbreak{}i\allowbreak{}r\allowbreak{}\char32\relax\allowbreak{}=\allowbreak{}\char32\relax\allowbreak{}m\allowbreak{}a\allowbreak{}x\allowbreak{}(\allowbreak{}s\allowbreak{}t\allowbreak{}a\allowbreak{}t\allowbreak{}s\allowbreak{},\allowbreak{}\char32\relax\allowbreak{}k\allowbreak{}e\allowbreak{}y\allowbreak{}=\allowbreak{}s\allowbreak{}t\allowbreak{}a\allowbreak{}t\allowbreak{}s\allowbreak{}.\allowbreak{}g\allowbreak{}e\allowbreak{}t\allowbreak{})} & 按 \BpeDocCodeInline{stats} 中的计数选最大的 \BpeDocCodeInline{pair}。若并列，这段代码采用字典遍历中先遇到者。 \\[3pt]
\SourceLineNumber{37} & \BpeDocCodeCell{\char32\relax\allowbreak{}\char32\relax\allowbreak{}\char32\relax\allowbreak{}\char32\relax\allowbreak{}\char32\relax\allowbreak{}\char32\relax\allowbreak{}\char32\relax\allowbreak{}\char32\relax\allowbreak{}\char32\relax\allowbreak{}\char32\relax\allowbreak{}\char32\relax\allowbreak{}\char32\relax\allowbreak{}i\allowbreak{}d\allowbreak{}x\allowbreak{}\char32\relax\allowbreak{}=\allowbreak{}\char32\relax\allowbreak{}2\allowbreak{}5\allowbreak{}6\allowbreak{}\char32\relax\allowbreak{}+\allowbreak{}\char32\relax\allowbreak{}i} & 新编号从 256 起，随后依次为 257、258……；它也记录规则产生的先后顺序。 \\[3pt]
\SourceLineNumber{39} & \BpeDocCodeCell{\char32\relax\allowbreak{}\char32\relax\allowbreak{}\char32\relax\allowbreak{}\char32\relax\allowbreak{}\char32\relax\allowbreak{}\char32\relax\allowbreak{}\char32\relax\allowbreak{}\char32\relax\allowbreak{}\char32\relax\allowbreak{}\char32\relax\allowbreak{}\char32\relax\allowbreak{}\char32\relax\allowbreak{}i\allowbreak{}d\allowbreak{}s\allowbreak{}\char32\relax\allowbreak{}=\allowbreak{}\char32\relax\allowbreak{}m\allowbreak{}e\allowbreak{}r\allowbreak{}g\allowbreak{}e\allowbreak{}(\allowbreak{}i\allowbreak{}d\allowbreak{}s\allowbreak{},\allowbreak{}\char32\relax\allowbreak{}p\allowbreak{}a\allowbreak{}i\allowbreak{}r\allowbreak{},\allowbreak{}\char32\relax\allowbreak{}i\allowbreak{}d\allowbreak{}x\allowbreak{})} & 将当前序列中可从左到右配对的 \BpeDocCodeInline{pair} 替换为 \BpeDocCodeInline{idx}，并更新 \BpeDocCodeInline{ids}。 \\[3pt]
\SourceLineNumber{41} & \BpeDocCodeCell{\char32\relax\allowbreak{}\char32\relax\allowbreak{}\char32\relax\allowbreak{}\char32\relax\allowbreak{}\char32\relax\allowbreak{}\char32\relax\allowbreak{}\char32\relax\allowbreak{}\char32\relax\allowbreak{}\char32\relax\allowbreak{}\char32\relax\allowbreak{}\char32\relax\allowbreak{}\char32\relax\allowbreak{}m\allowbreak{}e\allowbreak{}r\allowbreak{}g\allowbreak{}e\allowbreak{}s\allowbreak{}[\allowbreak{}p\allowbreak{}a\allowbreak{}i\allowbreak{}r\allowbreak{}]\allowbreak{}\char32\relax\allowbreak{}=\allowbreak{}\char32\relax\allowbreak{}i\allowbreak{}d\allowbreak{}x} & 把新规则写入合并表，供以后编码新文本时使用。 \\[3pt]
\SourceLineNumber{42} & \BpeDocCodeCell{\char32\relax\allowbreak{}\char32\relax\allowbreak{}\char32\relax\allowbreak{}\char32\relax\allowbreak{}\char32\relax\allowbreak{}\char32\relax\allowbreak{}\char32\relax\allowbreak{}\char32\relax\allowbreak{}\char32\relax\allowbreak{}\char32\relax\allowbreak{}\char32\relax\allowbreak{}\char32\relax\allowbreak{}v\allowbreak{}o\allowbreak{}c\allowbreak{}a\allowbreak{}b\allowbreak{}[\allowbreak{}i\allowbreak{}d\allowbreak{}x\allowbreak{}]\allowbreak{}\char32\relax\allowbreak{}=\allowbreak{}\char32\relax\allowbreak{}v\allowbreak{}o\allowbreak{}c\allowbreak{}a\allowbreak{}b\allowbreak{}[\allowbreak{}p\allowbreak{}a\allowbreak{}i\allowbreak{}r\allowbreak{}[\allowbreak{}0\allowbreak{}]\allowbreak{}]\allowbreak{}\char32\relax\allowbreak{}+\allowbreak{}\char32\relax\allowbreak{}v\allowbreak{}o\allowbreak{}c\allowbreak{}a\allowbreak{}b\allowbreak{}[\allowbreak{}p\allowbreak{}a\allowbreak{}i\allowbreak{}r\allowbreak{}[\allowbreak{}1\allowbreak{}]\allowbreak{}]} & 把两个旧 \BpeDocCodeInline{token} 对应的字节串相接，得到新 \BpeDocCodeInline{token} 的实际字节内容。这里的 + 是拼接字节串。 \\[3pt]
\SourceLineNumber{44} & \BpeDocCodeCell{\char32\relax\allowbreak{}\char32\relax\allowbreak{}\char32\relax\allowbreak{}\char32\relax\allowbreak{}\char32\relax\allowbreak{}\char32\relax\allowbreak{}\char32\relax\allowbreak{}\char32\relax\allowbreak{}\char32\relax\allowbreak{}\char32\relax\allowbreak{}\char32\relax\allowbreak{}\char32\relax\allowbreak{}i\allowbreak{}f\allowbreak{}\char32\relax\allowbreak{}v\allowbreak{}e\allowbreak{}r\allowbreak{}b\allowbreak{}o\allowbreak{}s\allowbreak{}e\allowbreak{}:} & 只有 \BpeDocCodeInline{verbose} 为真时，才执行下一行打印。 \\[3pt]
\SourceLineNumber{45} & \BpeDocCodeCell{\char32\relax\allowbreak{}\char32\relax\allowbreak{}\char32\relax\allowbreak{}\char32\relax\allowbreak{}\char32\relax\allowbreak{}\char32\relax\allowbreak{}\char32\relax\allowbreak{}\char32\relax\allowbreak{}\char32\relax\allowbreak{}\char32\relax\allowbreak{}\char32\relax\allowbreak{}\char32\relax\allowbreak{}\char32\relax\allowbreak{}\char32\relax\allowbreak{}\char32\relax\allowbreak{}\char32\relax\allowbreak{}p\allowbreak{}r\allowbreak{}i\allowbreak{}n\allowbreak{}t\allowbreak{}(\allowbreak{}f\allowbreak{}"\allowbreak{}m\allowbreak{}e\allowbreak{}r\allowbreak{}g\allowbreak{}e\allowbreak{}\char32\relax\allowbreak{}\char123\relax\allowbreak{}i\allowbreak{}+\allowbreak{}1\allowbreak{}\char125\relax\allowbreak{}/\allowbreak{}\char123\relax\allowbreak{}n\allowbreak{}u\allowbreak{}m\allowbreak{}\char95\relax\allowbreak{}m\allowbreak{}e\allowbreak{}r\allowbreak{}g\allowbreak{}e\allowbreak{}s\allowbreak{}\char125\relax\allowbreak{}:\allowbreak{}\char32\relax\allowbreak{}\char123\relax\allowbreak{}p\allowbreak{}a\allowbreak{}i\allowbreak{}r\allowbreak{}\char125\relax\allowbreak{}\char32\relax\allowbreak{}-\allowbreak{}>\allowbreak{}\char32\relax\allowbreak{}\char123\relax\allowbreak{}i\allowbreak{}d\allowbreak{}x\allowbreak{}\char125\relax\allowbreak{}\char32\relax\allowbreak{}(\allowbreak{}\char123\relax\allowbreak{}v\allowbreak{}o\allowbreak{}c\allowbreak{}a\allowbreak{}b\allowbreak{}[\allowbreak{}i\allowbreak{}d\allowbreak{}x\allowbreak{}]\allowbreak{}\char125\relax\allowbreak{})\allowbreak{}\char32\relax\allowbreak{}h\allowbreak{}a\allowbreak{}d\allowbreak{}\char32\relax\allowbreak{}\char123\relax\allowbreak{}s\allowbreak{}t\allowbreak{}a\allowbreak{}t\allowbreak{}s\allowbreak{}[\allowbreak{}p\allowbreak{}a\allowbreak{}i\allowbreak{}r\allowbreak{}]\allowbreak{}\char125\relax\allowbreak{}\char32\relax\allowbreak{}o\allowbreak{}c\allowbreak{}c\allowbreak{}u\allowbreak{}r\allowbreak{}r\allowbreak{}e\allowbreak{}n\allowbreak{}c\allowbreak{}e\allowbreak{}s\allowbreak{}"\allowbreak{})} & 打印轮次、合并对、新编号、字节内容和相邻窗口计数；最后这个计数不一定等于实际替换次数。 \\[3pt]
\SourceLineNumber{48} & \BpeDocCodeCell{\char32\relax\allowbreak{}\char32\relax\allowbreak{}\char32\relax\allowbreak{}\char32\relax\allowbreak{}\char32\relax\allowbreak{}\char32\relax\allowbreak{}\char32\relax\allowbreak{}\char32\relax\allowbreak{}s\allowbreak{}e\allowbreak{}l\allowbreak{}f\allowbreak{}.\allowbreak{}m\allowbreak{}e\allowbreak{}r\allowbreak{}g\allowbreak{}e\allowbreak{}s\allowbreak{}\char32\relax\allowbreak{}=\allowbreak{}\char32\relax\allowbreak{}m\allowbreak{}e\allowbreak{}r\allowbreak{}g\allowbreak{}e\allowbreak{}s\allowbreak{}\char32\relax\allowbreak{}\char35\relax\allowbreak{}\char32\relax\allowbreak{}u\allowbreak{}s\allowbreak{}e\allowbreak{}d\allowbreak{}\char32\relax\allowbreak{}i\allowbreak{}n\allowbreak{}\char32\relax\allowbreak{}e\allowbreak{}n\allowbreak{}c\allowbreak{}o\allowbreak{}d\allowbreak{}e\allowbreak{}(\allowbreak{})} & 把学到的合并表交给当前分词器，后续 \BpeDocCodeInline{encode} 会读取它。 \\[3pt]
\SourceLineNumber{49} & \BpeDocCodeCell{\char32\relax\allowbreak{}\char32\relax\allowbreak{}\char32\relax\allowbreak{}\char32\relax\allowbreak{}\char32\relax\allowbreak{}\char32\relax\allowbreak{}\char32\relax\allowbreak{}\char32\relax\allowbreak{}s\allowbreak{}e\allowbreak{}l\allowbreak{}f\allowbreak{}.\allowbreak{}v\allowbreak{}o\allowbreak{}c\allowbreak{}a\allowbreak{}b\allowbreak{}\char32\relax\allowbreak{}=\allowbreak{}\char32\relax\allowbreak{}v\allowbreak{}o\allowbreak{}c\allowbreak{}a\allowbreak{}b\allowbreak{}\char32\relax\allowbreak{}\char32\relax\allowbreak{}\char32\relax\allowbreak{}\char35\relax\allowbreak{}\char32\relax\allowbreak{}u\allowbreak{}s\allowbreak{}e\allowbreak{}d\allowbreak{}\char32\relax\allowbreak{}i\allowbreak{}n\allowbreak{}\char32\relax\allowbreak{}d\allowbreak{}e\allowbreak{}c\allowbreak{}o\allowbreak{}d\allowbreak{}e\allowbreak{}(\allowbreak{})} & 把完整词表交给当前分词器，后续 \BpeDocCodeInline{decode} 会读取它。 \\[3pt]
\end{longtable}
\endgroup


这里的“训练”没有神经网络反向传播，只是反复统计与合并。代码也没有在相邻对耗尽时自动提前结束：如果文本太短而目标词表设得过大，\BpeDocCodeInline{stats} 可能为空，\BpeDocCodeInline{max} 会失败。因此小例子只设 259，而不是直接要求训练出几万个 \BpeDocCodeInline{token}。实际实验应准备足够语料，并检查停止条件。


\Needspace{18\baselineskip}
\BpeDocSubsubsection{\BpeDocCodeInline{basic.\allowbreak{}py}：\BpeDocCodeInline{decode} 把编号还原成文字，51–55 行}{basic.py：decode 把编号还原成文字，51–55 行}


\Needspace{4\baselineskip}
\begingroup\fontsize{8.9}{12.2}\selectfont\setstretch{1}
\setlength{\BpeDocTableWidth}{\dimexpr\linewidth-6\tabcolsep\relax}
\renewcommand{\arraystretch}{1.08}
\par\addvspace{0.45em}
\noindent\begin{minipage}{\linewidth}
{\small \BpeDocCodeInline{minbpe/\allowbreak{}basic.\allowbreak{}py}，第 51–55 行（仅列执行语句）}\par\vspace{0.35em}
\begin{tabular}{>{\raggedleft\arraybackslash}p{0.065\BpeDocTableWidth}>{\raggedright\arraybackslash}p{0.475\BpeDocTableWidth}>{\raggedright\arraybackslash}p{0.460\BpeDocTableWidth}}
\toprule
\textbf{行号} & \textbf{原代码} & \textbf{逐行解释} \\
\midrule
\SourceLineNumber{51} & \BpeDocCodeCell{\char32\relax\allowbreak{}\char32\relax\allowbreak{}\char32\relax\allowbreak{}\char32\relax\allowbreak{}d\allowbreak{}e\allowbreak{}f\allowbreak{}\char32\relax\allowbreak{}d\allowbreak{}e\allowbreak{}c\allowbreak{}o\allowbreak{}d\allowbreak{}e\allowbreak{}(\allowbreak{}s\allowbreak{}e\allowbreak{}l\allowbreak{}f\allowbreak{},\allowbreak{}\char32\relax\allowbreak{}i\allowbreak{}d\allowbreak{}s\allowbreak{})\allowbreak{}:} & 定义解码方法，输入是一串 \BpeDocCodeInline{token} ID。 \\[3pt]
\SourceLineNumber{53} & \BpeDocCodeCell{\char32\relax\allowbreak{}\char32\relax\allowbreak{}\char32\relax\allowbreak{}\char32\relax\allowbreak{}\char32\relax\allowbreak{}\char32\relax\allowbreak{}\char32\relax\allowbreak{}\char32\relax\allowbreak{}t\allowbreak{}e\allowbreak{}x\allowbreak{}t\allowbreak{}\char95\relax\allowbreak{}b\allowbreak{}y\allowbreak{}t\allowbreak{}e\allowbreak{}s\allowbreak{}\char32\relax\allowbreak{}=\allowbreak{}\char32\relax\allowbreak{}b\allowbreak{}"\allowbreak{}"\allowbreak{}.\allowbreak{}j\allowbreak{}o\allowbreak{}i\allowbreak{}n\allowbreak{}(\allowbreak{}s\allowbreak{}e\allowbreak{}l\allowbreak{}f\allowbreak{}.\allowbreak{}v\allowbreak{}o\allowbreak{}c\allowbreak{}a\allowbreak{}b\allowbreak{}[\allowbreak{}i\allowbreak{}d\allowbreak{}x\allowbreak{}]\allowbreak{}\char32\relax\allowbreak{}f\allowbreak{}o\allowbreak{}r\allowbreak{}\char32\relax\allowbreak{}i\allowbreak{}d\allowbreak{}x\allowbreak{}\char32\relax\allowbreak{}i\allowbreak{}n\allowbreak{}\char32\relax\allowbreak{}i\allowbreak{}d\allowbreak{}s\allowbreak{})} & 依次按编号查词表，取出各 \BpeDocCodeInline{token} 的字节内容，用空字节串连接成完整字节序列。 \\[3pt]
\SourceLineNumber{54} & \BpeDocCodeCell{\char32\relax\allowbreak{}\char32\relax\allowbreak{}\char32\relax\allowbreak{}\char32\relax\allowbreak{}\char32\relax\allowbreak{}\char32\relax\allowbreak{}\char32\relax\allowbreak{}\char32\relax\allowbreak{}t\allowbreak{}e\allowbreak{}x\allowbreak{}t\allowbreak{}\char32\relax\allowbreak{}=\allowbreak{}\char32\relax\allowbreak{}t\allowbreak{}e\allowbreak{}x\allowbreak{}t\allowbreak{}\char95\relax\allowbreak{}b\allowbreak{}y\allowbreak{}t\allowbreak{}e\allowbreak{}s\allowbreak{}.\allowbreak{}d\allowbreak{}e\allowbreak{}c\allowbreak{}o\allowbreak{}d\allowbreak{}e\allowbreak{}(\allowbreak{}"\allowbreak{}u\allowbreak{}t\allowbreak{}f\allowbreak{}-\allowbreak{}8\allowbreak{}"\allowbreak{},\allowbreak{}\char32\relax\allowbreak{}e\allowbreak{}r\allowbreak{}r\allowbreak{}o\allowbreak{}r\allowbreak{}s\allowbreak{}=\allowbreak{}"\allowbreak{}r\allowbreak{}e\allowbreak{}p\allowbreak{}l\allowbreak{}a\allowbreak{}c\allowbreak{}e\allowbreak{}"\allowbreak{})} & 将完整字节序列按 UTF-8 还原为文字；非法字节序列用替代字符表示。 \\[3pt]
\SourceLineNumber{55} & \BpeDocCodeCell{\char32\relax\allowbreak{}\char32\relax\allowbreak{}\char32\relax\allowbreak{}\char32\relax\allowbreak{}\char32\relax\allowbreak{}\char32\relax\allowbreak{}\char32\relax\allowbreak{}\char32\relax\allowbreak{}r\allowbreak{}e\allowbreak{}t\allowbreak{}u\allowbreak{}r\allowbreak{}n\allowbreak{}\char32\relax\allowbreak{}t\allowbreak{}e\allowbreak{}x\allowbreak{}t} & 返回还原后的字符串。 \\[3pt]
\bottomrule
\end{tabular}
\end{minipage}\par\addvspace{0.65em}
\endgroup


必须先拼接字节，再解码成文字。原因是一个 \BpeDocCodeInline{token} 可能只保存某个汉字 UTF-8 编码的一部分，单独解码它未必合法。对合法文本正常 \BpeDocCodeInline{encode} 后再 \BpeDocCodeInline{decode}，通常可以恢复原文；任意拼凑 \BpeDocCodeInline{token} ID、截断序列或输入无效字节，则不保证仍是合法文字。未知 ID 查词表也会出错。


\Needspace{9\baselineskip}
\BpeDocSubsubsection{\BpeDocCodeInline{basic.\allowbreak{}py}：\BpeDocCodeInline{encode} 用固定规则处理新文本，57–74 行}{basic.py：encode 用固定规则处理新文本，57–74 行}


\Needspace{4\baselineskip}
\begingroup\fontsize{8.9}{12.2}\selectfont\setstretch{1}
\setlength{\BpeDocTableWidth}{\dimexpr\linewidth-6\tabcolsep\relax}
\renewcommand{\arraystretch}{1.08}
{\small \BpeDocCodeInline{minbpe/\allowbreak{}basic.\allowbreak{}py}，第 57–74 行（仅列执行语句）}

\begin{longtable}{>{\raggedleft\arraybackslash}p{0.065\BpeDocTableWidth}>{\raggedright\arraybackslash}p{0.475\BpeDocTableWidth}>{\raggedright\arraybackslash}p{0.460\BpeDocTableWidth}}
\toprule
\textbf{行号} & \textbf{原代码} & \textbf{逐行解释} \\
\midrule\endfirsthead
\toprule
\textbf{行号} & \textbf{原代码} & \textbf{逐行解释} \\
\midrule\endhead
\midrule\endfoot
\bottomrule\endlastfoot
\SourceLineNumber{57} & \BpeDocCodeCell{\char32\relax\allowbreak{}\char32\relax\allowbreak{}\char32\relax\allowbreak{}\char32\relax\allowbreak{}d\allowbreak{}e\allowbreak{}f\allowbreak{}\char32\relax\allowbreak{}e\allowbreak{}n\allowbreak{}c\allowbreak{}o\allowbreak{}d\allowbreak{}e\allowbreak{}(\allowbreak{}s\allowbreak{}e\allowbreak{}l\allowbreak{}f\allowbreak{},\allowbreak{}\char32\relax\allowbreak{}t\allowbreak{}e\allowbreak{}x\allowbreak{}t\allowbreak{})\allowbreak{}:} & 定义编码方法：输入普通字符串，返回 \BpeDocCodeInline{token} ID 列表。 \\[3pt]
\SourceLineNumber{59} & \BpeDocCodeCell{\char32\relax\allowbreak{}\char32\relax\allowbreak{}\char32\relax\allowbreak{}\char32\relax\allowbreak{}\char32\relax\allowbreak{}\char32\relax\allowbreak{}\char32\relax\allowbreak{}\char32\relax\allowbreak{}t\allowbreak{}e\allowbreak{}x\allowbreak{}t\allowbreak{}\char95\relax\allowbreak{}b\allowbreak{}y\allowbreak{}t\allowbreak{}e\allowbreak{}s\allowbreak{}\char32\relax\allowbreak{}=\allowbreak{}\char32\relax\allowbreak{}t\allowbreak{}e\allowbreak{}x\allowbreak{}t\allowbreak{}.\allowbreak{}e\allowbreak{}n\allowbreak{}c\allowbreak{}o\allowbreak{}d\allowbreak{}e\allowbreak{}(\allowbreak{}"\allowbreak{}u\allowbreak{}t\allowbreak{}f\allowbreak{}-\allowbreak{}8\allowbreak{}"\allowbreak{})\allowbreak{}\char32\relax\allowbreak{}\char35\relax\allowbreak{}\char32\relax\allowbreak{}r\allowbreak{}a\allowbreak{}w\allowbreak{}\char32\relax\allowbreak{}b\allowbreak{}y\allowbreak{}t\allowbreak{}e\allowbreak{}s} & 先将文本转换为 UTF-8 字节串。 \\[3pt]
\SourceLineNumber{60} & \BpeDocCodeCell{\char32\relax\allowbreak{}\char32\relax\allowbreak{}\char32\relax\allowbreak{}\char32\relax\allowbreak{}\char32\relax\allowbreak{}\char32\relax\allowbreak{}\char32\relax\allowbreak{}\char32\relax\allowbreak{}i\allowbreak{}d\allowbreak{}s\allowbreak{}\char32\relax\allowbreak{}=\allowbreak{}\char32\relax\allowbreak{}l\allowbreak{}i\allowbreak{}s\allowbreak{}t\allowbreak{}(\allowbreak{}t\allowbreak{}e\allowbreak{}x\allowbreak{}t\allowbreak{}\char95\relax\allowbreak{}b\allowbreak{}y\allowbreak{}t\allowbreak{}e\allowbreak{}s\allowbreak{})\allowbreak{}\char32\relax\allowbreak{}\char35\relax\allowbreak{}\char32\relax\allowbreak{}l\allowbreak{}i\allowbreak{}s\allowbreak{}t\allowbreak{}\char32\relax\allowbreak{}o\allowbreak{}f\allowbreak{}\char32\relax\allowbreak{}i\allowbreak{}n\allowbreak{}t\allowbreak{}e\allowbreak{}g\allowbreak{}e\allowbreak{}r\allowbreak{}s\allowbreak{}\char32\relax\allowbreak{}i\allowbreak{}n\allowbreak{}\char32\relax\allowbreak{}r\allowbreak{}a\allowbreak{}n\allowbreak{}g\allowbreak{}e\allowbreak{}\char32\relax\allowbreak{}0\allowbreak{}.\allowbreak{}.\allowbreak{}2\allowbreak{}5\allowbreak{}5} & 再把每个字节转换为 0–255 的初始整数编号。 \\[3pt]
\SourceLineNumber{61} & \BpeDocCodeCell{\char32\relax\allowbreak{}\char32\relax\allowbreak{}\char32\relax\allowbreak{}\char32\relax\allowbreak{}\char32\relax\allowbreak{}\char32\relax\allowbreak{}\char32\relax\allowbreak{}\char32\relax\allowbreak{}w\allowbreak{}h\allowbreak{}i\allowbreak{}l\allowbreak{}e\allowbreak{}\char32\relax\allowbreak{}l\allowbreak{}e\allowbreak{}n\allowbreak{}(\allowbreak{}i\allowbreak{}d\allowbreak{}s\allowbreak{})\allowbreak{}\char32\relax\allowbreak{}>\allowbreak{}=\allowbreak{}\char32\relax\allowbreak{}2\allowbreak{}:} & 至少有两个编号时才可能合并；只有 0 或 1 个时直接结束循环。 \\[3pt]
\SourceLineNumber{63} & \BpeDocCodeCell{\char32\relax\allowbreak{}\char32\relax\allowbreak{}\char32\relax\allowbreak{}\char32\relax\allowbreak{}\char32\relax\allowbreak{}\char32\relax\allowbreak{}\char32\relax\allowbreak{}\char32\relax\allowbreak{}\char32\relax\allowbreak{}\char32\relax\allowbreak{}\char32\relax\allowbreak{}\char32\relax\allowbreak{}s\allowbreak{}t\allowbreak{}a\allowbreak{}t\allowbreak{}s\allowbreak{}\char32\relax\allowbreak{}=\allowbreak{}\char32\relax\allowbreak{}g\allowbreak{}e\allowbreak{}t\allowbreak{}\char95\relax\allowbreak{}s\allowbreak{}t\allowbreak{}a\allowbreak{}t\allowbreak{}s\allowbreak{}(\allowbreak{}i\allowbreak{}d\allowbreak{}s\allowbreak{})} & 列出当前序列中出现的相邻对；此处虽然也计算次数，但下一行不使用次数作决策。 \\[3pt]
\SourceLineNumber{64} & \BpeDocCodeCell{\char32\relax\allowbreak{}\char32\relax\allowbreak{}\char32\relax\allowbreak{}\char32\relax\allowbreak{}\char32\relax\allowbreak{}\char32\relax\allowbreak{}\char32\relax\allowbreak{}\char32\relax\allowbreak{}\char32\relax\allowbreak{}\char32\relax\allowbreak{}\char32\relax\allowbreak{}\char32\relax\allowbreak{}p\allowbreak{}a\allowbreak{}i\allowbreak{}r\allowbreak{}\char32\relax\allowbreak{}=\allowbreak{}\char32\relax\allowbreak{}m\allowbreak{}i\allowbreak{}n\allowbreak{}(\allowbreak{}s\allowbreak{}t\allowbreak{}a\allowbreak{}t\allowbreak{}s\allowbreak{},\allowbreak{}\char32\relax\allowbreak{}k\allowbreak{}e\allowbreak{}y\allowbreak{}=\allowbreak{}l\allowbreak{}a\allowbreak{}m\allowbreak{}b\allowbreak{}d\allowbreak{}a\allowbreak{}\char32\relax\allowbreak{}p\allowbreak{}:\allowbreak{}\char32\relax\allowbreak{}s\allowbreak{}e\allowbreak{}l\allowbreak{}f\allowbreak{}.\allowbreak{}m\allowbreak{}e\allowbreak{}r\allowbreak{}g\allowbreak{}e\allowbreak{}s\allowbreak{}.\allowbreak{}g\allowbreak{}e\allowbreak{}t\allowbreak{}(\allowbreak{}p\allowbreak{},\allowbreak{}\char32\relax\allowbreak{}f\allowbreak{}l\allowbreak{}o\allowbreak{}a\allowbreak{}t\allowbreak{}(\allowbreak{}"\allowbreak{}i\allowbreak{}n\allowbreak{}f\allowbreak{}"\allowbreak{})\allowbreak{})\allowbreak{})} & 对每个相邻对 p 查已学合并编号，选编号最小者；没学过的 \BpeDocCodeInline{pair} 记为正无穷，使其排在可用规则之后。 \\[3pt]
\SourceLineNumber{69} & \BpeDocCodeCell{\char32\relax\allowbreak{}\char32\relax\allowbreak{}\char32\relax\allowbreak{}\char32\relax\allowbreak{}\char32\relax\allowbreak{}\char32\relax\allowbreak{}\char32\relax\allowbreak{}\char32\relax\allowbreak{}\char32\relax\allowbreak{}\char32\relax\allowbreak{}\char32\relax\allowbreak{}\char32\relax\allowbreak{}i\allowbreak{}f\allowbreak{}\char32\relax\allowbreak{}p\allowbreak{}a\allowbreak{}i\allowbreak{}r\allowbreak{}\char32\relax\allowbreak{}n\allowbreak{}o\allowbreak{}t\allowbreak{}\char32\relax\allowbreak{}i\allowbreak{}n\allowbreak{}\char32\relax\allowbreak{}s\allowbreak{}e\allowbreak{}l\allowbreak{}f\allowbreak{}.\allowbreak{}m\allowbreak{}e\allowbreak{}r\allowbreak{}g\allowbreak{}e\allowbreak{}s\allowbreak{}:} & 若选中的相邻对不在合并表，说明当前所有相邻对都没有可用合并规则。 \\[3pt]
\SourceLineNumber{70} & \BpeDocCodeCell{\char32\relax\allowbreak{}\char32\relax\allowbreak{}\char32\relax\allowbreak{}\char32\relax\allowbreak{}\char32\relax\allowbreak{}\char32\relax\allowbreak{}\char32\relax\allowbreak{}\char32\relax\allowbreak{}\char32\relax\allowbreak{}\char32\relax\allowbreak{}\char32\relax\allowbreak{}\char32\relax\allowbreak{}\char32\relax\allowbreak{}\char32\relax\allowbreak{}\char32\relax\allowbreak{}\char32\relax\allowbreak{}b\allowbreak{}r\allowbreak{}e\allowbreak{}a\allowbreak{}k\allowbreak{}\char32\relax\allowbreak{}\char35\relax\allowbreak{}\char32\relax\allowbreak{}n\allowbreak{}o\allowbreak{}t\allowbreak{}h\allowbreak{}i\allowbreak{}n\allowbreak{}g\allowbreak{}\char32\relax\allowbreak{}e\allowbreak{}l\allowbreak{}s\allowbreak{}e\allowbreak{}\char32\relax\allowbreak{}c\allowbreak{}a\allowbreak{}n\allowbreak{}\char32\relax\allowbreak{}b\allowbreak{}e\allowbreak{}\char32\relax\allowbreak{}m\allowbreak{}e\allowbreak{}r\allowbreak{}g\allowbreak{}e\allowbreak{}d\allowbreak{}\char32\relax\allowbreak{}a\allowbreak{}n\allowbreak{}y\allowbreak{}m\allowbreak{}o\allowbreak{}r\allowbreak{}e} & 结束循环，不再继续合并。 \\[3pt]
\SourceLineNumber{72} & \BpeDocCodeCell{\char32\relax\allowbreak{}\char32\relax\allowbreak{}\char32\relax\allowbreak{}\char32\relax\allowbreak{}\char32\relax\allowbreak{}\char32\relax\allowbreak{}\char32\relax\allowbreak{}\char32\relax\allowbreak{}\char32\relax\allowbreak{}\char32\relax\allowbreak{}\char32\relax\allowbreak{}\char32\relax\allowbreak{}i\allowbreak{}d\allowbreak{}x\allowbreak{}\char32\relax\allowbreak{}=\allowbreak{}\char32\relax\allowbreak{}s\allowbreak{}e\allowbreak{}l\allowbreak{}f\allowbreak{}.\allowbreak{}m\allowbreak{}e\allowbreak{}r\allowbreak{}g\allowbreak{}e\allowbreak{}s\allowbreak{}[\allowbreak{}p\allowbreak{}a\allowbreak{}i\allowbreak{}r\allowbreak{}]} & 读取这条既有规则对应的新 \BpeDocCodeInline{token} 编号。 \\[3pt]
\SourceLineNumber{73} & \BpeDocCodeCell{\char32\relax\allowbreak{}\char32\relax\allowbreak{}\char32\relax\allowbreak{}\char32\relax\allowbreak{}\char32\relax\allowbreak{}\char32\relax\allowbreak{}\char32\relax\allowbreak{}\char32\relax\allowbreak{}\char32\relax\allowbreak{}\char32\relax\allowbreak{}\char32\relax\allowbreak{}\char32\relax\allowbreak{}i\allowbreak{}d\allowbreak{}s\allowbreak{}\char32\relax\allowbreak{}=\allowbreak{}\char32\relax\allowbreak{}m\allowbreak{}e\allowbreak{}r\allowbreak{}g\allowbreak{}e\allowbreak{}(\allowbreak{}i\allowbreak{}d\allowbreak{}s\allowbreak{},\allowbreak{}\char32\relax\allowbreak{}p\allowbreak{}a\allowbreak{}i\allowbreak{}r\allowbreak{},\allowbreak{}\char32\relax\allowbreak{}i\allowbreak{}d\allowbreak{}x\allowbreak{})} & 执行不重叠替换，得到下一轮的序列。这里不创建新规则。 \\[3pt]
\SourceLineNumber{74} & \BpeDocCodeCell{\char32\relax\allowbreak{}\char32\relax\allowbreak{}\char32\relax\allowbreak{}\char32\relax\allowbreak{}\char32\relax\allowbreak{}\char32\relax\allowbreak{}\char32\relax\allowbreak{}\char32\relax\allowbreak{}r\allowbreak{}e\allowbreak{}t\allowbreak{}u\allowbreak{}r\allowbreak{}n\allowbreak{}\char32\relax\allowbreak{}i\allowbreak{}d\allowbreak{}s} & 返回最终 \BpeDocCodeInline{token} ID 列表。 \\[3pt]
\end{longtable}
\endgroup


训练阶段按“当前出现最多”来创造规则；编码阶段按“已有规则中最早学到”来应用规则。例如假设已有 ab\(\rightarrow\)256、bc\(\rightarrow\)257，处理 bcbcab 时，即使 bc 出现 2 次、ab 只出现 1 次，也先使用编号较小的 ab 规则。这正是 \BpeDocCodeInline{max}(\BpeDocCodeInline{stats}, key=\BpeDocCodeInline{stats.\allowbreak{}get}) 与 \BpeDocCodeInline{min}(..., key=\BpeDocCodeInline{lambda} ...) 的区别。不能把编码解释成“在每次用户输入里重新统计最常见对并训练”。


\Needspace{9\baselineskip}
\BpeDocSubsection{\BpeDocCodeInline{regex.\allowbreak{}py}：为什么还要在 BPE 前预切分？}{regex.py：为什么还要在 BPE 前预切分？}


\BpeDocCodeInline{RegexTokenizer} 先用正则表达式把文字分成若干块，然后在每一块内部独立做 BPE。正则表达式可以理解为一套按字符类型找片段的规则。它让字母、数字、标点、换行等先按预设模式分组；BPE 不能跨过这些分组边界随意合并。分组边界由具体模式决定，不等于简单按空格切词。


\begingroup\fontsize{9.3}{12.7}\selectfont\setstretch{1}
\setlength{\BpeDocTableWidth}{\dimexpr\linewidth-4\tabcolsep\relax}
\renewcommand{\arraystretch}{1.15}
\par\addvspace{0.45em}
\noindent\begin{minipage}{\linewidth}
\begin{tabular}{>{\raggedright\arraybackslash}p{0.290000\BpeDocTableWidth}>{\raggedright\arraybackslash}p{0.710000\BpeDocTableWidth}}
\toprule
\textbf{GPT-4 模式中的分支} & \textbf{通俗解释} \\

\midrule
'(?i:[sdmt]|ll|ve|re) & 匹配若干带撇号的英文缩写后缀，并局部忽略大小写。 \\[3pt]
[\textasciicircum{}\textbackslash{}r\textbackslash{}n\textbackslash{}p\{L\}\textbackslash{}p\{N\}]?+\textbackslash{}p\{L\}+ & 匹配连续字母，并允许前面附带一个符合条件的非字母数字字符；常见效果包括把前导空格与后面的词放在一块。 \\[3pt]
\textbackslash{}p\{N\}\{1,3\} & 把连续数字按每块 1–3 个数字匹配。 \\[3pt]
 ?[\textasciicircum{}\textbackslash{}s\textbackslash{}p\{L\}\textbackslash{}p\{N\}]++[\textbackslash{}r\textbackslash{}n]* & 匹配一段非空白、非字母、非数字字符，可含前导空格和末尾换行。 \\[3pt]
\textbackslash{}s*[\textbackslash{}r\textbackslash{}n] & 处理包含换行的空白片段。 \\[3pt]
\textbackslash{}s+(?!\textbackslash{}S) 与 \textbackslash{}s+ & 处理其余符合条件的空白。具体由模式顺序与匹配语义决定。 \\[3pt]
\bottomrule
\end{tabular}
\end{minipage}\par\addvspace{0.65em}
\endgroup


这些规则不是让分词器理解语法或词义，而是设定允许合并的范围。中文字符也可属于 Unicode 字母类别；不能把“字母分支”仅理解为 a 到 z。\BpeDocCodeInline{regex.\allowbreak{}py} 第 12 行导入的是第三方 regex 包并命名为 re，不是 Python 自带 re；这与模式中的 Unicode 属性写法有关。


\Needspace{18\baselineskip}
\BpeDocSubsubsection{\BpeDocCodeInline{regex.\allowbreak{}py}：初始化预切分与特殊 \BpeDocCodeInline{token} 状态，24–34 行}{regex.py：初始化预切分与特殊 token 状态，24–34 行}


\Needspace{4\baselineskip}
\begingroup\fontsize{8.9}{12.2}\selectfont\setstretch{1}
\setlength{\BpeDocTableWidth}{\dimexpr\linewidth-6\tabcolsep\relax}
\renewcommand{\arraystretch}{1.08}
\par\addvspace{0.45em}
\noindent\begin{minipage}{\linewidth}
{\small \BpeDocCodeInline{minbpe/\allowbreak{}regex.\allowbreak{}py}，第 24–34 行（仅列执行语句）}\par\vspace{0.35em}
\begin{tabular}{>{\raggedleft\arraybackslash}p{0.065\BpeDocTableWidth}>{\raggedright\arraybackslash}p{0.475\BpeDocTableWidth}>{\raggedright\arraybackslash}p{0.460\BpeDocTableWidth}}
\toprule
\textbf{行号} & \textbf{原代码} & \textbf{逐行解释} \\
\midrule
\SourceLineNumber{24} & \BpeDocCodeCell{\char32\relax\allowbreak{}\char32\relax\allowbreak{}\char32\relax\allowbreak{}\char32\relax\allowbreak{}d\allowbreak{}e\allowbreak{}f\allowbreak{}\char32\relax\allowbreak{}\char95\relax\allowbreak{}\char95\relax\allowbreak{}i\allowbreak{}n\allowbreak{}i\allowbreak{}t\allowbreak{}\char95\relax\allowbreak{}\char95\relax\allowbreak{}(\allowbreak{}s\allowbreak{}e\allowbreak{}l\allowbreak{}f\allowbreak{},\allowbreak{}\char32\relax\allowbreak{}p\allowbreak{}a\allowbreak{}t\allowbreak{}t\allowbreak{}e\allowbreak{}r\allowbreak{}n\allowbreak{}=\allowbreak{}N\allowbreak{}o\allowbreak{}n\allowbreak{}e\allowbreak{})\allowbreak{}:} & 定义初始化方法，可传入自定义正则模式；不传时使用默认值。 \\[3pt]
\SourceLineNumber{30} & \BpeDocCodeCell{\char32\relax\allowbreak{}\char32\relax\allowbreak{}\char32\relax\allowbreak{}\char32\relax\allowbreak{}\char32\relax\allowbreak{}\char32\relax\allowbreak{}\char32\relax\allowbreak{}\char32\relax\allowbreak{}s\allowbreak{}u\allowbreak{}p\allowbreak{}e\allowbreak{}r\allowbreak{}(\allowbreak{})\allowbreak{}.\allowbreak{}\char95\relax\allowbreak{}\char95\relax\allowbreak{}i\allowbreak{}n\allowbreak{}i\allowbreak{}t\allowbreak{}\char95\relax\allowbreak{}\char95\relax\allowbreak{}(\allowbreak{})} & 先完成公共 \BpeDocCodeInline{Tokenizer} 的初始化。 \\[3pt]
\SourceLineNumber{31} & \BpeDocCodeCell{\char32\relax\allowbreak{}\char32\relax\allowbreak{}\char32\relax\allowbreak{}\char32\relax\allowbreak{}\char32\relax\allowbreak{}\char32\relax\allowbreak{}\char32\relax\allowbreak{}\char32\relax\allowbreak{}s\allowbreak{}e\allowbreak{}l\allowbreak{}f\allowbreak{}.\allowbreak{}p\allowbreak{}a\allowbreak{}t\allowbreak{}t\allowbreak{}e\allowbreak{}r\allowbreak{}n\allowbreak{}\char32\relax\allowbreak{}=\allowbreak{}\char32\relax\allowbreak{}G\allowbreak{}P\allowbreak{}T\allowbreak{}4\allowbreak{}\char95\relax\allowbreak{}S\allowbreak{}P\allowbreak{}L\allowbreak{}I\allowbreak{}T\allowbreak{}\char95\relax\allowbreak{}P\allowbreak{}A\allowbreak{}T\allowbreak{}T\allowbreak{}E\allowbreak{}R\allowbreak{}N\allowbreak{}\char32\relax\allowbreak{}i\allowbreak{}f\allowbreak{}\char32\relax\allowbreak{}p\allowbreak{}a\allowbreak{}t\allowbreak{}t\allowbreak{}e\allowbreak{}r\allowbreak{}n\allowbreak{}\char32\relax\allowbreak{}i\allowbreak{}s\allowbreak{}\char32\relax\allowbreak{}N\allowbreak{}o\allowbreak{}n\allowbreak{}e\allowbreak{}\char32\relax\allowbreak{}e\allowbreak{}l\allowbreak{}s\allowbreak{}e\allowbreak{}\char32\relax\allowbreak{}p\allowbreak{}a\allowbreak{}t\allowbreak{}t\allowbreak{}e\allowbreak{}r\allowbreak{}n} & 若 \BpeDocCodeInline{pattern} 没传，就选 \BpeDocCodeInline{GPT4\char95\relax\allowbreak{}SPLIT\char95\relax\allowbreak{}PATTERN}；否则使用调用者提供的模式。 \\[3pt]
\SourceLineNumber{32} & \BpeDocCodeCell{\char32\relax\allowbreak{}\char32\relax\allowbreak{}\char32\relax\allowbreak{}\char32\relax\allowbreak{}\char32\relax\allowbreak{}\char32\relax\allowbreak{}\char32\relax\allowbreak{}\char32\relax\allowbreak{}s\allowbreak{}e\allowbreak{}l\allowbreak{}f\allowbreak{}.\allowbreak{}c\allowbreak{}o\allowbreak{}m\allowbreak{}p\allowbreak{}i\allowbreak{}l\allowbreak{}e\allowbreak{}d\allowbreak{}\char95\relax\allowbreak{}p\allowbreak{}a\allowbreak{}t\allowbreak{}t\allowbreak{}e\allowbreak{}r\allowbreak{}n\allowbreak{}\char32\relax\allowbreak{}=\allowbreak{}\char32\relax\allowbreak{}r\allowbreak{}e\allowbreak{}.\allowbreak{}c\allowbreak{}o\allowbreak{}m\allowbreak{}p\allowbreak{}i\allowbreak{}l\allowbreak{}e\allowbreak{}(\allowbreak{}s\allowbreak{}e\allowbreak{}l\allowbreak{}f\allowbreak{}.\allowbreak{}p\allowbreak{}a\allowbreak{}t\allowbreak{}t\allowbreak{}e\allowbreak{}r\allowbreak{}n\allowbreak{})} & 把模式编译成可重复使用的正则对象。 \\[3pt]
\SourceLineNumber{33} & \BpeDocCodeCell{\char32\relax\allowbreak{}\char32\relax\allowbreak{}\char32\relax\allowbreak{}\char32\relax\allowbreak{}\char32\relax\allowbreak{}\char32\relax\allowbreak{}\char32\relax\allowbreak{}\char32\relax\allowbreak{}s\allowbreak{}e\allowbreak{}l\allowbreak{}f\allowbreak{}.\allowbreak{}s\allowbreak{}p\allowbreak{}e\allowbreak{}c\allowbreak{}i\allowbreak{}a\allowbreak{}l\allowbreak{}\char95\relax\allowbreak{}t\allowbreak{}o\allowbreak{}k\allowbreak{}e\allowbreak{}n\allowbreak{}s\allowbreak{}\char32\relax\allowbreak{}=\allowbreak{}\char32\relax\allowbreak{}\char123\relax\allowbreak{}\char125\relax} & 建立“特殊字符串 \(\rightarrow\) 编号”的空表。 \\[3pt]
\SourceLineNumber{34} & \BpeDocCodeCell{\char32\relax\allowbreak{}\char32\relax\allowbreak{}\char32\relax\allowbreak{}\char32\relax\allowbreak{}\char32\relax\allowbreak{}\char32\relax\allowbreak{}\char32\relax\allowbreak{}\char32\relax\allowbreak{}s\allowbreak{}e\allowbreak{}l\allowbreak{}f\allowbreak{}.\allowbreak{}i\allowbreak{}n\allowbreak{}v\allowbreak{}e\allowbreak{}r\allowbreak{}s\allowbreak{}e\allowbreak{}\char95\relax\allowbreak{}s\allowbreak{}p\allowbreak{}e\allowbreak{}c\allowbreak{}i\allowbreak{}a\allowbreak{}l\allowbreak{}\char95\relax\allowbreak{}t\allowbreak{}o\allowbreak{}k\allowbreak{}e\allowbreak{}n\allowbreak{}s\allowbreak{}\char32\relax\allowbreak{}=\allowbreak{}\char32\relax\allowbreak{}\char123\relax\allowbreak{}\char125\relax} & 建立“编号 \(\rightarrow\) 特殊字符串”的空表，供反向恢复时使用。 \\[3pt]
\bottomrule
\end{tabular}
\end{minipage}\par\addvspace{0.65em}
\endgroup


\Needspace{9\baselineskip}
\BpeDocSubsubsection{\BpeDocCodeInline{regex.\allowbreak{}py}：训练时共享统计，但不跨块合并，36–70 行}{regex.py：训练时共享统计，但不跨块合并，36–70 行}


\Needspace{4\baselineskip}
\begingroup\fontsize{8.9}{12.2}\selectfont\setstretch{1}
\setlength{\BpeDocTableWidth}{\dimexpr\linewidth-6\tabcolsep\relax}
\renewcommand{\arraystretch}{1.08}
{\small \BpeDocCodeInline{minbpe/\allowbreak{}regex.\allowbreak{}py}，第 36–70 行（仅列执行语句）}

\begin{longtable}{>{\raggedleft\arraybackslash}p{0.065\BpeDocTableWidth}>{\raggedright\arraybackslash}p{0.475\BpeDocTableWidth}>{\raggedright\arraybackslash}p{0.460\BpeDocTableWidth}}
\toprule
\textbf{行号} & \textbf{原代码} & \textbf{逐行解释} \\
\midrule\endfirsthead
\toprule
\textbf{行号} & \textbf{原代码} & \textbf{逐行解释} \\
\midrule\endhead
\midrule\endfoot
\bottomrule\endlastfoot
\SourceLineNumber{36} & \BpeDocCodeCell{\char32\relax\allowbreak{}\char32\relax\allowbreak{}\char32\relax\allowbreak{}\char32\relax\allowbreak{}d\allowbreak{}e\allowbreak{}f\allowbreak{}\char32\relax\allowbreak{}t\allowbreak{}r\allowbreak{}a\allowbreak{}i\allowbreak{}n\allowbreak{}(\allowbreak{}s\allowbreak{}e\allowbreak{}l\allowbreak{}f\allowbreak{},\allowbreak{}\char32\relax\allowbreak{}t\allowbreak{}e\allowbreak{}x\allowbreak{}t\allowbreak{},\allowbreak{}\char32\relax\allowbreak{}v\allowbreak{}o\allowbreak{}c\allowbreak{}a\allowbreak{}b\allowbreak{}\char95\relax\allowbreak{}s\allowbreak{}i\allowbreak{}z\allowbreak{}e\allowbreak{},\allowbreak{}\char32\relax\allowbreak{}v\allowbreak{}e\allowbreak{}r\allowbreak{}b\allowbreak{}o\allowbreak{}s\allowbreak{}e\allowbreak{}=\allowbreak{}F\allowbreak{}a\allowbreak{}l\allowbreak{}s\allowbreak{}e\allowbreak{})\allowbreak{}:} & 定义训练方法，输入训练文本、目标词表大小和打印开关。 \\[3pt]
\SourceLineNumber{37} & \BpeDocCodeCell{\char32\relax\allowbreak{}\char32\relax\allowbreak{}\char32\relax\allowbreak{}\char32\relax\allowbreak{}\char32\relax\allowbreak{}\char32\relax\allowbreak{}\char32\relax\allowbreak{}\char32\relax\allowbreak{}a\allowbreak{}s\allowbreak{}s\allowbreak{}e\allowbreak{}r\allowbreak{}t\allowbreak{}\char32\relax\allowbreak{}v\allowbreak{}o\allowbreak{}c\allowbreak{}a\allowbreak{}b\allowbreak{}\char95\relax\allowbreak{}s\allowbreak{}i\allowbreak{}z\allowbreak{}e\allowbreak{}\char32\relax\allowbreak{}>\allowbreak{}=\allowbreak{}\char32\relax\allowbreak{}2\allowbreak{}5\allowbreak{}6} & 要求至少保留 256 个单字节 \BpeDocCodeInline{token}。 \\[3pt]
\SourceLineNumber{38} & \BpeDocCodeCell{\char32\relax\allowbreak{}\char32\relax\allowbreak{}\char32\relax\allowbreak{}\char32\relax\allowbreak{}\char32\relax\allowbreak{}\char32\relax\allowbreak{}\char32\relax\allowbreak{}\char32\relax\allowbreak{}n\allowbreak{}u\allowbreak{}m\allowbreak{}\char95\relax\allowbreak{}m\allowbreak{}e\allowbreak{}r\allowbreak{}g\allowbreak{}e\allowbreak{}s\allowbreak{}\char32\relax\allowbreak{}=\allowbreak{}\char32\relax\allowbreak{}v\allowbreak{}o\allowbreak{}c\allowbreak{}a\allowbreak{}b\allowbreak{}\char95\relax\allowbreak{}s\allowbreak{}i\allowbreak{}z\allowbreak{}e\allowbreak{}\char32\relax\allowbreak{}-\allowbreak{}\char32\relax\allowbreak{}2\allowbreak{}5\allowbreak{}6} & 计算需要学习多少条新合并规则。 \\[3pt]
\SourceLineNumber{41} & \BpeDocCodeCell{\char32\relax\allowbreak{}\char32\relax\allowbreak{}\char32\relax\allowbreak{}\char32\relax\allowbreak{}\char32\relax\allowbreak{}\char32\relax\allowbreak{}\char32\relax\allowbreak{}\char32\relax\allowbreak{}t\allowbreak{}e\allowbreak{}x\allowbreak{}t\allowbreak{}\char95\relax\allowbreak{}c\allowbreak{}h\allowbreak{}u\allowbreak{}n\allowbreak{}k\allowbreak{}s\allowbreak{}\char32\relax\allowbreak{}=\allowbreak{}\char32\relax\allowbreak{}r\allowbreak{}e\allowbreak{}.\allowbreak{}f\allowbreak{}i\allowbreak{}n\allowbreak{}d\allowbreak{}a\allowbreak{}l\allowbreak{}l\allowbreak{}(\allowbreak{}s\allowbreak{}e\allowbreak{}l\allowbreak{}f\allowbreak{}.\allowbreak{}c\allowbreak{}o\allowbreak{}m\allowbreak{}p\allowbreak{}i\allowbreak{}l\allowbreak{}e\allowbreak{}d\allowbreak{}\char95\relax\allowbreak{}p\allowbreak{}a\allowbreak{}t\allowbreak{}t\allowbreak{}e\allowbreak{}r\allowbreak{}n\allowbreak{},\allowbreak{}\char32\relax\allowbreak{}t\allowbreak{}e\allowbreak{}x\allowbreak{}t\allowbreak{})} & 按编译后的正则模式找到全部文本块。 \\[3pt]
\SourceLineNumber{44} & \BpeDocCodeCell{\char32\relax\allowbreak{}\char32\relax\allowbreak{}\char32\relax\allowbreak{}\char32\relax\allowbreak{}\char32\relax\allowbreak{}\char32\relax\allowbreak{}\char32\relax\allowbreak{}\char32\relax\allowbreak{}i\allowbreak{}d\allowbreak{}s\allowbreak{}\char32\relax\allowbreak{}=\allowbreak{}\char32\relax\allowbreak{}[\allowbreak{}l\allowbreak{}i\allowbreak{}s\allowbreak{}t\allowbreak{}(\allowbreak{}c\allowbreak{}h\allowbreak{}.\allowbreak{}e\allowbreak{}n\allowbreak{}c\allowbreak{}o\allowbreak{}d\allowbreak{}e\allowbreak{}(\allowbreak{}"\allowbreak{}u\allowbreak{}t\allowbreak{}f\allowbreak{}-\allowbreak{}8\allowbreak{}"\allowbreak{})\allowbreak{})\allowbreak{}\char32\relax\allowbreak{}f\allowbreak{}o\allowbreak{}r\allowbreak{}\char32\relax\allowbreak{}c\allowbreak{}h\allowbreak{}\char32\relax\allowbreak{}i\allowbreak{}n\allowbreak{}\char32\relax\allowbreak{}t\allowbreak{}e\allowbreak{}x\allowbreak{}t\allowbreak{}\char95\relax\allowbreak{}c\allowbreak{}h\allowbreak{}u\allowbreak{}n\allowbreak{}k\allowbreak{}s\allowbreak{}]} & 每个块先 UTF-8 编码再转列表，得到“列表的列表”；块边界仍然保留。 \\[3pt]
\SourceLineNumber{47} & \BpeDocCodeCell{\char32\relax\allowbreak{}\char32\relax\allowbreak{}\char32\relax\allowbreak{}\char32\relax\allowbreak{}\char32\relax\allowbreak{}\char32\relax\allowbreak{}\char32\relax\allowbreak{}\char32\relax\allowbreak{}m\allowbreak{}e\allowbreak{}r\allowbreak{}g\allowbreak{}e\allowbreak{}s\allowbreak{}\char32\relax\allowbreak{}=\allowbreak{}\char32\relax\allowbreak{}\char123\relax\allowbreak{}\char125\relax\allowbreak{}\char32\relax\allowbreak{}\char35\relax\allowbreak{}\char32\relax\allowbreak{}(\allowbreak{}i\allowbreak{}n\allowbreak{}t\allowbreak{},\allowbreak{}\char32\relax\allowbreak{}i\allowbreak{}n\allowbreak{}t\allowbreak{})\allowbreak{}\char32\relax\allowbreak{}-\allowbreak{}>\allowbreak{}\char32\relax\allowbreak{}i\allowbreak{}n\allowbreak{}t} & 建立本次训练的新合并表。 \\[3pt]
\SourceLineNumber{48} & \BpeDocCodeCell{\char32\relax\allowbreak{}\char32\relax\allowbreak{}\char32\relax\allowbreak{}\char32\relax\allowbreak{}\char32\relax\allowbreak{}\char32\relax\allowbreak{}\char32\relax\allowbreak{}\char32\relax\allowbreak{}v\allowbreak{}o\allowbreak{}c\allowbreak{}a\allowbreak{}b\allowbreak{}\char32\relax\allowbreak{}=\allowbreak{}\char32\relax\allowbreak{}\char123\relax\allowbreak{}i\allowbreak{}d\allowbreak{}x\allowbreak{}:\allowbreak{}\char32\relax\allowbreak{}b\allowbreak{}y\allowbreak{}t\allowbreak{}e\allowbreak{}s\allowbreak{}(\allowbreak{}[\allowbreak{}i\allowbreak{}d\allowbreak{}x\allowbreak{}]\allowbreak{})\allowbreak{}\char32\relax\allowbreak{}f\allowbreak{}o\allowbreak{}r\allowbreak{}\char32\relax\allowbreak{}i\allowbreak{}d\allowbreak{}x\allowbreak{}\char32\relax\allowbreak{}i\allowbreak{}n\allowbreak{}\char32\relax\allowbreak{}r\allowbreak{}a\allowbreak{}n\allowbreak{}g\allowbreak{}e\allowbreak{}(\allowbreak{}2\allowbreak{}5\allowbreak{}6\allowbreak{})\allowbreak{}\char125\relax\allowbreak{}\char32\relax\allowbreak{}\char35\relax\allowbreak{}\char32\relax\allowbreak{}i\allowbreak{}d\allowbreak{}x\allowbreak{}\char32\relax\allowbreak{}-\allowbreak{}>\allowbreak{}\char32\relax\allowbreak{}b\allowbreak{}y\allowbreak{}t\allowbreak{}e\allowbreak{}s} & 建立初始 256 项字节词表。 \\[3pt]
\SourceLineNumber{49} & \BpeDocCodeCell{\char32\relax\allowbreak{}\char32\relax\allowbreak{}\char32\relax\allowbreak{}\char32\relax\allowbreak{}\char32\relax\allowbreak{}\char32\relax\allowbreak{}\char32\relax\allowbreak{}\char32\relax\allowbreak{}f\allowbreak{}o\allowbreak{}r\allowbreak{}\char32\relax\allowbreak{}i\allowbreak{}\char32\relax\allowbreak{}i\allowbreak{}n\allowbreak{}\char32\relax\allowbreak{}r\allowbreak{}a\allowbreak{}n\allowbreak{}g\allowbreak{}e\allowbreak{}(\allowbreak{}n\allowbreak{}u\allowbreak{}m\allowbreak{}\char95\relax\allowbreak{}m\allowbreak{}e\allowbreak{}r\allowbreak{}g\allowbreak{}e\allowbreak{}s\allowbreak{})\allowbreak{}:} & 每轮产生一个新 \BpeDocCodeInline{token}。 \\[3pt]
\SourceLineNumber{51} & \BpeDocCodeCell{\char32\relax\allowbreak{}\char32\relax\allowbreak{}\char32\relax\allowbreak{}\char32\relax\allowbreak{}\char32\relax\allowbreak{}\char32\relax\allowbreak{}\char32\relax\allowbreak{}\char32\relax\allowbreak{}\char32\relax\allowbreak{}\char32\relax\allowbreak{}\char32\relax\allowbreak{}\char32\relax\allowbreak{}s\allowbreak{}t\allowbreak{}a\allowbreak{}t\allowbreak{}s\allowbreak{}\char32\relax\allowbreak{}=\allowbreak{}\char32\relax\allowbreak{}\char123\relax\allowbreak{}\char125\relax} & 每一轮先清空总统计表。 \\[3pt]
\SourceLineNumber{52} & \BpeDocCodeCell{\char32\relax\allowbreak{}\char32\relax\allowbreak{}\char32\relax\allowbreak{}\char32\relax\allowbreak{}\char32\relax\allowbreak{}\char32\relax\allowbreak{}\char32\relax\allowbreak{}\char32\relax\allowbreak{}\char32\relax\allowbreak{}\char32\relax\allowbreak{}\char32\relax\allowbreak{}\char32\relax\allowbreak{}f\allowbreak{}o\allowbreak{}r\allowbreak{}\char32\relax\allowbreak{}c\allowbreak{}h\allowbreak{}u\allowbreak{}n\allowbreak{}k\allowbreak{}\char95\relax\allowbreak{}i\allowbreak{}d\allowbreak{}s\allowbreak{}\char32\relax\allowbreak{}i\allowbreak{}n\allowbreak{}\char32\relax\allowbreak{}i\allowbreak{}d\allowbreak{}s\allowbreak{}:} & 逐个取出文本块的编号列表。 \\[3pt]
\SourceLineNumber{54} & \BpeDocCodeCell{\char32\relax\allowbreak{}\char32\relax\allowbreak{}\char32\relax\allowbreak{}\char32\relax\allowbreak{}\char32\relax\allowbreak{}\char32\relax\allowbreak{}\char32\relax\allowbreak{}\char32\relax\allowbreak{}\char32\relax\allowbreak{}\char32\relax\allowbreak{}\char32\relax\allowbreak{}\char32\relax\allowbreak{}\char32\relax\allowbreak{}\char32\relax\allowbreak{}\char32\relax\allowbreak{}\char32\relax\allowbreak{}g\allowbreak{}e\allowbreak{}t\allowbreak{}\char95\relax\allowbreak{}s\allowbreak{}t\allowbreak{}a\allowbreak{}t\allowbreak{}s\allowbreak{}(\allowbreak{}c\allowbreak{}h\allowbreak{}u\allowbreak{}n\allowbreak{}k\allowbreak{}\char95\relax\allowbreak{}i\allowbreak{}d\allowbreak{}s\allowbreak{},\allowbreak{}\char32\relax\allowbreak{}s\allowbreak{}t\allowbreak{}a\allowbreak{}t\allowbreak{}s\allowbreak{})} & 累计当前块内部的相邻对；不会把前块末尾与后块开头算成一对。 \\[3pt]
\SourceLineNumber{56} & \BpeDocCodeCell{\char32\relax\allowbreak{}\char32\relax\allowbreak{}\char32\relax\allowbreak{}\char32\relax\allowbreak{}\char32\relax\allowbreak{}\char32\relax\allowbreak{}\char32\relax\allowbreak{}\char32\relax\allowbreak{}\char32\relax\allowbreak{}\char32\relax\allowbreak{}\char32\relax\allowbreak{}\char32\relax\allowbreak{}p\allowbreak{}a\allowbreak{}i\allowbreak{}r\allowbreak{}\char32\relax\allowbreak{}=\allowbreak{}\char32\relax\allowbreak{}m\allowbreak{}a\allowbreak{}x\allowbreak{}(\allowbreak{}s\allowbreak{}t\allowbreak{}a\allowbreak{}t\allowbreak{}s\allowbreak{},\allowbreak{}\char32\relax\allowbreak{}k\allowbreak{}e\allowbreak{}y\allowbreak{}=\allowbreak{}s\allowbreak{}t\allowbreak{}a\allowbreak{}t\allowbreak{}s\allowbreak{}.\allowbreak{}g\allowbreak{}e\allowbreak{}t\allowbreak{})} & 在各块累加后的统计中选频次最高的 \BpeDocCodeInline{pair}。 \\[3pt]
\SourceLineNumber{58} & \BpeDocCodeCell{\char32\relax\allowbreak{}\char32\relax\allowbreak{}\char32\relax\allowbreak{}\char32\relax\allowbreak{}\char32\relax\allowbreak{}\char32\relax\allowbreak{}\char32\relax\allowbreak{}\char32\relax\allowbreak{}\char32\relax\allowbreak{}\char32\relax\allowbreak{}\char32\relax\allowbreak{}\char32\relax\allowbreak{}i\allowbreak{}d\allowbreak{}x\allowbreak{}\char32\relax\allowbreak{}=\allowbreak{}\char32\relax\allowbreak{}2\allowbreak{}5\allowbreak{}6\allowbreak{}\char32\relax\allowbreak{}+\allowbreak{}\char32\relax\allowbreak{}i} & 新编号按 256、257……依次递增。 \\[3pt]
\SourceLineNumber{60} & \BpeDocCodeCell{\char32\relax\allowbreak{}\char32\relax\allowbreak{}\char32\relax\allowbreak{}\char32\relax\allowbreak{}\char32\relax\allowbreak{}\char32\relax\allowbreak{}\char32\relax\allowbreak{}\char32\relax\allowbreak{}\char32\relax\allowbreak{}\char32\relax\allowbreak{}\char32\relax\allowbreak{}\char32\relax\allowbreak{}i\allowbreak{}d\allowbreak{}s\allowbreak{}\char32\relax\allowbreak{}=\allowbreak{}\char32\relax\allowbreak{}[\allowbreak{}m\allowbreak{}e\allowbreak{}r\allowbreak{}g\allowbreak{}e\allowbreak{}(\allowbreak{}c\allowbreak{}h\allowbreak{}u\allowbreak{}n\allowbreak{}k\allowbreak{}\char95\relax\allowbreak{}i\allowbreak{}d\allowbreak{}s\allowbreak{},\allowbreak{}\char32\relax\allowbreak{}p\allowbreak{}a\allowbreak{}i\allowbreak{}r\allowbreak{},\allowbreak{}\char32\relax\allowbreak{}i\allowbreak{}d\allowbreak{}x\allowbreak{})\allowbreak{}\char32\relax\allowbreak{}f\allowbreak{}o\allowbreak{}r\allowbreak{}\char32\relax\allowbreak{}c\allowbreak{}h\allowbreak{}u\allowbreak{}n\allowbreak{}k\allowbreak{}\char95\relax\allowbreak{}i\allowbreak{}d\allowbreak{}s\allowbreak{}\char32\relax\allowbreak{}i\allowbreak{}n\allowbreak{}\char32\relax\allowbreak{}i\allowbreak{}d\allowbreak{}s\allowbreak{}]} & 对每块分别执行 \BpeDocCodeInline{merge}，保留块边界并替换整个 \BpeDocCodeInline{ids} 列表。 \\[3pt]
\SourceLineNumber{62} & \BpeDocCodeCell{\char32\relax\allowbreak{}\char32\relax\allowbreak{}\char32\relax\allowbreak{}\char32\relax\allowbreak{}\char32\relax\allowbreak{}\char32\relax\allowbreak{}\char32\relax\allowbreak{}\char32\relax\allowbreak{}\char32\relax\allowbreak{}\char32\relax\allowbreak{}\char32\relax\allowbreak{}\char32\relax\allowbreak{}m\allowbreak{}e\allowbreak{}r\allowbreak{}g\allowbreak{}e\allowbreak{}s\allowbreak{}[\allowbreak{}p\allowbreak{}a\allowbreak{}i\allowbreak{}r\allowbreak{}]\allowbreak{}\char32\relax\allowbreak{}=\allowbreak{}\char32\relax\allowbreak{}i\allowbreak{}d\allowbreak{}x} & 记录两个旧编号合成新编号的关系。 \\[3pt]
\SourceLineNumber{63} & \BpeDocCodeCell{\char32\relax\allowbreak{}\char32\relax\allowbreak{}\char32\relax\allowbreak{}\char32\relax\allowbreak{}\char32\relax\allowbreak{}\char32\relax\allowbreak{}\char32\relax\allowbreak{}\char32\relax\allowbreak{}\char32\relax\allowbreak{}\char32\relax\allowbreak{}\char32\relax\allowbreak{}\char32\relax\allowbreak{}v\allowbreak{}o\allowbreak{}c\allowbreak{}a\allowbreak{}b\allowbreak{}[\allowbreak{}i\allowbreak{}d\allowbreak{}x\allowbreak{}]\allowbreak{}\char32\relax\allowbreak{}=\allowbreak{}\char32\relax\allowbreak{}v\allowbreak{}o\allowbreak{}c\allowbreak{}a\allowbreak{}b\allowbreak{}[\allowbreak{}p\allowbreak{}a\allowbreak{}i\allowbreak{}r\allowbreak{}[\allowbreak{}0\allowbreak{}]\allowbreak{}]\allowbreak{}\char32\relax\allowbreak{}+\allowbreak{}\char32\relax\allowbreak{}v\allowbreak{}o\allowbreak{}c\allowbreak{}a\allowbreak{}b\allowbreak{}[\allowbreak{}p\allowbreak{}a\allowbreak{}i\allowbreak{}r\allowbreak{}[\allowbreak{}1\allowbreak{}]\allowbreak{}]} & 拼接旧 \BpeDocCodeInline{token} 的字节内容，生成新词表项。 \\[3pt]
\SourceLineNumber{65} & \BpeDocCodeCell{\char32\relax\allowbreak{}\char32\relax\allowbreak{}\char32\relax\allowbreak{}\char32\relax\allowbreak{}\char32\relax\allowbreak{}\char32\relax\allowbreak{}\char32\relax\allowbreak{}\char32\relax\allowbreak{}\char32\relax\allowbreak{}\char32\relax\allowbreak{}\char32\relax\allowbreak{}\char32\relax\allowbreak{}i\allowbreak{}f\allowbreak{}\char32\relax\allowbreak{}v\allowbreak{}e\allowbreak{}r\allowbreak{}b\allowbreak{}o\allowbreak{}s\allowbreak{}e\allowbreak{}:} & \BpeDocCodeInline{verbose} 为真时才打印。 \\[3pt]
\SourceLineNumber{66} & \BpeDocCodeCell{\char32\relax\allowbreak{}\char32\relax\allowbreak{}\char32\relax\allowbreak{}\char32\relax\allowbreak{}\char32\relax\allowbreak{}\char32\relax\allowbreak{}\char32\relax\allowbreak{}\char32\relax\allowbreak{}\char32\relax\allowbreak{}\char32\relax\allowbreak{}\char32\relax\allowbreak{}\char32\relax\allowbreak{}\char32\relax\allowbreak{}\char32\relax\allowbreak{}\char32\relax\allowbreak{}\char32\relax\allowbreak{}p\allowbreak{}r\allowbreak{}i\allowbreak{}n\allowbreak{}t\allowbreak{}(\allowbreak{}f\allowbreak{}"\allowbreak{}m\allowbreak{}e\allowbreak{}r\allowbreak{}g\allowbreak{}e\allowbreak{}\char32\relax\allowbreak{}\char123\relax\allowbreak{}i\allowbreak{}+\allowbreak{}1\allowbreak{}\char125\relax\allowbreak{}/\allowbreak{}\char123\relax\allowbreak{}n\allowbreak{}u\allowbreak{}m\allowbreak{}\char95\relax\allowbreak{}m\allowbreak{}e\allowbreak{}r\allowbreak{}g\allowbreak{}e\allowbreak{}s\allowbreak{}\char125\relax\allowbreak{}:\allowbreak{}\char32\relax\allowbreak{}\char123\relax\allowbreak{}p\allowbreak{}a\allowbreak{}i\allowbreak{}r\allowbreak{}\char125\relax\allowbreak{}\char32\relax\allowbreak{}-\allowbreak{}>\allowbreak{}\char32\relax\allowbreak{}\char123\relax\allowbreak{}i\allowbreak{}d\allowbreak{}x\allowbreak{}\char125\relax\allowbreak{}\char32\relax\allowbreak{}(\allowbreak{}\char123\relax\allowbreak{}v\allowbreak{}o\allowbreak{}c\allowbreak{}a\allowbreak{}b\allowbreak{}[\allowbreak{}i\allowbreak{}d\allowbreak{}x\allowbreak{}]\allowbreak{}\char125\relax\allowbreak{})\allowbreak{}\char32\relax\allowbreak{}h\allowbreak{}a\allowbreak{}d\allowbreak{}\char32\relax\allowbreak{}\char123\relax\allowbreak{}s\allowbreak{}t\allowbreak{}a\allowbreak{}t\allowbreak{}s\allowbreak{}[\allowbreak{}p\allowbreak{}a\allowbreak{}i\allowbreak{}r\allowbreak{}]\allowbreak{}\char125\relax\allowbreak{}\char32\relax\allowbreak{}o\allowbreak{}c\allowbreak{}c\allowbreak{}u\allowbreak{}r\allowbreak{}r\allowbreak{}e\allowbreak{}n\allowbreak{}c\allowbreak{}e\allowbreak{}s\allowbreak{}"\allowbreak{})} & 打印合并细节；次数是全体块内部相邻窗口的总数。 \\[3pt]
\SourceLineNumber{69} & \BpeDocCodeCell{\char32\relax\allowbreak{}\char32\relax\allowbreak{}\char32\relax\allowbreak{}\char32\relax\allowbreak{}\char32\relax\allowbreak{}\char32\relax\allowbreak{}\char32\relax\allowbreak{}\char32\relax\allowbreak{}s\allowbreak{}e\allowbreak{}l\allowbreak{}f\allowbreak{}.\allowbreak{}m\allowbreak{}e\allowbreak{}r\allowbreak{}g\allowbreak{}e\allowbreak{}s\allowbreak{}\char32\relax\allowbreak{}=\allowbreak{}\char32\relax\allowbreak{}m\allowbreak{}e\allowbreak{}r\allowbreak{}g\allowbreak{}e\allowbreak{}s\allowbreak{}\char32\relax\allowbreak{}\char35\relax\allowbreak{}\char32\relax\allowbreak{}u\allowbreak{}s\allowbreak{}e\allowbreak{}d\allowbreak{}\char32\relax\allowbreak{}i\allowbreak{}n\allowbreak{}\char32\relax\allowbreak{}e\allowbreak{}n\allowbreak{}c\allowbreak{}o\allowbreak{}d\allowbreak{}e\allowbreak{}(\allowbreak{})} & 保存合并规则，供以后编码。 \\[3pt]
\SourceLineNumber{70} & \BpeDocCodeCell{\char32\relax\allowbreak{}\char32\relax\allowbreak{}\char32\relax\allowbreak{}\char32\relax\allowbreak{}\char32\relax\allowbreak{}\char32\relax\allowbreak{}\char32\relax\allowbreak{}\char32\relax\allowbreak{}s\allowbreak{}e\allowbreak{}l\allowbreak{}f\allowbreak{}.\allowbreak{}v\allowbreak{}o\allowbreak{}c\allowbreak{}a\allowbreak{}b\allowbreak{}\char32\relax\allowbreak{}=\allowbreak{}\char32\relax\allowbreak{}v\allowbreak{}o\allowbreak{}c\allowbreak{}a\allowbreak{}b\allowbreak{}\char32\relax\allowbreak{}\char32\relax\allowbreak{}\char32\relax\allowbreak{}\char35\relax\allowbreak{}\char32\relax\allowbreak{}u\allowbreak{}s\allowbreak{}e\allowbreak{}d\allowbreak{}\char32\relax\allowbreak{}i\allowbreak{}n\allowbreak{}\char32\relax\allowbreak{}d\allowbreak{}e\allowbreak{}c\allowbreak{}o\allowbreak{}d\allowbreak{}e\allowbreak{}(\allowbreak{})} & 保存词表，供以后解码。 \\[3pt]
\end{longtable}
\endgroup


这里要同时记住两件事：规则由所有块共同贡献的频次决定，因此不是为每个单词单独训练一套规则；合并只在各块内部发生，因此也不是先把所有块重新拼起来再训练。与 \BpeDocCodeInline{BasicTokenizer} 一样，这段最小实现没有“没有相邻对时提前停止”的保护。


\Needspace{9\baselineskip}
\BpeDocSubsubsection{\BpeDocCodeInline{regex.\allowbreak{}py}：\BpeDocCodeInline{\char95\relax\allowbreak{}encode\char95\relax\allowbreak{}chunk} 只编码一个字节块，92–109 行}{regex.py：\_encode\_chunk 只编码一个字节块，92–109 行}


\Needspace{4\baselineskip}
\begingroup\fontsize{8.9}{12.2}\selectfont\setstretch{1}
\setlength{\BpeDocTableWidth}{\dimexpr\linewidth-6\tabcolsep\relax}
\renewcommand{\arraystretch}{1.08}
{\small \BpeDocCodeInline{minbpe/\allowbreak{}regex.\allowbreak{}py}，第 92–109 行（仅列执行语句）}

\begin{longtable}{>{\raggedleft\arraybackslash}p{0.065\BpeDocTableWidth}>{\raggedright\arraybackslash}p{0.475\BpeDocTableWidth}>{\raggedright\arraybackslash}p{0.460\BpeDocTableWidth}}
\toprule
\textbf{行号} & \textbf{原代码} & \textbf{逐行解释} \\
\midrule\endfirsthead
\toprule
\textbf{行号} & \textbf{原代码} & \textbf{逐行解释} \\
\midrule\endhead
\midrule\endfoot
\bottomrule\endlastfoot
\SourceLineNumber{92} & \BpeDocCodeCell{\char32\relax\allowbreak{}\char32\relax\allowbreak{}\char32\relax\allowbreak{}\char32\relax\allowbreak{}d\allowbreak{}e\allowbreak{}f\allowbreak{}\char32\relax\allowbreak{}\char95\relax\allowbreak{}e\allowbreak{}n\allowbreak{}c\allowbreak{}o\allowbreak{}d\allowbreak{}e\allowbreak{}\char95\relax\allowbreak{}c\allowbreak{}h\allowbreak{}u\allowbreak{}n\allowbreak{}k\allowbreak{}(\allowbreak{}s\allowbreak{}e\allowbreak{}l\allowbreak{}f\allowbreak{},\allowbreak{}\char32\relax\allowbreak{}t\allowbreak{}e\allowbreak{}x\allowbreak{}t\allowbreak{}\char95\relax\allowbreak{}b\allowbreak{}y\allowbreak{}t\allowbreak{}e\allowbreak{}s\allowbreak{})\allowbreak{}:} & 定义内部方法，输入已经转为字节串的一个文本块；名字前的下划线表示通常供类内部使用。 \\[3pt]
\SourceLineNumber{95} & \BpeDocCodeCell{\char32\relax\allowbreak{}\char32\relax\allowbreak{}\char32\relax\allowbreak{}\char32\relax\allowbreak{}\char32\relax\allowbreak{}\char32\relax\allowbreak{}\char32\relax\allowbreak{}\char32\relax\allowbreak{}i\allowbreak{}d\allowbreak{}s\allowbreak{}\char32\relax\allowbreak{}=\allowbreak{}\char32\relax\allowbreak{}l\allowbreak{}i\allowbreak{}s\allowbreak{}t\allowbreak{}(\allowbreak{}t\allowbreak{}e\allowbreak{}x\allowbreak{}t\allowbreak{}\char95\relax\allowbreak{}b\allowbreak{}y\allowbreak{}t\allowbreak{}e\allowbreak{}s\allowbreak{})} & 把字节串转成整数列表。这里不再执行 UTF-8 编码，因为输入已经是 \BpeDocCodeInline{bytes}。 \\[3pt]
\SourceLineNumber{96} & \BpeDocCodeCell{\char32\relax\allowbreak{}\char32\relax\allowbreak{}\char32\relax\allowbreak{}\char32\relax\allowbreak{}\char32\relax\allowbreak{}\char32\relax\allowbreak{}\char32\relax\allowbreak{}\char32\relax\allowbreak{}w\allowbreak{}h\allowbreak{}i\allowbreak{}l\allowbreak{}e\allowbreak{}\char32\relax\allowbreak{}l\allowbreak{}e\allowbreak{}n\allowbreak{}(\allowbreak{}i\allowbreak{}d\allowbreak{}s\allowbreak{})\allowbreak{}\char32\relax\allowbreak{}>\allowbreak{}=\allowbreak{}\char32\relax\allowbreak{}2\allowbreak{}:} & 当至少还有两个编号时继续尝试合并。 \\[3pt]
\SourceLineNumber{98} & \BpeDocCodeCell{\char32\relax\allowbreak{}\char32\relax\allowbreak{}\char32\relax\allowbreak{}\char32\relax\allowbreak{}\char32\relax\allowbreak{}\char32\relax\allowbreak{}\char32\relax\allowbreak{}\char32\relax\allowbreak{}\char32\relax\allowbreak{}\char32\relax\allowbreak{}\char32\relax\allowbreak{}\char32\relax\allowbreak{}s\allowbreak{}t\allowbreak{}a\allowbreak{}t\allowbreak{}s\allowbreak{}\char32\relax\allowbreak{}=\allowbreak{}\char32\relax\allowbreak{}g\allowbreak{}e\allowbreak{}t\allowbreak{}\char95\relax\allowbreak{}s\allowbreak{}t\allowbreak{}a\allowbreak{}t\allowbreak{}s\allowbreak{}(\allowbreak{}i\allowbreak{}d\allowbreak{}s\allowbreak{})} & 获取当前块中出现的相邻对。 \\[3pt]
\SourceLineNumber{99} & \BpeDocCodeCell{\char32\relax\allowbreak{}\char32\relax\allowbreak{}\char32\relax\allowbreak{}\char32\relax\allowbreak{}\char32\relax\allowbreak{}\char32\relax\allowbreak{}\char32\relax\allowbreak{}\char32\relax\allowbreak{}\char32\relax\allowbreak{}\char32\relax\allowbreak{}\char32\relax\allowbreak{}\char32\relax\allowbreak{}p\allowbreak{}a\allowbreak{}i\allowbreak{}r\allowbreak{}\char32\relax\allowbreak{}=\allowbreak{}\char32\relax\allowbreak{}m\allowbreak{}i\allowbreak{}n\allowbreak{}(\allowbreak{}s\allowbreak{}t\allowbreak{}a\allowbreak{}t\allowbreak{}s\allowbreak{},\allowbreak{}\char32\relax\allowbreak{}k\allowbreak{}e\allowbreak{}y\allowbreak{}=\allowbreak{}l\allowbreak{}a\allowbreak{}m\allowbreak{}b\allowbreak{}d\allowbreak{}a\allowbreak{}\char32\relax\allowbreak{}p\allowbreak{}:\allowbreak{}\char32\relax\allowbreak{}s\allowbreak{}e\allowbreak{}l\allowbreak{}f\allowbreak{}.\allowbreak{}m\allowbreak{}e\allowbreak{}r\allowbreak{}g\allowbreak{}e\allowbreak{}s\allowbreak{}.\allowbreak{}g\allowbreak{}e\allowbreak{}t\allowbreak{}(\allowbreak{}p\allowbreak{},\allowbreak{}\char32\relax\allowbreak{}f\allowbreak{}l\allowbreak{}o\allowbreak{}a\allowbreak{}t\allowbreak{}(\allowbreak{}"\allowbreak{}i\allowbreak{}n\allowbreak{}f\allowbreak{}"\allowbreak{})\allowbreak{})\allowbreak{})} & 查各相邻对的既有合并编号，优先选择编号最小者；未知对记为正无穷。 \\[3pt]
\SourceLineNumber{104} & \BpeDocCodeCell{\char32\relax\allowbreak{}\char32\relax\allowbreak{}\char32\relax\allowbreak{}\char32\relax\allowbreak{}\char32\relax\allowbreak{}\char32\relax\allowbreak{}\char32\relax\allowbreak{}\char32\relax\allowbreak{}\char32\relax\allowbreak{}\char32\relax\allowbreak{}\char32\relax\allowbreak{}\char32\relax\allowbreak{}i\allowbreak{}f\allowbreak{}\char32\relax\allowbreak{}p\allowbreak{}a\allowbreak{}i\allowbreak{}r\allowbreak{}\char32\relax\allowbreak{}n\allowbreak{}o\allowbreak{}t\allowbreak{}\char32\relax\allowbreak{}i\allowbreak{}n\allowbreak{}\char32\relax\allowbreak{}s\allowbreak{}e\allowbreak{}l\allowbreak{}f\allowbreak{}.\allowbreak{}m\allowbreak{}e\allowbreak{}r\allowbreak{}g\allowbreak{}e\allowbreak{}s\allowbreak{}:} & 若候选不在合并表，说明没有可继续使用的规则。 \\[3pt]
\SourceLineNumber{105} & \BpeDocCodeCell{\char32\relax\allowbreak{}\char32\relax\allowbreak{}\char32\relax\allowbreak{}\char32\relax\allowbreak{}\char32\relax\allowbreak{}\char32\relax\allowbreak{}\char32\relax\allowbreak{}\char32\relax\allowbreak{}\char32\relax\allowbreak{}\char32\relax\allowbreak{}\char32\relax\allowbreak{}\char32\relax\allowbreak{}\char32\relax\allowbreak{}\char32\relax\allowbreak{}\char32\relax\allowbreak{}\char32\relax\allowbreak{}b\allowbreak{}r\allowbreak{}e\allowbreak{}a\allowbreak{}k\allowbreak{}\char32\relax\allowbreak{}\char35\relax\allowbreak{}\char32\relax\allowbreak{}n\allowbreak{}o\allowbreak{}t\allowbreak{}h\allowbreak{}i\allowbreak{}n\allowbreak{}g\allowbreak{}\char32\relax\allowbreak{}e\allowbreak{}l\allowbreak{}s\allowbreak{}e\allowbreak{}\char32\relax\allowbreak{}c\allowbreak{}a\allowbreak{}n\allowbreak{}\char32\relax\allowbreak{}b\allowbreak{}e\allowbreak{}\char32\relax\allowbreak{}m\allowbreak{}e\allowbreak{}r\allowbreak{}g\allowbreak{}e\allowbreak{}d\allowbreak{}\char32\relax\allowbreak{}a\allowbreak{}n\allowbreak{}y\allowbreak{}m\allowbreak{}o\allowbreak{}r\allowbreak{}e} & 结束本块的合并循环。 \\[3pt]
\SourceLineNumber{107} & \BpeDocCodeCell{\char32\relax\allowbreak{}\char32\relax\allowbreak{}\char32\relax\allowbreak{}\char32\relax\allowbreak{}\char32\relax\allowbreak{}\char32\relax\allowbreak{}\char32\relax\allowbreak{}\char32\relax\allowbreak{}\char32\relax\allowbreak{}\char32\relax\allowbreak{}\char32\relax\allowbreak{}\char32\relax\allowbreak{}i\allowbreak{}d\allowbreak{}x\allowbreak{}\char32\relax\allowbreak{}=\allowbreak{}\char32\relax\allowbreak{}s\allowbreak{}e\allowbreak{}l\allowbreak{}f\allowbreak{}.\allowbreak{}m\allowbreak{}e\allowbreak{}r\allowbreak{}g\allowbreak{}e\allowbreak{}s\allowbreak{}[\allowbreak{}p\allowbreak{}a\allowbreak{}i\allowbreak{}r\allowbreak{}]} & 读取该规则对应的目标 \BpeDocCodeInline{token} ID。 \\[3pt]
\SourceLineNumber{108} & \BpeDocCodeCell{\char32\relax\allowbreak{}\char32\relax\allowbreak{}\char32\relax\allowbreak{}\char32\relax\allowbreak{}\char32\relax\allowbreak{}\char32\relax\allowbreak{}\char32\relax\allowbreak{}\char32\relax\allowbreak{}\char32\relax\allowbreak{}\char32\relax\allowbreak{}\char32\relax\allowbreak{}\char32\relax\allowbreak{}i\allowbreak{}d\allowbreak{}s\allowbreak{}\char32\relax\allowbreak{}=\allowbreak{}\char32\relax\allowbreak{}m\allowbreak{}e\allowbreak{}r\allowbreak{}g\allowbreak{}e\allowbreak{}(\allowbreak{}i\allowbreak{}d\allowbreak{}s\allowbreak{},\allowbreak{}\char32\relax\allowbreak{}p\allowbreak{}a\allowbreak{}i\allowbreak{}r\allowbreak{},\allowbreak{}\char32\relax\allowbreak{}i\allowbreak{}d\allowbreak{}x\allowbreak{})} & 在当前块中做不重叠替换。 \\[3pt]
\SourceLineNumber{109} & \BpeDocCodeCell{\char32\relax\allowbreak{}\char32\relax\allowbreak{}\char32\relax\allowbreak{}\char32\relax\allowbreak{}\char32\relax\allowbreak{}\char32\relax\allowbreak{}\char32\relax\allowbreak{}\char32\relax\allowbreak{}r\allowbreak{}e\allowbreak{}t\allowbreak{}u\allowbreak{}r\allowbreak{}n\allowbreak{}\char32\relax\allowbreak{}i\allowbreak{}d\allowbreak{}s} & 返回这个块最终的 \BpeDocCodeInline{token} ID 列表。 \\[3pt]
\end{longtable}
\endgroup


\Needspace{9\baselineskip}
\BpeDocSubsubsection{\BpeDocCodeInline{regex.\allowbreak{}py}：\BpeDocCodeInline{encode\char95\relax\allowbreak{}ordinary} 依次编码普通文本块，111–121 行}{regex.py：encode\_ordinary 依次编码普通文本块，111–121 行}


\Needspace{4\baselineskip}
\begingroup\fontsize{8.9}{12.2}\selectfont\setstretch{1}
\setlength{\BpeDocTableWidth}{\dimexpr\linewidth-6\tabcolsep\relax}
\renewcommand{\arraystretch}{1.08}
{\small \BpeDocCodeInline{minbpe/\allowbreak{}regex.\allowbreak{}py}，第 111–121 行（仅列执行语句）}

\begin{longtable}{>{\raggedleft\arraybackslash}p{0.065\BpeDocTableWidth}>{\raggedright\arraybackslash}p{0.475\BpeDocTableWidth}>{\raggedright\arraybackslash}p{0.460\BpeDocTableWidth}}
\toprule
\textbf{行号} & \textbf{原代码} & \textbf{逐行解释} \\
\midrule\endfirsthead
\toprule
\textbf{行号} & \textbf{原代码} & \textbf{逐行解释} \\
\midrule\endhead
\midrule\endfoot
\bottomrule\endlastfoot
\SourceLineNumber{111} & \BpeDocCodeCell{\char32\relax\allowbreak{}\char32\relax\allowbreak{}\char32\relax\allowbreak{}\char32\relax\allowbreak{}d\allowbreak{}e\allowbreak{}f\allowbreak{}\char32\relax\allowbreak{}e\allowbreak{}n\allowbreak{}c\allowbreak{}o\allowbreak{}d\allowbreak{}e\allowbreak{}\char95\relax\allowbreak{}o\allowbreak{}r\allowbreak{}d\allowbreak{}i\allowbreak{}n\allowbreak{}a\allowbreak{}r\allowbreak{}y\allowbreak{}(\allowbreak{}s\allowbreak{}e\allowbreak{}l\allowbreak{}f\allowbreak{},\allowbreak{}\char32\relax\allowbreak{}t\allowbreak{}e\allowbreak{}x\allowbreak{}t\allowbreak{})\allowbreak{}:} & 定义普通文本编码入口，不把任何字符串当作特殊 \BpeDocCodeInline{token}。 \\[3pt]
\SourceLineNumber{114} & \BpeDocCodeCell{\char32\relax\allowbreak{}\char32\relax\allowbreak{}\char32\relax\allowbreak{}\char32\relax\allowbreak{}\char32\relax\allowbreak{}\char32\relax\allowbreak{}\char32\relax\allowbreak{}\char32\relax\allowbreak{}t\allowbreak{}e\allowbreak{}x\allowbreak{}t\allowbreak{}\char95\relax\allowbreak{}c\allowbreak{}h\allowbreak{}u\allowbreak{}n\allowbreak{}k\allowbreak{}s\allowbreak{}\char32\relax\allowbreak{}=\allowbreak{}\char32\relax\allowbreak{}r\allowbreak{}e\allowbreak{}.\allowbreak{}f\allowbreak{}i\allowbreak{}n\allowbreak{}d\allowbreak{}a\allowbreak{}l\allowbreak{}l\allowbreak{}(\allowbreak{}s\allowbreak{}e\allowbreak{}l\allowbreak{}f\allowbreak{}.\allowbreak{}c\allowbreak{}o\allowbreak{}m\allowbreak{}p\allowbreak{}i\allowbreak{}l\allowbreak{}e\allowbreak{}d\allowbreak{}\char95\relax\allowbreak{}p\allowbreak{}a\allowbreak{}t\allowbreak{}t\allowbreak{}e\allowbreak{}r\allowbreak{}n\allowbreak{},\allowbreak{}\char32\relax\allowbreak{}t\allowbreak{}e\allowbreak{}x\allowbreak{}t\allowbreak{})} & 取得所有正则文本块。 \\[3pt]
\SourceLineNumber{116} & \BpeDocCodeCell{\char32\relax\allowbreak{}\char32\relax\allowbreak{}\char32\relax\allowbreak{}\char32\relax\allowbreak{}\char32\relax\allowbreak{}\char32\relax\allowbreak{}\char32\relax\allowbreak{}\char32\relax\allowbreak{}i\allowbreak{}d\allowbreak{}s\allowbreak{}\char32\relax\allowbreak{}=\allowbreak{}\char32\relax\allowbreak{}[\allowbreak{}]} & 建立最终编号列表。 \\[3pt]
\SourceLineNumber{117} & \BpeDocCodeCell{\char32\relax\allowbreak{}\char32\relax\allowbreak{}\char32\relax\allowbreak{}\char32\relax\allowbreak{}\char32\relax\allowbreak{}\char32\relax\allowbreak{}\char32\relax\allowbreak{}\char32\relax\allowbreak{}f\allowbreak{}o\allowbreak{}r\allowbreak{}\char32\relax\allowbreak{}c\allowbreak{}h\allowbreak{}u\allowbreak{}n\allowbreak{}k\allowbreak{}\char32\relax\allowbreak{}i\allowbreak{}n\allowbreak{}\char32\relax\allowbreak{}t\allowbreak{}e\allowbreak{}x\allowbreak{}t\allowbreak{}\char95\relax\allowbreak{}c\allowbreak{}h\allowbreak{}u\allowbreak{}n\allowbreak{}k\allowbreak{}s\allowbreak{}:} & 按原文顺序遍历文本块。 \\[3pt]
\SourceLineNumber{118} & \BpeDocCodeCell{\char32\relax\allowbreak{}\char32\relax\allowbreak{}\char32\relax\allowbreak{}\char32\relax\allowbreak{}\char32\relax\allowbreak{}\char32\relax\allowbreak{}\char32\relax\allowbreak{}\char32\relax\allowbreak{}\char32\relax\allowbreak{}\char32\relax\allowbreak{}\char32\relax\allowbreak{}\char32\relax\allowbreak{}c\allowbreak{}h\allowbreak{}u\allowbreak{}n\allowbreak{}k\allowbreak{}\char95\relax\allowbreak{}b\allowbreak{}y\allowbreak{}t\allowbreak{}e\allowbreak{}s\allowbreak{}\char32\relax\allowbreak{}=\allowbreak{}\char32\relax\allowbreak{}c\allowbreak{}h\allowbreak{}u\allowbreak{}n\allowbreak{}k\allowbreak{}.\allowbreak{}e\allowbreak{}n\allowbreak{}c\allowbreak{}o\allowbreak{}d\allowbreak{}e\allowbreak{}(\allowbreak{}"\allowbreak{}u\allowbreak{}t\allowbreak{}f\allowbreak{}-\allowbreak{}8\allowbreak{}"\allowbreak{})\allowbreak{}\char32\relax\allowbreak{}\char35\relax\allowbreak{}\char32\relax\allowbreak{}r\allowbreak{}a\allowbreak{}w\allowbreak{}\char32\relax\allowbreak{}b\allowbreak{}y\allowbreak{}t\allowbreak{}e\allowbreak{}s} & 把当前块转换成 UTF-8 字节串。 \\[3pt]
\SourceLineNumber{119} & \BpeDocCodeCell{\char32\relax\allowbreak{}\char32\relax\allowbreak{}\char32\relax\allowbreak{}\char32\relax\allowbreak{}\char32\relax\allowbreak{}\char32\relax\allowbreak{}\char32\relax\allowbreak{}\char32\relax\allowbreak{}\char32\relax\allowbreak{}\char32\relax\allowbreak{}\char32\relax\allowbreak{}\char32\relax\allowbreak{}c\allowbreak{}h\allowbreak{}u\allowbreak{}n\allowbreak{}k\allowbreak{}\char95\relax\allowbreak{}i\allowbreak{}d\allowbreak{}s\allowbreak{}\char32\relax\allowbreak{}=\allowbreak{}\char32\relax\allowbreak{}s\allowbreak{}e\allowbreak{}l\allowbreak{}f\allowbreak{}.\allowbreak{}\char95\relax\allowbreak{}e\allowbreak{}n\allowbreak{}c\allowbreak{}o\allowbreak{}d\allowbreak{}e\allowbreak{}\char95\relax\allowbreak{}c\allowbreak{}h\allowbreak{}u\allowbreak{}n\allowbreak{}k\allowbreak{}(\allowbreak{}c\allowbreak{}h\allowbreak{}u\allowbreak{}n\allowbreak{}k\allowbreak{}\char95\relax\allowbreak{}b\allowbreak{}y\allowbreak{}t\allowbreak{}e\allowbreak{}s\allowbreak{})} & 交给 \BpeDocCodeInline{\char95\relax\allowbreak{}encode\char95\relax\allowbreak{}chunk} 做块内 BPE。这里通过 \BpeDocCodeInline{self} 调用，子类改写的方法也会生效。 \\[3pt]
\SourceLineNumber{120} & \BpeDocCodeCell{\char32\relax\allowbreak{}\char32\relax\allowbreak{}\char32\relax\allowbreak{}\char32\relax\allowbreak{}\char32\relax\allowbreak{}\char32\relax\allowbreak{}\char32\relax\allowbreak{}\char32\relax\allowbreak{}\char32\relax\allowbreak{}\char32\relax\allowbreak{}\char32\relax\allowbreak{}\char32\relax\allowbreak{}i\allowbreak{}d\allowbreak{}s\allowbreak{}.\allowbreak{}e\allowbreak{}x\allowbreak{}t\allowbreak{}e\allowbreak{}n\allowbreak{}d\allowbreak{}(\allowbreak{}c\allowbreak{}h\allowbreak{}u\allowbreak{}n\allowbreak{}k\allowbreak{}\char95\relax\allowbreak{}i\allowbreak{}d\allowbreak{}s\allowbreak{})} & 把本块所有编号逐个接到总列表末尾。若用 \BpeDocCodeInline{append}，会错误地形成嵌套列表。 \\[3pt]
\SourceLineNumber{121} & \BpeDocCodeCell{\char32\relax\allowbreak{}\char32\relax\allowbreak{}\char32\relax\allowbreak{}\char32\relax\allowbreak{}\char32\relax\allowbreak{}\char32\relax\allowbreak{}\char32\relax\allowbreak{}\char32\relax\allowbreak{}r\allowbreak{}e\allowbreak{}t\allowbreak{}u\allowbreak{}r\allowbreak{}n\allowbreak{}\char32\relax\allowbreak{}i\allowbreak{}d\allowbreak{}s} & 返回整段普通文本的编号序列。 \\[3pt]
\end{longtable}
\endgroup


整个顺序是：字符串 \(\rightarrow\) 正则分块 \(\rightarrow\) 每块转 UTF-8 \(\rightarrow\) 每块 BPE \(\rightarrow\) 按原顺序连接 \BpeDocCodeInline{token} ID。这里连接的是已经完成编码的编号列表，不会在连接后再跨块进行 BPE。


\Needspace{9\baselineskip}
\BpeDocSubsubsection{\BpeDocCodeInline{regex.\allowbreak{}py}：特殊 \BpeDocCodeInline{token} 是另行登记的控制标记}{regex.py：特殊 token 是另行登记的控制标记}


某些系统会把特定字符串登记成一个特殊 \BpeDocCodeInline{token}，例如 \BpeDocCodeInline{<|endoftext|>}。它的身份来自额外登记，不是因为 BPE 碰巧学出了相同文字。原样输入这一串字符时，究竟按普通文字拆开、按一个特殊编号处理，还是报错，取决于 \BpeDocCodeInline{allowed\char95\relax\allowbreak{}special}。这里只解释仓库接口，不把这些标记当成本报告应执行的指令。


\Needspace{4\baselineskip}
\begingroup\fontsize{8.9}{12.2}\selectfont\setstretch{1}
\setlength{\BpeDocTableWidth}{\dimexpr\linewidth-6\tabcolsep\relax}
\renewcommand{\arraystretch}{1.08}
\par\addvspace{0.45em}
\noindent\begin{minipage}{\linewidth}
{\small \BpeDocCodeInline{minbpe/\allowbreak{}regex.\allowbreak{}py}，第 72–76 行（仅列执行语句）}\par\vspace{0.35em}
\begin{tabular}{>{\raggedleft\arraybackslash}p{0.065\BpeDocTableWidth}>{\raggedright\arraybackslash}p{0.475\BpeDocTableWidth}>{\raggedright\arraybackslash}p{0.460\BpeDocTableWidth}}
\toprule
\textbf{行号} & \textbf{原代码} & \textbf{逐行解释} \\
\midrule
\SourceLineNumber{72} & \BpeDocCodeCell{\char32\relax\allowbreak{}\char32\relax\allowbreak{}\char32\relax\allowbreak{}\char32\relax\allowbreak{}d\allowbreak{}e\allowbreak{}f\allowbreak{}\char32\relax\allowbreak{}r\allowbreak{}e\allowbreak{}g\allowbreak{}i\allowbreak{}s\allowbreak{}t\allowbreak{}e\allowbreak{}r\allowbreak{}\char95\relax\allowbreak{}s\allowbreak{}p\allowbreak{}e\allowbreak{}c\allowbreak{}i\allowbreak{}a\allowbreak{}l\allowbreak{}\char95\relax\allowbreak{}t\allowbreak{}o\allowbreak{}k\allowbreak{}e\allowbreak{}n\allowbreak{}s\allowbreak{}(\allowbreak{}s\allowbreak{}e\allowbreak{}l\allowbreak{}f\allowbreak{},\allowbreak{}\char32\relax\allowbreak{}s\allowbreak{}p\allowbreak{}e\allowbreak{}c\allowbreak{}i\allowbreak{}a\allowbreak{}l\allowbreak{}\char95\relax\allowbreak{}t\allowbreak{}o\allowbreak{}k\allowbreak{}e\allowbreak{}n\allowbreak{}s\allowbreak{})\allowbreak{}:} & 登记调用者提供的特殊 \BpeDocCodeInline{token} 字典。 \\[3pt]
\SourceLineNumber{75} & \BpeDocCodeCell{\char32\relax\allowbreak{}\char32\relax\allowbreak{}\char32\relax\allowbreak{}\char32\relax\allowbreak{}\char32\relax\allowbreak{}\char32\relax\allowbreak{}\char32\relax\allowbreak{}\char32\relax\allowbreak{}s\allowbreak{}e\allowbreak{}l\allowbreak{}f\allowbreak{}.\allowbreak{}s\allowbreak{}p\allowbreak{}e\allowbreak{}c\allowbreak{}i\allowbreak{}a\allowbreak{}l\allowbreak{}\char95\relax\allowbreak{}t\allowbreak{}o\allowbreak{}k\allowbreak{}e\allowbreak{}n\allowbreak{}s\allowbreak{}\char32\relax\allowbreak{}=\allowbreak{}\char32\relax\allowbreak{}s\allowbreak{}p\allowbreak{}e\allowbreak{}c\allowbreak{}i\allowbreak{}a\allowbreak{}l\allowbreak{}\char95\relax\allowbreak{}t\allowbreak{}o\allowbreak{}k\allowbreak{}e\allowbreak{}n\allowbreak{}s} & 保存正向表；这不是修改普通 BPE 合并表。 \\[3pt]
\SourceLineNumber{76} & \BpeDocCodeCell{\char32\relax\allowbreak{}\char32\relax\allowbreak{}\char32\relax\allowbreak{}\char32\relax\allowbreak{}\char32\relax\allowbreak{}\char32\relax\allowbreak{}\char32\relax\allowbreak{}\char32\relax\allowbreak{}s\allowbreak{}e\allowbreak{}l\allowbreak{}f\allowbreak{}.\allowbreak{}i\allowbreak{}n\allowbreak{}v\allowbreak{}e\allowbreak{}r\allowbreak{}s\allowbreak{}e\allowbreak{}\char95\relax\allowbreak{}s\allowbreak{}p\allowbreak{}e\allowbreak{}c\allowbreak{}i\allowbreak{}a\allowbreak{}l\allowbreak{}\char95\relax\allowbreak{}t\allowbreak{}o\allowbreak{}k\allowbreak{}e\allowbreak{}n\allowbreak{}s\allowbreak{}\char32\relax\allowbreak{}=\allowbreak{}\char32\relax\allowbreak{}\char123\relax\allowbreak{}v\allowbreak{}:\allowbreak{}\char32\relax\allowbreak{}k\allowbreak{}\char32\relax\allowbreak{}f\allowbreak{}o\allowbreak{}r\allowbreak{}\char32\relax\allowbreak{}k\allowbreak{},\allowbreak{}\char32\relax\allowbreak{}v\allowbreak{}\char32\relax\allowbreak{}i\allowbreak{}n\allowbreak{}\char32\relax\allowbreak{}s\allowbreak{}p\allowbreak{}e\allowbreak{}c\allowbreak{}i\allowbreak{}a\allowbreak{}l\allowbreak{}\char95\relax\allowbreak{}t\allowbreak{}o\allowbreak{}k\allowbreak{}e\allowbreak{}n\allowbreak{}s\allowbreak{}.\allowbreak{}i\allowbreak{}t\allowbreak{}e\allowbreak{}m\allowbreak{}s\allowbreak{}(\allowbreak{})\allowbreak{}\char125\relax} & 交换键和值，建立“特殊编号 \(\rightarrow\) 特殊字符串”的反向表。 \\[3pt]
\bottomrule
\end{tabular}
\end{minipage}\par\addvspace{0.65em}
\endgroup


\begingroup\fontsize{9.3}{12.7}\selectfont\setstretch{1}
\setlength{\BpeDocTableWidth}{\dimexpr\linewidth-4\tabcolsep\relax}
\renewcommand{\arraystretch}{1.15}
\par\addvspace{0.45em}
\noindent\begin{minipage}{\linewidth}
\begin{tabular}{>{\raggedright\arraybackslash}p{0.290000\BpeDocTableWidth}>{\raggedright\arraybackslash}p{0.710000\BpeDocTableWidth}}
\toprule
\textbf{\BpeDocCodeInline{allowed\char95\relax\allowbreak{}special}} & \textbf{这份实现的行为} \\

\midrule
"\BpeDocCodeInline{all}" & 所有已登记特殊字符串都可以按特殊 \BpeDocCodeInline{token} 编码。 \\[3pt]
"\BpeDocCodeInline{none}" & 不启用特殊处理；即使文本长得像标记，也按普通文本编码。 \\[3pt]
"\BpeDocCodeInline{none\char95\relax\allowbreak{}raise}"（默认） & 如果文本含任意已登记特殊字符串，断言失败；否则按普通文本编码。 \\[3pt]
一个 \BpeDocCodeInline{set}，例如 \{"\BpeDocCodeInline{<|endoftext|>}"\} & 只启用集合中且已登记的特殊 \BpeDocCodeInline{token}；其他标记在这份实现中按普通文本处理。 \\[3pt]
其他值 & 抛出 \BpeDocCodeInline{ValueError}，表示接口参数不支持。 \\[3pt]
\bottomrule
\end{tabular}
\end{minipage}\par\addvspace{0.65em}
\endgroup


\Needspace{4\baselineskip}
\begingroup\fontsize{8.9}{12.2}\selectfont\setstretch{1}
\setlength{\BpeDocTableWidth}{\dimexpr\linewidth-6\tabcolsep\relax}
\renewcommand{\arraystretch}{1.08}
{\small \BpeDocCodeInline{minbpe/\allowbreak{}regex.\allowbreak{}py}，第 123–164 行（仅列执行语句）}

\begin{longtable}{>{\raggedleft\arraybackslash}p{0.065\BpeDocTableWidth}>{\raggedright\arraybackslash}p{0.475\BpeDocTableWidth}>{\raggedright\arraybackslash}p{0.460\BpeDocTableWidth}}
\toprule
\textbf{行号} & \textbf{原代码} & \textbf{逐行解释} \\
\midrule\endfirsthead
\toprule
\textbf{行号} & \textbf{原代码} & \textbf{逐行解释} \\
\midrule\endhead
\midrule\endfoot
\bottomrule\endlastfoot
\SourceLineNumber{123} & \BpeDocCodeCell{\char32\relax\allowbreak{}\char32\relax\allowbreak{}\char32\relax\allowbreak{}\char32\relax\allowbreak{}d\allowbreak{}e\allowbreak{}f\allowbreak{}\char32\relax\allowbreak{}e\allowbreak{}n\allowbreak{}c\allowbreak{}o\allowbreak{}d\allowbreak{}e\allowbreak{}(\allowbreak{}s\allowbreak{}e\allowbreak{}l\allowbreak{}f\allowbreak{},\allowbreak{}\char32\relax\allowbreak{}t\allowbreak{}e\allowbreak{}x\allowbreak{}t\allowbreak{},\allowbreak{}\char32\relax\allowbreak{}a\allowbreak{}l\allowbreak{}l\allowbreak{}o\allowbreak{}w\allowbreak{}e\allowbreak{}d\allowbreak{}\char95\relax\allowbreak{}s\allowbreak{}p\allowbreak{}e\allowbreak{}c\allowbreak{}i\allowbreak{}a\allowbreak{}l\allowbreak{}=\allowbreak{}"\allowbreak{}n\allowbreak{}o\allowbreak{}n\allowbreak{}e\allowbreak{}\char95\relax\allowbreak{}r\allowbreak{}a\allowbreak{}i\allowbreak{}s\allowbreak{}e\allowbreak{}"\allowbreak{})\allowbreak{}:} & 定义对外编码入口，默认策略为 \BpeDocCodeInline{none\char95\relax\allowbreak{}raise}。 \\[3pt]
\SourceLineNumber{132} & \BpeDocCodeCell{\char32\relax\allowbreak{}\char32\relax\allowbreak{}\char32\relax\allowbreak{}\char32\relax\allowbreak{}\char32\relax\allowbreak{}\char32\relax\allowbreak{}\char32\relax\allowbreak{}\char32\relax\allowbreak{}s\allowbreak{}p\allowbreak{}e\allowbreak{}c\allowbreak{}i\allowbreak{}a\allowbreak{}l\allowbreak{}\char32\relax\allowbreak{}=\allowbreak{}\char32\relax\allowbreak{}N\allowbreak{}o\allowbreak{}n\allowbreak{}e} & 先把待启用的特殊 \BpeDocCodeInline{token} 表置为 \BpeDocCodeInline{None}。 \\[3pt]
\SourceLineNumber{133} & \BpeDocCodeCell{\char32\relax\allowbreak{}\char32\relax\allowbreak{}\char32\relax\allowbreak{}\char32\relax\allowbreak{}\char32\relax\allowbreak{}\char32\relax\allowbreak{}\char32\relax\allowbreak{}\char32\relax\allowbreak{}i\allowbreak{}f\allowbreak{}\char32\relax\allowbreak{}a\allowbreak{}l\allowbreak{}l\allowbreak{}o\allowbreak{}w\allowbreak{}e\allowbreak{}d\allowbreak{}\char95\relax\allowbreak{}s\allowbreak{}p\allowbreak{}e\allowbreak{}c\allowbreak{}i\allowbreak{}a\allowbreak{}l\allowbreak{}\char32\relax\allowbreak{}=\allowbreak{}=\allowbreak{}\char32\relax\allowbreak{}"\allowbreak{}a\allowbreak{}l\allowbreak{}l\allowbreak{}"\allowbreak{}:} & 判断是否要求启用所有特殊 \BpeDocCodeInline{token}。 \\[3pt]
\SourceLineNumber{134} & \BpeDocCodeCell{\char32\relax\allowbreak{}\char32\relax\allowbreak{}\char32\relax\allowbreak{}\char32\relax\allowbreak{}\char32\relax\allowbreak{}\char32\relax\allowbreak{}\char32\relax\allowbreak{}\char32\relax\allowbreak{}\char32\relax\allowbreak{}\char32\relax\allowbreak{}\char32\relax\allowbreak{}\char32\relax\allowbreak{}s\allowbreak{}p\allowbreak{}e\allowbreak{}c\allowbreak{}i\allowbreak{}a\allowbreak{}l\allowbreak{}\char32\relax\allowbreak{}=\allowbreak{}\char32\relax\allowbreak{}s\allowbreak{}e\allowbreak{}l\allowbreak{}f\allowbreak{}.\allowbreak{}s\allowbreak{}p\allowbreak{}e\allowbreak{}c\allowbreak{}i\allowbreak{}a\allowbreak{}l\allowbreak{}\char95\relax\allowbreak{}t\allowbreak{}o\allowbreak{}k\allowbreak{}e\allowbreak{}n\allowbreak{}s} & 直接使用全部已登记特殊 \BpeDocCodeInline{token}。 \\[3pt]
\SourceLineNumber{135} & \BpeDocCodeCell{\char32\relax\allowbreak{}\char32\relax\allowbreak{}\char32\relax\allowbreak{}\char32\relax\allowbreak{}\char32\relax\allowbreak{}\char32\relax\allowbreak{}\char32\relax\allowbreak{}\char32\relax\allowbreak{}e\allowbreak{}l\allowbreak{}i\allowbreak{}f\allowbreak{}\char32\relax\allowbreak{}a\allowbreak{}l\allowbreak{}l\allowbreak{}o\allowbreak{}w\allowbreak{}e\allowbreak{}d\allowbreak{}\char95\relax\allowbreak{}s\allowbreak{}p\allowbreak{}e\allowbreak{}c\allowbreak{}i\allowbreak{}a\allowbreak{}l\allowbreak{}\char32\relax\allowbreak{}=\allowbreak{}=\allowbreak{}\char32\relax\allowbreak{}"\allowbreak{}n\allowbreak{}o\allowbreak{}n\allowbreak{}e\allowbreak{}"\allowbreak{}:} & 判断是否要求把一切当作普通文本。 \\[3pt]
\SourceLineNumber{136} & \BpeDocCodeCell{\char32\relax\allowbreak{}\char32\relax\allowbreak{}\char32\relax\allowbreak{}\char32\relax\allowbreak{}\char32\relax\allowbreak{}\char32\relax\allowbreak{}\char32\relax\allowbreak{}\char32\relax\allowbreak{}\char32\relax\allowbreak{}\char32\relax\allowbreak{}\char32\relax\allowbreak{}\char32\relax\allowbreak{}s\allowbreak{}p\allowbreak{}e\allowbreak{}c\allowbreak{}i\allowbreak{}a\allowbreak{}l\allowbreak{}\char32\relax\allowbreak{}=\allowbreak{}\char32\relax\allowbreak{}\char123\relax\allowbreak{}\char125\relax} & 使用空表，不启用任何特殊 \BpeDocCodeInline{token}。 \\[3pt]
\SourceLineNumber{137} & \BpeDocCodeCell{\char32\relax\allowbreak{}\char32\relax\allowbreak{}\char32\relax\allowbreak{}\char32\relax\allowbreak{}\char32\relax\allowbreak{}\char32\relax\allowbreak{}\char32\relax\allowbreak{}\char32\relax\allowbreak{}e\allowbreak{}l\allowbreak{}i\allowbreak{}f\allowbreak{}\char32\relax\allowbreak{}a\allowbreak{}l\allowbreak{}l\allowbreak{}o\allowbreak{}w\allowbreak{}e\allowbreak{}d\allowbreak{}\char95\relax\allowbreak{}s\allowbreak{}p\allowbreak{}e\allowbreak{}c\allowbreak{}i\allowbreak{}a\allowbreak{}l\allowbreak{}\char32\relax\allowbreak{}=\allowbreak{}=\allowbreak{}\char32\relax\allowbreak{}"\allowbreak{}n\allowbreak{}o\allowbreak{}n\allowbreak{}e\allowbreak{}\char95\relax\allowbreak{}r\allowbreak{}a\allowbreak{}i\allowbreak{}s\allowbreak{}e\allowbreak{}"\allowbreak{}:} & 判断是否采用默认的拒绝标记策略。 \\[3pt]
\SourceLineNumber{138} & \BpeDocCodeCell{\char32\relax\allowbreak{}\char32\relax\allowbreak{}\char32\relax\allowbreak{}\char32\relax\allowbreak{}\char32\relax\allowbreak{}\char32\relax\allowbreak{}\char32\relax\allowbreak{}\char32\relax\allowbreak{}\char32\relax\allowbreak{}\char32\relax\allowbreak{}\char32\relax\allowbreak{}\char32\relax\allowbreak{}s\allowbreak{}p\allowbreak{}e\allowbreak{}c\allowbreak{}i\allowbreak{}a\allowbreak{}l\allowbreak{}\char32\relax\allowbreak{}=\allowbreak{}\char32\relax\allowbreak{}\char123\relax\allowbreak{}\char125\relax} & 仍使用空表，后续走普通编码。 \\[3pt]
\SourceLineNumber{139} & \BpeDocCodeCell{\char32\relax\allowbreak{}\char32\relax\allowbreak{}\char32\relax\allowbreak{}\char32\relax\allowbreak{}\char32\relax\allowbreak{}\char32\relax\allowbreak{}\char32\relax\allowbreak{}\char32\relax\allowbreak{}\char32\relax\allowbreak{}\char32\relax\allowbreak{}\char32\relax\allowbreak{}\char32\relax\allowbreak{}a\allowbreak{}s\allowbreak{}s\allowbreak{}e\allowbreak{}r\allowbreak{}t\allowbreak{}\char32\relax\allowbreak{}a\allowbreak{}l\allowbreak{}l\allowbreak{}(\allowbreak{}t\allowbreak{}o\allowbreak{}k\allowbreak{}e\allowbreak{}n\allowbreak{}\char32\relax\allowbreak{}n\allowbreak{}o\allowbreak{}t\allowbreak{}\char32\relax\allowbreak{}i\allowbreak{}n\allowbreak{}\char32\relax\allowbreak{}t\allowbreak{}e\allowbreak{}x\allowbreak{}t\allowbreak{}\char32\relax\allowbreak{}f\allowbreak{}o\allowbreak{}r\allowbreak{}\char32\relax\allowbreak{}t\allowbreak{}o\allowbreak{}k\allowbreak{}e\allowbreak{}n\allowbreak{}\char32\relax\allowbreak{}i\allowbreak{}n\allowbreak{}\char32\relax\allowbreak{}s\allowbreak{}e\allowbreak{}l\allowbreak{}f\allowbreak{}.\allowbreak{}s\allowbreak{}p\allowbreak{}e\allowbreak{}c\allowbreak{}i\allowbreak{}a\allowbreak{}l\allowbreak{}\char95\relax\allowbreak{}t\allowbreak{}o\allowbreak{}k\allowbreak{}e\allowbreak{}n\allowbreak{}s\allowbreak{})} & 检查所有已登记特殊字符串都没有出现在输入中；只要出现一个，断言失败。 \\[3pt]
\SourceLineNumber{140} & \BpeDocCodeCell{\char32\relax\allowbreak{}\char32\relax\allowbreak{}\char32\relax\allowbreak{}\char32\relax\allowbreak{}\char32\relax\allowbreak{}\char32\relax\allowbreak{}\char32\relax\allowbreak{}\char32\relax\allowbreak{}e\allowbreak{}l\allowbreak{}i\allowbreak{}f\allowbreak{}\char32\relax\allowbreak{}i\allowbreak{}s\allowbreak{}i\allowbreak{}n\allowbreak{}s\allowbreak{}t\allowbreak{}a\allowbreak{}n\allowbreak{}c\allowbreak{}e\allowbreak{}(\allowbreak{}a\allowbreak{}l\allowbreak{}l\allowbreak{}o\allowbreak{}w\allowbreak{}e\allowbreak{}d\allowbreak{}\char95\relax\allowbreak{}s\allowbreak{}p\allowbreak{}e\allowbreak{}c\allowbreak{}i\allowbreak{}a\allowbreak{}l\allowbreak{},\allowbreak{}\char32\relax\allowbreak{}s\allowbreak{}e\allowbreak{}t\allowbreak{})\allowbreak{}:} & 判断参数是否为 Python 集合 \BpeDocCodeInline{set}。 \\[3pt]
\SourceLineNumber{141} & \BpeDocCodeCell{\char32\relax\allowbreak{}\char32\relax\allowbreak{}\char32\relax\allowbreak{}\char32\relax\allowbreak{}\char32\relax\allowbreak{}\char32\relax\allowbreak{}\char32\relax\allowbreak{}\char32\relax\allowbreak{}\char32\relax\allowbreak{}\char32\relax\allowbreak{}\char32\relax\allowbreak{}\char32\relax\allowbreak{}s\allowbreak{}p\allowbreak{}e\allowbreak{}c\allowbreak{}i\allowbreak{}a\allowbreak{}l\allowbreak{}\char32\relax\allowbreak{}=\allowbreak{}\char32\relax\allowbreak{}\char123\relax\allowbreak{}k\allowbreak{}:\allowbreak{}\char32\relax\allowbreak{}v\allowbreak{}\char32\relax\allowbreak{}f\allowbreak{}o\allowbreak{}r\allowbreak{}\char32\relax\allowbreak{}k\allowbreak{},\allowbreak{}\char32\relax\allowbreak{}v\allowbreak{}\char32\relax\allowbreak{}i\allowbreak{}n\allowbreak{}\char32\relax\allowbreak{}s\allowbreak{}e\allowbreak{}l\allowbreak{}f\allowbreak{}.\allowbreak{}s\allowbreak{}p\allowbreak{}e\allowbreak{}c\allowbreak{}i\allowbreak{}a\allowbreak{}l\allowbreak{}\char95\relax\allowbreak{}t\allowbreak{}o\allowbreak{}k\allowbreak{}e\allowbreak{}n\allowbreak{}s\allowbreak{}.\allowbreak{}i\allowbreak{}t\allowbreak{}e\allowbreak{}m\allowbreak{}s\allowbreak{}(\allowbreak{})\allowbreak{}\char32\relax\allowbreak{}i\allowbreak{}f\allowbreak{}\char32\relax\allowbreak{}k\allowbreak{}\char32\relax\allowbreak{}i\allowbreak{}n\allowbreak{}\char32\relax\allowbreak{}a\allowbreak{}l\allowbreak{}l\allowbreak{}o\allowbreak{}w\allowbreak{}e\allowbreak{}d\allowbreak{}\char95\relax\allowbreak{}s\allowbreak{}p\allowbreak{}e\allowbreak{}c\allowbreak{}i\allowbreak{}a\allowbreak{}l\allowbreak{}\char125\relax} & 从已登记表中筛出集合允许的条目，生成本次使用的特殊表。 \\[3pt]
\SourceLineNumber{142} & \BpeDocCodeCell{\char32\relax\allowbreak{}\char32\relax\allowbreak{}\char32\relax\allowbreak{}\char32\relax\allowbreak{}\char32\relax\allowbreak{}\char32\relax\allowbreak{}\char32\relax\allowbreak{}\char32\relax\allowbreak{}e\allowbreak{}l\allowbreak{}s\allowbreak{}e\allowbreak{}:} & 其他参数类型或值进入这个分支。 \\[3pt]
\SourceLineNumber{143} & \BpeDocCodeCell{\char32\relax\allowbreak{}\char32\relax\allowbreak{}\char32\relax\allowbreak{}\char32\relax\allowbreak{}\char32\relax\allowbreak{}\char32\relax\allowbreak{}\char32\relax\allowbreak{}\char32\relax\allowbreak{}\char32\relax\allowbreak{}\char32\relax\allowbreak{}\char32\relax\allowbreak{}\char32\relax\allowbreak{}r\allowbreak{}a\allowbreak{}i\allowbreak{}s\allowbreak{}e\allowbreak{}\char32\relax\allowbreak{}V\allowbreak{}a\allowbreak{}l\allowbreak{}u\allowbreak{}e\allowbreak{}E\allowbreak{}r\allowbreak{}r\allowbreak{}o\allowbreak{}r\allowbreak{}(\allowbreak{}f\allowbreak{}"\allowbreak{}a\allowbreak{}l\allowbreak{}l\allowbreak{}o\allowbreak{}w\allowbreak{}e\allowbreak{}d\allowbreak{}\char95\relax\allowbreak{}s\allowbreak{}p\allowbreak{}e\allowbreak{}c\allowbreak{}i\allowbreak{}a\allowbreak{}l\allowbreak{}=\allowbreak{}\char123\relax\allowbreak{}a\allowbreak{}l\allowbreak{}l\allowbreak{}o\allowbreak{}w\allowbreak{}e\allowbreak{}d\allowbreak{}\char95\relax\allowbreak{}s\allowbreak{}p\allowbreak{}e\allowbreak{}c\allowbreak{}i\allowbreak{}a\allowbreak{}l\allowbreak{}\char125\relax\allowbreak{}\char32\relax\allowbreak{}n\allowbreak{}o\allowbreak{}t\allowbreak{}\char32\relax\allowbreak{}u\allowbreak{}n\allowbreak{}d\allowbreak{}e\allowbreak{}r\allowbreak{}s\allowbreak{}t\allowbreak{}o\allowbreak{}o\allowbreak{}d\allowbreak{}"\allowbreak{})} & 抛出参数错误，并在信息中显示不被理解的值。 \\[3pt]
\SourceLineNumber{144} & \BpeDocCodeCell{\char32\relax\allowbreak{}\char32\relax\allowbreak{}\char32\relax\allowbreak{}\char32\relax\allowbreak{}\char32\relax\allowbreak{}\char32\relax\allowbreak{}\char32\relax\allowbreak{}\char32\relax\allowbreak{}i\allowbreak{}f\allowbreak{}\char32\relax\allowbreak{}n\allowbreak{}o\allowbreak{}t\allowbreak{}\char32\relax\allowbreak{}s\allowbreak{}p\allowbreak{}e\allowbreak{}c\allowbreak{}i\allowbreak{}a\allowbreak{}l\allowbreak{}:} & 若启用表为空，就无需做特殊文本切分。 \\[3pt]
\SourceLineNumber{146} & \BpeDocCodeCell{\char32\relax\allowbreak{}\char32\relax\allowbreak{}\char32\relax\allowbreak{}\char32\relax\allowbreak{}\char32\relax\allowbreak{}\char32\relax\allowbreak{}\char32\relax\allowbreak{}\char32\relax\allowbreak{}\char32\relax\allowbreak{}\char32\relax\allowbreak{}\char32\relax\allowbreak{}\char32\relax\allowbreak{}r\allowbreak{}e\allowbreak{}t\allowbreak{}u\allowbreak{}r\allowbreak{}n\allowbreak{}\char32\relax\allowbreak{}s\allowbreak{}e\allowbreak{}l\allowbreak{}f\allowbreak{}.\allowbreak{}e\allowbreak{}n\allowbreak{}c\allowbreak{}o\allowbreak{}d\allowbreak{}e\allowbreak{}\char95\relax\allowbreak{}o\allowbreak{}r\allowbreak{}d\allowbreak{}i\allowbreak{}n\allowbreak{}a\allowbreak{}r\allowbreak{}y\allowbreak{}(\allowbreak{}t\allowbreak{}e\allowbreak{}x\allowbreak{}t\allowbreak{})} & 直接返回 \BpeDocCodeInline{encode\char95\relax\allowbreak{}ordinary} 的结果，结束当前函数。 \\[3pt]
\SourceLineNumber{152} & \BpeDocCodeCell{\char32\relax\allowbreak{}\char32\relax\allowbreak{}\char32\relax\allowbreak{}\char32\relax\allowbreak{}\char32\relax\allowbreak{}\char32\relax\allowbreak{}\char32\relax\allowbreak{}\char32\relax\allowbreak{}s\allowbreak{}p\allowbreak{}e\allowbreak{}c\allowbreak{}i\allowbreak{}a\allowbreak{}l\allowbreak{}\char95\relax\allowbreak{}p\allowbreak{}a\allowbreak{}t\allowbreak{}t\allowbreak{}e\allowbreak{}r\allowbreak{}n\allowbreak{}\char32\relax\allowbreak{}=\allowbreak{}\char32\relax\allowbreak{}"\allowbreak{}(\allowbreak{}"\allowbreak{}\char32\relax\allowbreak{}+\allowbreak{}\char32\relax\allowbreak{}"\allowbreak{}|\allowbreak{}"\allowbreak{}.\allowbreak{}j\allowbreak{}o\allowbreak{}i\allowbreak{}n\allowbreak{}(\allowbreak{}r\allowbreak{}e\allowbreak{}.\allowbreak{}e\allowbreak{}s\allowbreak{}c\allowbreak{}a\allowbreak{}p\allowbreak{}e\allowbreak{}(\allowbreak{}k\allowbreak{})\allowbreak{}\char32\relax\allowbreak{}f\allowbreak{}o\allowbreak{}r\allowbreak{}\char32\relax\allowbreak{}k\allowbreak{}\char32\relax\allowbreak{}i\allowbreak{}n\allowbreak{}\char32\relax\allowbreak{}s\allowbreak{}p\allowbreak{}e\allowbreak{}c\allowbreak{}i\allowbreak{}a\allowbreak{}l\allowbreak{})\allowbreak{}\char32\relax\allowbreak{}+\allowbreak{}\char32\relax\allowbreak{}"\allowbreak{})\allowbreak{}"} & 先转义各特殊字符串，避免其符号被当作正则语法；用 | 表示“任一标记”，再加括号保留匹配内容。 \\[3pt]
\SourceLineNumber{153} & \BpeDocCodeCell{\char32\relax\allowbreak{}\char32\relax\allowbreak{}\char32\relax\allowbreak{}\char32\relax\allowbreak{}\char32\relax\allowbreak{}\char32\relax\allowbreak{}\char32\relax\allowbreak{}\char32\relax\allowbreak{}s\allowbreak{}p\allowbreak{}e\allowbreak{}c\allowbreak{}i\allowbreak{}a\allowbreak{}l\allowbreak{}\char95\relax\allowbreak{}c\allowbreak{}h\allowbreak{}u\allowbreak{}n\allowbreak{}k\allowbreak{}s\allowbreak{}\char32\relax\allowbreak{}=\allowbreak{}\char32\relax\allowbreak{}r\allowbreak{}e\allowbreak{}.\allowbreak{}s\allowbreak{}p\allowbreak{}l\allowbreak{}i\allowbreak{}t\allowbreak{}(\allowbreak{}s\allowbreak{}p\allowbreak{}e\allowbreak{}c\allowbreak{}i\allowbreak{}a\allowbreak{}l\allowbreak{}\char95\relax\allowbreak{}p\allowbreak{}a\allowbreak{}t\allowbreak{}t\allowbreak{}e\allowbreak{}r\allowbreak{}n\allowbreak{},\allowbreak{}\char32\relax\allowbreak{}t\allowbreak{}e\allowbreak{}x\allowbreak{}t\allowbreak{})} & 按这个模式切开输入，得到普通片段和特殊片段交替组成的列表。 \\[3pt]
\SourceLineNumber{156} & \BpeDocCodeCell{\char32\relax\allowbreak{}\char32\relax\allowbreak{}\char32\relax\allowbreak{}\char32\relax\allowbreak{}\char32\relax\allowbreak{}\char32\relax\allowbreak{}\char32\relax\allowbreak{}\char32\relax\allowbreak{}i\allowbreak{}d\allowbreak{}s\allowbreak{}\char32\relax\allowbreak{}=\allowbreak{}\char32\relax\allowbreak{}[\allowbreak{}]} & 建立最终编号列表。 \\[3pt]
\SourceLineNumber{157} & \BpeDocCodeCell{\char32\relax\allowbreak{}\char32\relax\allowbreak{}\char32\relax\allowbreak{}\char32\relax\allowbreak{}\char32\relax\allowbreak{}\char32\relax\allowbreak{}\char32\relax\allowbreak{}\char32\relax\allowbreak{}f\allowbreak{}o\allowbreak{}r\allowbreak{}\char32\relax\allowbreak{}p\allowbreak{}a\allowbreak{}r\allowbreak{}t\allowbreak{}\char32\relax\allowbreak{}i\allowbreak{}n\allowbreak{}\char32\relax\allowbreak{}s\allowbreak{}p\allowbreak{}e\allowbreak{}c\allowbreak{}i\allowbreak{}a\allowbreak{}l\allowbreak{}\char95\relax\allowbreak{}c\allowbreak{}h\allowbreak{}u\allowbreak{}n\allowbreak{}k\allowbreak{}s\allowbreak{}:} & 按原顺序逐个处理片段。 \\[3pt]
\SourceLineNumber{158} & \BpeDocCodeCell{\char32\relax\allowbreak{}\char32\relax\allowbreak{}\char32\relax\allowbreak{}\char32\relax\allowbreak{}\char32\relax\allowbreak{}\char32\relax\allowbreak{}\char32\relax\allowbreak{}\char32\relax\allowbreak{}\char32\relax\allowbreak{}\char32\relax\allowbreak{}\char32\relax\allowbreak{}\char32\relax\allowbreak{}i\allowbreak{}f\allowbreak{}\char32\relax\allowbreak{}p\allowbreak{}a\allowbreak{}r\allowbreak{}t\allowbreak{}\char32\relax\allowbreak{}i\allowbreak{}n\allowbreak{}\char32\relax\allowbreak{}s\allowbreak{}p\allowbreak{}e\allowbreak{}c\allowbreak{}i\allowbreak{}a\allowbreak{}l\allowbreak{}:} & 若该片段正好是启用表中的特殊字符串，走特殊分支。 \\[3pt]
\SourceLineNumber{160} & \BpeDocCodeCell{\char32\relax\allowbreak{}\char32\relax\allowbreak{}\char32\relax\allowbreak{}\char32\relax\allowbreak{}\char32\relax\allowbreak{}\char32\relax\allowbreak{}\char32\relax\allowbreak{}\char32\relax\allowbreak{}\char32\relax\allowbreak{}\char32\relax\allowbreak{}\char32\relax\allowbreak{}\char32\relax\allowbreak{}\char32\relax\allowbreak{}\char32\relax\allowbreak{}\char32\relax\allowbreak{}\char32\relax\allowbreak{}i\allowbreak{}d\allowbreak{}s\allowbreak{}.\allowbreak{}a\allowbreak{}p\allowbreak{}p\allowbreak{}e\allowbreak{}n\allowbreak{}d\allowbreak{}(\allowbreak{}s\allowbreak{}p\allowbreak{}e\allowbreak{}c\allowbreak{}i\allowbreak{}a\allowbreak{}l\allowbreak{}[\allowbreak{}p\allowbreak{}a\allowbreak{}r\allowbreak{}t\allowbreak{}]\allowbreak{})} & 把它对应的单个特殊 \BpeDocCodeInline{token} ID 加入列表，不对它做普通 BPE。 \\[3pt]
\SourceLineNumber{161} & \BpeDocCodeCell{\char32\relax\allowbreak{}\char32\relax\allowbreak{}\char32\relax\allowbreak{}\char32\relax\allowbreak{}\char32\relax\allowbreak{}\char32\relax\allowbreak{}\char32\relax\allowbreak{}\char32\relax\allowbreak{}\char32\relax\allowbreak{}\char32\relax\allowbreak{}\char32\relax\allowbreak{}\char32\relax\allowbreak{}e\allowbreak{}l\allowbreak{}s\allowbreak{}e\allowbreak{}:} & 否则按普通片段处理。 \\[3pt]
\SourceLineNumber{163} & \BpeDocCodeCell{\char32\relax\allowbreak{}\char32\relax\allowbreak{}\char32\relax\allowbreak{}\char32\relax\allowbreak{}\char32\relax\allowbreak{}\char32\relax\allowbreak{}\char32\relax\allowbreak{}\char32\relax\allowbreak{}\char32\relax\allowbreak{}\char32\relax\allowbreak{}\char32\relax\allowbreak{}\char32\relax\allowbreak{}\char32\relax\allowbreak{}\char32\relax\allowbreak{}\char32\relax\allowbreak{}\char32\relax\allowbreak{}i\allowbreak{}d\allowbreak{}s\allowbreak{}.\allowbreak{}e\allowbreak{}x\allowbreak{}t\allowbreak{}e\allowbreak{}n\allowbreak{}d\allowbreak{}(\allowbreak{}s\allowbreak{}e\allowbreak{}l\allowbreak{}f\allowbreak{}.\allowbreak{}e\allowbreak{}n\allowbreak{}c\allowbreak{}o\allowbreak{}d\allowbreak{}e\allowbreak{}\char95\relax\allowbreak{}o\allowbreak{}r\allowbreak{}d\allowbreak{}i\allowbreak{}n\allowbreak{}a\allowbreak{}r\allowbreak{}y\allowbreak{}(\allowbreak{}p\allowbreak{}a\allowbreak{}r\allowbreak{}t\allowbreak{})\allowbreak{})} & 调用 \BpeDocCodeInline{encode\char95\relax\allowbreak{}ordinary}，并把返回的各编号接到总列表末尾。 \\[3pt]
\SourceLineNumber{164} & \BpeDocCodeCell{\char32\relax\allowbreak{}\char32\relax\allowbreak{}\char32\relax\allowbreak{}\char32\relax\allowbreak{}\char32\relax\allowbreak{}\char32\relax\allowbreak{}\char32\relax\allowbreak{}\char32\relax\allowbreak{}r\allowbreak{}e\allowbreak{}t\allowbreak{}u\allowbreak{}r\allowbreak{}n\allowbreak{}\char32\relax\allowbreak{}i\allowbreak{}d\allowbreak{}s} & 返回最终编号序列。 \\[3pt]
\end{longtable}
\endgroup


对一个空的普通片段，\BpeDocCodeInline{encode\char95\relax\allowbreak{}ordinary} 会返回空列表，因此 \BpeDocCodeInline{re.\allowbreak{}split} 产生的首尾空片段不会额外制造 \BpeDocCodeInline{token}。默认 \BpeDocCodeInline{none\char95\relax\allowbreak{}raise} 在这份代码里用 \BpeDocCodeInline{assert} 实现；报告准确写“触发断言”即可，不需要把它描述成 \BpeDocCodeInline{ValueError}。


\Needspace{9\baselineskip}
\BpeDocSubsubsection{\BpeDocCodeInline{regex.\allowbreak{}py}：普通类如何解码特殊编号，78–90 行}{regex.py：普通类如何解码特殊编号，78–90 行}


\Needspace{4\baselineskip}
\begingroup\fontsize{8.9}{12.2}\selectfont\setstretch{1}
\setlength{\BpeDocTableWidth}{\dimexpr\linewidth-6\tabcolsep\relax}
\renewcommand{\arraystretch}{1.08}
{\small \BpeDocCodeInline{minbpe/\allowbreak{}regex.\allowbreak{}py}，第 78–90 行（仅列执行语句）}

\begin{longtable}{>{\raggedleft\arraybackslash}p{0.065\BpeDocTableWidth}>{\raggedright\arraybackslash}p{0.475\BpeDocTableWidth}>{\raggedright\arraybackslash}p{0.460\BpeDocTableWidth}}
\toprule
\textbf{行号} & \textbf{原代码} & \textbf{逐行解释} \\
\midrule\endfirsthead
\toprule
\textbf{行号} & \textbf{原代码} & \textbf{逐行解释} \\
\midrule\endhead
\midrule\endfoot
\bottomrule\endlastfoot
\SourceLineNumber{78} & \BpeDocCodeCell{\char32\relax\allowbreak{}\char32\relax\allowbreak{}\char32\relax\allowbreak{}\char32\relax\allowbreak{}d\allowbreak{}e\allowbreak{}f\allowbreak{}\char32\relax\allowbreak{}d\allowbreak{}e\allowbreak{}c\allowbreak{}o\allowbreak{}d\allowbreak{}e\allowbreak{}(\allowbreak{}s\allowbreak{}e\allowbreak{}l\allowbreak{}f\allowbreak{},\allowbreak{}\char32\relax\allowbreak{}i\allowbreak{}d\allowbreak{}s\allowbreak{})\allowbreak{}:} & 定义 \BpeDocCodeInline{RegexTokenizer} 的解码方法。 \\[3pt]
\SourceLineNumber{80} & \BpeDocCodeCell{\char32\relax\allowbreak{}\char32\relax\allowbreak{}\char32\relax\allowbreak{}\char32\relax\allowbreak{}\char32\relax\allowbreak{}\char32\relax\allowbreak{}\char32\relax\allowbreak{}\char32\relax\allowbreak{}p\allowbreak{}a\allowbreak{}r\allowbreak{}t\allowbreak{}\char95\relax\allowbreak{}b\allowbreak{}y\allowbreak{}t\allowbreak{}e\allowbreak{}s\allowbreak{}\char32\relax\allowbreak{}=\allowbreak{}\char32\relax\allowbreak{}[\allowbreak{}]} & 准备收集每个编号对应的字节片段。 \\[3pt]
\SourceLineNumber{81} & \BpeDocCodeCell{\char32\relax\allowbreak{}\char32\relax\allowbreak{}\char32\relax\allowbreak{}\char32\relax\allowbreak{}\char32\relax\allowbreak{}\char32\relax\allowbreak{}\char32\relax\allowbreak{}\char32\relax\allowbreak{}f\allowbreak{}o\allowbreak{}r\allowbreak{}\char32\relax\allowbreak{}i\allowbreak{}d\allowbreak{}x\allowbreak{}\char32\relax\allowbreak{}i\allowbreak{}n\allowbreak{}\char32\relax\allowbreak{}i\allowbreak{}d\allowbreak{}s\allowbreak{}:} & 按顺序遍历输入编号。 \\[3pt]
\SourceLineNumber{82} & \BpeDocCodeCell{\char32\relax\allowbreak{}\char32\relax\allowbreak{}\char32\relax\allowbreak{}\char32\relax\allowbreak{}\char32\relax\allowbreak{}\char32\relax\allowbreak{}\char32\relax\allowbreak{}\char32\relax\allowbreak{}\char32\relax\allowbreak{}\char32\relax\allowbreak{}\char32\relax\allowbreak{}\char32\relax\allowbreak{}i\allowbreak{}f\allowbreak{}\char32\relax\allowbreak{}i\allowbreak{}d\allowbreak{}x\allowbreak{}\char32\relax\allowbreak{}i\allowbreak{}n\allowbreak{}\char32\relax\allowbreak{}s\allowbreak{}e\allowbreak{}l\allowbreak{}f\allowbreak{}.\allowbreak{}v\allowbreak{}o\allowbreak{}c\allowbreak{}a\allowbreak{}b\allowbreak{}:} & 先判断编号是否在普通 BPE 词表中。 \\[3pt]
\SourceLineNumber{83} & \BpeDocCodeCell{\char32\relax\allowbreak{}\char32\relax\allowbreak{}\char32\relax\allowbreak{}\char32\relax\allowbreak{}\char32\relax\allowbreak{}\char32\relax\allowbreak{}\char32\relax\allowbreak{}\char32\relax\allowbreak{}\char32\relax\allowbreak{}\char32\relax\allowbreak{}\char32\relax\allowbreak{}\char32\relax\allowbreak{}\char32\relax\allowbreak{}\char32\relax\allowbreak{}\char32\relax\allowbreak{}\char32\relax\allowbreak{}p\allowbreak{}a\allowbreak{}r\allowbreak{}t\allowbreak{}\char95\relax\allowbreak{}b\allowbreak{}y\allowbreak{}t\allowbreak{}e\allowbreak{}s\allowbreak{}.\allowbreak{}a\allowbreak{}p\allowbreak{}p\allowbreak{}e\allowbreak{}n\allowbreak{}d\allowbreak{}(\allowbreak{}s\allowbreak{}e\allowbreak{}l\allowbreak{}f\allowbreak{}.\allowbreak{}v\allowbreak{}o\allowbreak{}c\allowbreak{}a\allowbreak{}b\allowbreak{}[\allowbreak{}i\allowbreak{}d\allowbreak{}x\allowbreak{}]\allowbreak{})} & 如果是，加入普通词表对应的字节串。 \\[3pt]
\SourceLineNumber{84} & \BpeDocCodeCell{\char32\relax\allowbreak{}\char32\relax\allowbreak{}\char32\relax\allowbreak{}\char32\relax\allowbreak{}\char32\relax\allowbreak{}\char32\relax\allowbreak{}\char32\relax\allowbreak{}\char32\relax\allowbreak{}\char32\relax\allowbreak{}\char32\relax\allowbreak{}\char32\relax\allowbreak{}\char32\relax\allowbreak{}e\allowbreak{}l\allowbreak{}i\allowbreak{}f\allowbreak{}\char32\relax\allowbreak{}i\allowbreak{}d\allowbreak{}x\allowbreak{}\char32\relax\allowbreak{}i\allowbreak{}n\allowbreak{}\char32\relax\allowbreak{}s\allowbreak{}e\allowbreak{}l\allowbreak{}f\allowbreak{}.\allowbreak{}i\allowbreak{}n\allowbreak{}v\allowbreak{}e\allowbreak{}r\allowbreak{}s\allowbreak{}e\allowbreak{}\char95\relax\allowbreak{}s\allowbreak{}p\allowbreak{}e\allowbreak{}c\allowbreak{}i\allowbreak{}a\allowbreak{}l\allowbreak{}\char95\relax\allowbreak{}t\allowbreak{}o\allowbreak{}k\allowbreak{}e\allowbreak{}n\allowbreak{}s\allowbreak{}:} & 否则判断它是否是已登记的特殊编号。 \\[3pt]
\SourceLineNumber{85} & \BpeDocCodeCell{\char32\relax\allowbreak{}\char32\relax\allowbreak{}\char32\relax\allowbreak{}\char32\relax\allowbreak{}\char32\relax\allowbreak{}\char32\relax\allowbreak{}\char32\relax\allowbreak{}\char32\relax\allowbreak{}\char32\relax\allowbreak{}\char32\relax\allowbreak{}\char32\relax\allowbreak{}\char32\relax\allowbreak{}\char32\relax\allowbreak{}\char32\relax\allowbreak{}\char32\relax\allowbreak{}\char32\relax\allowbreak{}p\allowbreak{}a\allowbreak{}r\allowbreak{}t\allowbreak{}\char95\relax\allowbreak{}b\allowbreak{}y\allowbreak{}t\allowbreak{}e\allowbreak{}s\allowbreak{}.\allowbreak{}a\allowbreak{}p\allowbreak{}p\allowbreak{}e\allowbreak{}n\allowbreak{}d\allowbreak{}(\allowbreak{}s\allowbreak{}e\allowbreak{}l\allowbreak{}f\allowbreak{}.\allowbreak{}i\allowbreak{}n\allowbreak{}v\allowbreak{}e\allowbreak{}r\allowbreak{}s\allowbreak{}e\allowbreak{}\char95\relax\allowbreak{}s\allowbreak{}p\allowbreak{}e\allowbreak{}c\allowbreak{}i\allowbreak{}a\allowbreak{}l\allowbreak{}\char95\relax\allowbreak{}t\allowbreak{}o\allowbreak{}k\allowbreak{}e\allowbreak{}n\allowbreak{}s\allowbreak{}[\allowbreak{}i\allowbreak{}d\allowbreak{}x\allowbreak{}]\allowbreak{}.\allowbreak{}e\allowbreak{}n\allowbreak{}c\allowbreak{}o\allowbreak{}d\allowbreak{}e\allowbreak{}(\allowbreak{}"\allowbreak{}u\allowbreak{}t\allowbreak{}f\allowbreak{}-\allowbreak{}8\allowbreak{}"\allowbreak{})\allowbreak{})} & 若是特殊编号，把对应特殊字符串转为 UTF-8 字节并加入。 \\[3pt]
\SourceLineNumber{86} & \BpeDocCodeCell{\char32\relax\allowbreak{}\char32\relax\allowbreak{}\char32\relax\allowbreak{}\char32\relax\allowbreak{}\char32\relax\allowbreak{}\char32\relax\allowbreak{}\char32\relax\allowbreak{}\char32\relax\allowbreak{}\char32\relax\allowbreak{}\char32\relax\allowbreak{}\char32\relax\allowbreak{}\char32\relax\allowbreak{}e\allowbreak{}l\allowbreak{}s\allowbreak{}e\allowbreak{}:} & 两种表都查不到时进入错误分支。 \\[3pt]
\SourceLineNumber{87} & \BpeDocCodeCell{\char32\relax\allowbreak{}\char32\relax\allowbreak{}\char32\relax\allowbreak{}\char32\relax\allowbreak{}\char32\relax\allowbreak{}\char32\relax\allowbreak{}\char32\relax\allowbreak{}\char32\relax\allowbreak{}\char32\relax\allowbreak{}\char32\relax\allowbreak{}\char32\relax\allowbreak{}\char32\relax\allowbreak{}\char32\relax\allowbreak{}\char32\relax\allowbreak{}\char32\relax\allowbreak{}\char32\relax\allowbreak{}r\allowbreak{}a\allowbreak{}i\allowbreak{}s\allowbreak{}e\allowbreak{}\char32\relax\allowbreak{}V\allowbreak{}a\allowbreak{}l\allowbreak{}u\allowbreak{}e\allowbreak{}E\allowbreak{}r\allowbreak{}r\allowbreak{}o\allowbreak{}r\allowbreak{}(\allowbreak{}f\allowbreak{}"\allowbreak{}i\allowbreak{}n\allowbreak{}v\allowbreak{}a\allowbreak{}l\allowbreak{}i\allowbreak{}d\allowbreak{}\char32\relax\allowbreak{}t\allowbreak{}o\allowbreak{}k\allowbreak{}e\allowbreak{}n\allowbreak{}\char32\relax\allowbreak{}i\allowbreak{}d\allowbreak{}:\allowbreak{}\char32\relax\allowbreak{}\char123\relax\allowbreak{}i\allowbreak{}d\allowbreak{}x\allowbreak{}\char125\relax\allowbreak{}"\allowbreak{})} & 抛出 \BpeDocCodeInline{ValueError}，指出未知编号。 \\[3pt]
\SourceLineNumber{88} & \BpeDocCodeCell{\char32\relax\allowbreak{}\char32\relax\allowbreak{}\char32\relax\allowbreak{}\char32\relax\allowbreak{}\char32\relax\allowbreak{}\char32\relax\allowbreak{}\char32\relax\allowbreak{}\char32\relax\allowbreak{}t\allowbreak{}e\allowbreak{}x\allowbreak{}t\allowbreak{}\char95\relax\allowbreak{}b\allowbreak{}y\allowbreak{}t\allowbreak{}e\allowbreak{}s\allowbreak{}\char32\relax\allowbreak{}=\allowbreak{}\char32\relax\allowbreak{}b\allowbreak{}"\allowbreak{}"\allowbreak{}.\allowbreak{}j\allowbreak{}o\allowbreak{}i\allowbreak{}n\allowbreak{}(\allowbreak{}p\allowbreak{}a\allowbreak{}r\allowbreak{}t\allowbreak{}\char95\relax\allowbreak{}b\allowbreak{}y\allowbreak{}t\allowbreak{}e\allowbreak{}s\allowbreak{})} & 拼接所有字节片段。 \\[3pt]
\SourceLineNumber{89} & \BpeDocCodeCell{\char32\relax\allowbreak{}\char32\relax\allowbreak{}\char32\relax\allowbreak{}\char32\relax\allowbreak{}\char32\relax\allowbreak{}\char32\relax\allowbreak{}\char32\relax\allowbreak{}\char32\relax\allowbreak{}t\allowbreak{}e\allowbreak{}x\allowbreak{}t\allowbreak{}\char32\relax\allowbreak{}=\allowbreak{}\char32\relax\allowbreak{}t\allowbreak{}e\allowbreak{}x\allowbreak{}t\allowbreak{}\char95\relax\allowbreak{}b\allowbreak{}y\allowbreak{}t\allowbreak{}e\allowbreak{}s\allowbreak{}.\allowbreak{}d\allowbreak{}e\allowbreak{}c\allowbreak{}o\allowbreak{}d\allowbreak{}e\allowbreak{}(\allowbreak{}"\allowbreak{}u\allowbreak{}t\allowbreak{}f\allowbreak{}-\allowbreak{}8\allowbreak{}"\allowbreak{},\allowbreak{}\char32\relax\allowbreak{}e\allowbreak{}r\allowbreak{}r\allowbreak{}o\allowbreak{}r\allowbreak{}s\allowbreak{}=\allowbreak{}"\allowbreak{}r\allowbreak{}e\allowbreak{}p\allowbreak{}l\allowbreak{}a\allowbreak{}c\allowbreak{}e\allowbreak{}"\allowbreak{})} & 对完整字节序列做 UTF-8 解码，非法序列用替代字符处理。 \\[3pt]
\SourceLineNumber{90} & \BpeDocCodeCell{\char32\relax\allowbreak{}\char32\relax\allowbreak{}\char32\relax\allowbreak{}\char32\relax\allowbreak{}\char32\relax\allowbreak{}\char32\relax\allowbreak{}\char32\relax\allowbreak{}\char32\relax\allowbreak{}r\allowbreak{}e\allowbreak{}t\allowbreak{}u\allowbreak{}r\allowbreak{}n\allowbreak{}\char32\relax\allowbreak{}t\allowbreak{}e\allowbreak{}x\allowbreak{}t} & 返回文本。 \\[3pt]
\end{longtable}
\endgroup


这段代码说明“支持特殊 \BpeDocCodeInline{token} 解码”需要专门分支。后面 \BpeDocCodeInline{GPT4Tokenizer} 改写了 \BpeDocCodeInline{decode}，不能仅因它继承 \BpeDocCodeInline{RegexTokenizer}，就推断它仍完整保留了这段特殊编号处理。这个细节会影响往返测试。


\Needspace{9\baselineskip}
\BpeDocSubsection{\BpeDocCodeInline{gpt4.\allowbreak{}py}：它复现的是已有 \BpeDocCodeInline{cl100k\char95\relax\allowbreak{}base} 分词规则}{gpt4.py：它复现的是已有 cl100k\_base 分词规则}


MinBPE 中名为 \BpeDocCodeInline{GPT4Tokenizer} 的类读取 tiktoken 已有的 \BpeDocCodeInline{cl100k\char95\relax\allowbreak{}base} 编码表，再把规则接到 \BpeDocCodeInline{RegexTokenizer} 流程上。本项目研究这一固定编码的分词复现、规则恢复与字节编号映射；它不等同于训练 GPT-4 模型，也不代表所有后续 GPT 系列模型共享同一分词器。


\begingroup\fontsize{9.3}{12.7}\selectfont\setstretch{1}
\setlength{\BpeDocTableWidth}{\dimexpr\linewidth-4\tabcolsep\relax}
\renewcommand{\arraystretch}{1.15}
\par\addvspace{0.45em}
\noindent\begin{minipage}{\linewidth}
\begin{tabular}{>{\raggedright\arraybackslash}p{0.290000\BpeDocTableWidth}>{\raggedright\arraybackslash}p{0.710000\BpeDocTableWidth}}
\toprule
\textbf{调用位置} & \textbf{实际发生的事} \\

\midrule
\BpeDocCodeInline{GPT4Tokenizer} 初始化 & 读取预训练编码表，恢复合并关系，准备字节编号映射和特殊 \BpeDocCodeInline{token} 表。 \\[3pt]
继承的 \BpeDocCodeInline{RegexTokenizer.\allowbreak{}encode} & 根据 \BpeDocCodeInline{allowed\char95\relax\allowbreak{}special} 处理特殊片段；普通片段交给 \BpeDocCodeInline{encode\char95\relax\allowbreak{}ordinary}。 \\[3pt]
继承的 \BpeDocCodeInline{encode\char95\relax\allowbreak{}ordinary} & 做 GPT-4 正则预切分，然后调用 \BpeDocCodeInline{self.\allowbreak{}\char95\relax\allowbreak{}encode\char95\relax\allowbreak{}chunk}。 \\[3pt]
\BpeDocCodeInline{GPT4Tokenizer.\allowbreak{}\char95\relax\allowbreak{}encode\char95\relax\allowbreak{}chunk} & 先把原始字节换成官方单字节 \BpeDocCodeInline{token} 编号，再调用父类的 BPE 编码。 \\[3pt]
\BpeDocCodeInline{GPT4Tokenizer.\allowbreak{}decode} & 先查词表拼接内部字节表示，再把单字节编号映射回原始字节，最后 UTF-8 解码。 \\[3pt]
\bottomrule
\end{tabular}
\end{minipage}\par\addvspace{0.65em}
\endgroup


\Needspace{9\baselineskip}
\BpeDocSubsubsection{\BpeDocCodeInline{gpt4.\allowbreak{}py}：初始化预训练规则，60–79 行}{gpt4.py：初始化预训练规则，60–79 行}


\Needspace{4\baselineskip}
\begingroup\fontsize{8.9}{12.2}\selectfont\setstretch{1}
\setlength{\BpeDocTableWidth}{\dimexpr\linewidth-6\tabcolsep\relax}
\renewcommand{\arraystretch}{1.08}
{\small \BpeDocCodeInline{minbpe/\allowbreak{}gpt4.\allowbreak{}py}，第 60–79 行（仅列执行语句）}

\begin{longtable}{>{\raggedleft\arraybackslash}p{0.065\BpeDocTableWidth}>{\raggedright\arraybackslash}p{0.475\BpeDocTableWidth}>{\raggedright\arraybackslash}p{0.460\BpeDocTableWidth}}
\toprule
\textbf{行号} & \textbf{原代码} & \textbf{逐行解释} \\
\midrule\endfirsthead
\toprule
\textbf{行号} & \textbf{原代码} & \textbf{逐行解释} \\
\midrule\endhead
\midrule\endfoot
\bottomrule\endlastfoot
\SourceLineNumber{60} & \BpeDocCodeCell{\char32\relax\allowbreak{}\char32\relax\allowbreak{}\char32\relax\allowbreak{}\char32\relax\allowbreak{}d\allowbreak{}e\allowbreak{}f\allowbreak{}\char32\relax\allowbreak{}\char95\relax\allowbreak{}\char95\relax\allowbreak{}i\allowbreak{}n\allowbreak{}i\allowbreak{}t\allowbreak{}\char95\relax\allowbreak{}\char95\relax\allowbreak{}(\allowbreak{}s\allowbreak{}e\allowbreak{}l\allowbreak{}f\allowbreak{})\allowbreak{}:} & 创建 \BpeDocCodeInline{GPT4Tokenizer} 对象时执行此初始化方法。 \\[3pt]
\SourceLineNumber{61} & \BpeDocCodeCell{\char32\relax\allowbreak{}\char32\relax\allowbreak{}\char32\relax\allowbreak{}\char32\relax\allowbreak{}\char32\relax\allowbreak{}\char32\relax\allowbreak{}\char32\relax\allowbreak{}\char32\relax\allowbreak{}s\allowbreak{}u\allowbreak{}p\allowbreak{}e\allowbreak{}r\allowbreak{}(\allowbreak{})\allowbreak{}.\allowbreak{}\char95\relax\allowbreak{}\char95\relax\allowbreak{}i\allowbreak{}n\allowbreak{}i\allowbreak{}t\allowbreak{}\char95\relax\allowbreak{}\char95\relax\allowbreak{}(\allowbreak{}p\allowbreak{}a\allowbreak{}t\allowbreak{}t\allowbreak{}e\allowbreak{}r\allowbreak{}n\allowbreak{}=\allowbreak{}G\allowbreak{}P\allowbreak{}T\allowbreak{}4\allowbreak{}\char95\relax\allowbreak{}S\allowbreak{}P\allowbreak{}L\allowbreak{}I\allowbreak{}T\allowbreak{}\char95\relax\allowbreak{}P\allowbreak{}A\allowbreak{}T\allowbreak{}T\allowbreak{}E\allowbreak{}R\allowbreak{}N\allowbreak{})} & 调用 \BpeDocCodeInline{RegexTokenizer} 初始化，并明确使用 \BpeDocCodeInline{GPT4\char95\relax\allowbreak{}SPLIT\char95\relax\allowbreak{}PATTERN}。 \\[3pt]
\SourceLineNumber{63} & \BpeDocCodeCell{\char32\relax\allowbreak{}\char32\relax\allowbreak{}\char32\relax\allowbreak{}\char32\relax\allowbreak{}\char32\relax\allowbreak{}\char32\relax\allowbreak{}\char32\relax\allowbreak{}\char32\relax\allowbreak{}e\allowbreak{}n\allowbreak{}c\allowbreak{}\char32\relax\allowbreak{}=\allowbreak{}\char32\relax\allowbreak{}t\allowbreak{}i\allowbreak{}k\allowbreak{}t\allowbreak{}o\allowbreak{}k\allowbreak{}e\allowbreak{}n\allowbreak{}.\allowbreak{}g\allowbreak{}e\allowbreak{}t\allowbreak{}\char95\relax\allowbreak{}e\allowbreak{}n\allowbreak{}c\allowbreak{}o\allowbreak{}d\allowbreak{}i\allowbreak{}n\allowbreak{}g\allowbreak{}(\allowbreak{}"\allowbreak{}c\allowbreak{}l\allowbreak{}1\allowbreak{}0\allowbreak{}0\allowbreak{}k\allowbreak{}\char95\relax\allowbreak{}b\allowbreak{}a\allowbreak{}s\allowbreak{}e\allowbreak{}"\allowbreak{})} & 从 tiktoken 获取 \BpeDocCodeInline{cl100k\char95\relax\allowbreak{}base} 编码器。这一步不是训练模型。 \\[3pt]
\SourceLineNumber{64} & \BpeDocCodeCell{\char32\relax\allowbreak{}\char32\relax\allowbreak{}\char32\relax\allowbreak{}\char32\relax\allowbreak{}\char32\relax\allowbreak{}\char32\relax\allowbreak{}\char32\relax\allowbreak{}\char32\relax\allowbreak{}m\allowbreak{}e\allowbreak{}r\allowbreak{}g\allowbreak{}e\allowbreak{}a\allowbreak{}b\allowbreak{}l\allowbreak{}e\allowbreak{}\char95\relax\allowbreak{}r\allowbreak{}a\allowbreak{}n\allowbreak{}k\allowbreak{}s\allowbreak{}\char32\relax\allowbreak{}=\allowbreak{}\char32\relax\allowbreak{}e\allowbreak{}n\allowbreak{}c\allowbreak{}.\allowbreak{}\char95\relax\allowbreak{}m\allowbreak{}e\allowbreak{}r\allowbreak{}g\allowbreak{}e\allowbreak{}a\allowbreak{}b\allowbreak{}l\allowbreak{}e\allowbreak{}\char95\relax\allowbreak{}r\allowbreak{}a\allowbreak{}n\allowbreak{}k\allowbreak{}s} & 读取其“字节串 \(\rightarrow\) 合并等级／编号”表。属性名以单下划线开头，属于实现细节，版本变化时需核对。 \\[3pt]
\SourceLineNumber{66} & \BpeDocCodeCell{\char32\relax\allowbreak{}\char32\relax\allowbreak{}\char32\relax\allowbreak{}\char32\relax\allowbreak{}\char32\relax\allowbreak{}\char32\relax\allowbreak{}\char32\relax\allowbreak{}\char32\relax\allowbreak{}s\allowbreak{}e\allowbreak{}l\allowbreak{}f\allowbreak{}.\allowbreak{}m\allowbreak{}e\allowbreak{}r\allowbreak{}g\allowbreak{}e\allowbreak{}s\allowbreak{}\char32\relax\allowbreak{}=\allowbreak{}\char32\relax\allowbreak{}r\allowbreak{}e\allowbreak{}c\allowbreak{}o\allowbreak{}v\allowbreak{}e\allowbreak{}r\allowbreak{}\char95\relax\allowbreak{}m\allowbreak{}e\allowbreak{}r\allowbreak{}g\allowbreak{}e\allowbreak{}s\allowbreak{}(\allowbreak{}m\allowbreak{}e\allowbreak{}r\allowbreak{}g\allowbreak{}e\allowbreak{}a\allowbreak{}b\allowbreak{}l\allowbreak{}e\allowbreak{}\char95\relax\allowbreak{}r\allowbreak{}a\allowbreak{}n\allowbreak{}k\allowbreak{}s\allowbreak{})} & 调用 \BpeDocCodeInline{recover\char95\relax\allowbreak{}merges}，生成 MinBPE 使用的“两个 \BpeDocCodeInline{token} ID \(\rightarrow\) 新 ID”表。 \\[3pt]
\SourceLineNumber{68} & \BpeDocCodeCell{\char32\relax\allowbreak{}\char32\relax\allowbreak{}\char32\relax\allowbreak{}\char32\relax\allowbreak{}\char32\relax\allowbreak{}\char32\relax\allowbreak{}\char32\relax\allowbreak{}\char32\relax\allowbreak{}v\allowbreak{}o\allowbreak{}c\allowbreak{}a\allowbreak{}b\allowbreak{}\char32\relax\allowbreak{}=\allowbreak{}\char32\relax\allowbreak{}\char123\relax\allowbreak{}i\allowbreak{}d\allowbreak{}x\allowbreak{}:\allowbreak{}\char32\relax\allowbreak{}b\allowbreak{}y\allowbreak{}t\allowbreak{}e\allowbreak{}s\allowbreak{}(\allowbreak{}[\allowbreak{}i\allowbreak{}d\allowbreak{}x\allowbreak{}]\allowbreak{})\allowbreak{}\char32\relax\allowbreak{}f\allowbreak{}o\allowbreak{}r\allowbreak{}\char32\relax\allowbreak{}i\allowbreak{}d\allowbreak{}x\allowbreak{}\char32\relax\allowbreak{}i\allowbreak{}n\allowbreak{}\char32\relax\allowbreak{}r\allowbreak{}a\allowbreak{}n\allowbreak{}g\allowbreak{}e\allowbreak{}(\allowbreak{}2\allowbreak{}5\allowbreak{}6\allowbreak{})\allowbreak{}\char125\relax} & 先把 0–255 编号各记为同数值单字节；注意这里暂时是内部编号表示，并非都等于实际原始字节。 \\[3pt]
\SourceLineNumber{69} & \BpeDocCodeCell{\char32\relax\allowbreak{}\char32\relax\allowbreak{}\char32\relax\allowbreak{}\char32\relax\allowbreak{}\char32\relax\allowbreak{}\char32\relax\allowbreak{}\char32\relax\allowbreak{}\char32\relax\allowbreak{}f\allowbreak{}o\allowbreak{}r\allowbreak{}\char32\relax\allowbreak{}(\allowbreak{}p\allowbreak{}0\allowbreak{},\allowbreak{}\char32\relax\allowbreak{}p\allowbreak{}1\allowbreak{})\allowbreak{},\allowbreak{}\char32\relax\allowbreak{}i\allowbreak{}d\allowbreak{}x\allowbreak{}\char32\relax\allowbreak{}i\allowbreak{}n\allowbreak{}\char32\relax\allowbreak{}s\allowbreak{}e\allowbreak{}l\allowbreak{}f\allowbreak{}.\allowbreak{}m\allowbreak{}e\allowbreak{}r\allowbreak{}g\allowbreak{}e\allowbreak{}s\allowbreak{}.\allowbreak{}i\allowbreak{}t\allowbreak{}e\allowbreak{}m\allowbreak{}s\allowbreak{}(\allowbreak{})\allowbreak{}:} & 按合并表顺序遍历两个子 \BpeDocCodeInline{token} 与新编号。 \\[3pt]
\SourceLineNumber{70} & \BpeDocCodeCell{\char32\relax\allowbreak{}\char32\relax\allowbreak{}\char32\relax\allowbreak{}\char32\relax\allowbreak{}\char32\relax\allowbreak{}\char32\relax\allowbreak{}\char32\relax\allowbreak{}\char32\relax\allowbreak{}\char32\relax\allowbreak{}\char32\relax\allowbreak{}\char32\relax\allowbreak{}\char32\relax\allowbreak{}v\allowbreak{}o\allowbreak{}c\allowbreak{}a\allowbreak{}b\allowbreak{}[\allowbreak{}i\allowbreak{}d\allowbreak{}x\allowbreak{}]\allowbreak{}\char32\relax\allowbreak{}=\allowbreak{}\char32\relax\allowbreak{}v\allowbreak{}o\allowbreak{}c\allowbreak{}a\allowbreak{}b\allowbreak{}[\allowbreak{}p\allowbreak{}0\allowbreak{}]\allowbreak{}\char32\relax\allowbreak{}+\allowbreak{}\char32\relax\allowbreak{}v\allowbreak{}o\allowbreak{}c\allowbreak{}a\allowbreak{}b\allowbreak{}[\allowbreak{}p\allowbreak{}1\allowbreak{}]} & 把两个子项的内部字节表示拼接，建立新 \BpeDocCodeInline{token} 项。 \\[3pt]
\SourceLineNumber{71} & \BpeDocCodeCell{\char32\relax\allowbreak{}\char32\relax\allowbreak{}\char32\relax\allowbreak{}\char32\relax\allowbreak{}\char32\relax\allowbreak{}\char32\relax\allowbreak{}\char32\relax\allowbreak{}\char32\relax\allowbreak{}s\allowbreak{}e\allowbreak{}l\allowbreak{}f\allowbreak{}.\allowbreak{}v\allowbreak{}o\allowbreak{}c\allowbreak{}a\allowbreak{}b\allowbreak{}\char32\relax\allowbreak{}=\allowbreak{}\char32\relax\allowbreak{}v\allowbreak{}o\allowbreak{}c\allowbreak{}a\allowbreak{}b} & 把构造好的内部词表存到对象。 \\[3pt]
\SourceLineNumber{76} & \BpeDocCodeCell{\char32\relax\allowbreak{}\char32\relax\allowbreak{}\char32\relax\allowbreak{}\char32\relax\allowbreak{}\char32\relax\allowbreak{}\char32\relax\allowbreak{}\char32\relax\allowbreak{}\char32\relax\allowbreak{}s\allowbreak{}e\allowbreak{}l\allowbreak{}f\allowbreak{}.\allowbreak{}b\allowbreak{}y\allowbreak{}t\allowbreak{}e\allowbreak{}\char95\relax\allowbreak{}s\allowbreak{}h\allowbreak{}u\allowbreak{}f\allowbreak{}f\allowbreak{}l\allowbreak{}e\allowbreak{}\char32\relax\allowbreak{}=\allowbreak{}\char32\relax\allowbreak{}\char123\relax\allowbreak{}i\allowbreak{}:\allowbreak{}\char32\relax\allowbreak{}m\allowbreak{}e\allowbreak{}r\allowbreak{}g\allowbreak{}e\allowbreak{}a\allowbreak{}b\allowbreak{}l\allowbreak{}e\allowbreak{}\char95\relax\allowbreak{}r\allowbreak{}a\allowbreak{}n\allowbreak{}k\allowbreak{}s\allowbreak{}[\allowbreak{}b\allowbreak{}y\allowbreak{}t\allowbreak{}e\allowbreak{}s\allowbreak{}(\allowbreak{}[\allowbreak{}i\allowbreak{}]\allowbreak{})\allowbreak{}]\allowbreak{}\char32\relax\allowbreak{}f\allowbreak{}o\allowbreak{}r\allowbreak{}\char32\relax\allowbreak{}i\allowbreak{}\char32\relax\allowbreak{}i\allowbreak{}n\allowbreak{}\char32\relax\allowbreak{}r\allowbreak{}a\allowbreak{}n\allowbreak{}g\allowbreak{}e\allowbreak{}(\allowbreak{}2\allowbreak{}5\allowbreak{}6\allowbreak{})\allowbreak{}\char125\relax} & 为每个原始字节 i 查官方单字节 \BpeDocCodeInline{token} 编号，建立正向重排表。 \\[3pt]
\SourceLineNumber{77} & \BpeDocCodeCell{\char32\relax\allowbreak{}\char32\relax\allowbreak{}\char32\relax\allowbreak{}\char32\relax\allowbreak{}\char32\relax\allowbreak{}\char32\relax\allowbreak{}\char32\relax\allowbreak{}\char32\relax\allowbreak{}s\allowbreak{}e\allowbreak{}l\allowbreak{}f\allowbreak{}.\allowbreak{}i\allowbreak{}n\allowbreak{}v\allowbreak{}e\allowbreak{}r\allowbreak{}s\allowbreak{}e\allowbreak{}\char95\relax\allowbreak{}b\allowbreak{}y\allowbreak{}t\allowbreak{}e\allowbreak{}\char95\relax\allowbreak{}s\allowbreak{}h\allowbreak{}u\allowbreak{}f\allowbreak{}f\allowbreak{}l\allowbreak{}e\allowbreak{}\char32\relax\allowbreak{}=\allowbreak{}\char32\relax\allowbreak{}\char123\relax\allowbreak{}v\allowbreak{}:\allowbreak{}\char32\relax\allowbreak{}k\allowbreak{}\char32\relax\allowbreak{}f\allowbreak{}o\allowbreak{}r\allowbreak{}\char32\relax\allowbreak{}k\allowbreak{},\allowbreak{}\char32\relax\allowbreak{}v\allowbreak{}\char32\relax\allowbreak{}i\allowbreak{}n\allowbreak{}\char32\relax\allowbreak{}s\allowbreak{}e\allowbreak{}l\allowbreak{}f\allowbreak{}.\allowbreak{}b\allowbreak{}y\allowbreak{}t\allowbreak{}e\allowbreak{}\char95\relax\allowbreak{}s\allowbreak{}h\allowbreak{}u\allowbreak{}f\allowbreak{}f\allowbreak{}l\allowbreak{}e\allowbreak{}.\allowbreak{}i\allowbreak{}t\allowbreak{}e\allowbreak{}m\allowbreak{}s\allowbreak{}(\allowbreak{})\allowbreak{}\char125\relax} & 交换键和值，建立“官方单字节编号 \(\rightarrow\) 原始字节”的逆向映射。 \\[3pt]
\SourceLineNumber{79} & \BpeDocCodeCell{\char32\relax\allowbreak{}\char32\relax\allowbreak{}\char32\relax\allowbreak{}\char32\relax\allowbreak{}\char32\relax\allowbreak{}\char32\relax\allowbreak{}\char32\relax\allowbreak{}\char32\relax\allowbreak{}s\allowbreak{}e\allowbreak{}l\allowbreak{}f\allowbreak{}.\allowbreak{}r\allowbreak{}e\allowbreak{}g\allowbreak{}i\allowbreak{}s\allowbreak{}t\allowbreak{}e\allowbreak{}r\allowbreak{}\char95\relax\allowbreak{}s\allowbreak{}p\allowbreak{}e\allowbreak{}c\allowbreak{}i\allowbreak{}a\allowbreak{}l\allowbreak{}\char95\relax\allowbreak{}t\allowbreak{}o\allowbreak{}k\allowbreak{}e\allowbreak{}n\allowbreak{}s\allowbreak{}(\allowbreak{}G\allowbreak{}P\allowbreak{}T\allowbreak{}4\allowbreak{}\char95\relax\allowbreak{}S\allowbreak{}P\allowbreak{}E\allowbreak{}C\allowbreak{}I\allowbreak{}A\allowbreak{}L\allowbreak{}\char95\relax\allowbreak{}T\allowbreak{}O\allowbreak{}K\allowbreak{}E\allowbreak{}N\allowbreak{}S\allowbreak{})} & 调用父类提供的登记方法，保存 \BpeDocCodeInline{GPT4\char95\relax\allowbreak{}SPECIAL\char95\relax\allowbreak{}TOKENS} 及其反向表。 \\[3pt]
\end{longtable}
\endgroup


初学时可以把 \BpeDocCodeInline{byte\char95\relax\allowbreak{}shuffle} 理解为“两套编号体系之间的换号表”。同一个原始字节，在 \BpeDocCodeInline{BasicTokenizer} 中初始编号直接取它的数值，在 \BpeDocCodeInline{cl100k\char95\relax\allowbreak{}base} 中则要查官方编号。重排并没有改变输入文字的含义；它只是让后续合并表看到自己预期的编号。构建出来的 \BpeDocCodeInline{self.\allowbreak{}vocab} 暂存的是这套内部表示，因此解码时还必须反向换回来。


\Needspace{18\baselineskip}
\BpeDocSubsubsection{\BpeDocCodeInline{gpt4.\allowbreak{}py}：编码前先换号，81–85 行}{gpt4.py：编码前先换号，81–85 行}


\Needspace{4\baselineskip}
\begingroup\fontsize{8.9}{12.2}\selectfont\setstretch{1}
\setlength{\BpeDocTableWidth}{\dimexpr\linewidth-6\tabcolsep\relax}
\renewcommand{\arraystretch}{1.08}
\par\addvspace{0.45em}
\noindent\begin{minipage}{\linewidth}
{\small \BpeDocCodeInline{minbpe/\allowbreak{}gpt4.\allowbreak{}py}，第 81–85 行（仅列执行语句）}\par\vspace{0.35em}
\begin{tabular}{>{\raggedleft\arraybackslash}p{0.065\BpeDocTableWidth}>{\raggedright\arraybackslash}p{0.475\BpeDocTableWidth}>{\raggedright\arraybackslash}p{0.460\BpeDocTableWidth}}
\toprule
\textbf{行号} & \textbf{原代码} & \textbf{逐行解释} \\
\midrule
\SourceLineNumber{81} & \BpeDocCodeCell{\char32\relax\allowbreak{}\char32\relax\allowbreak{}\char32\relax\allowbreak{}\char32\relax\allowbreak{}d\allowbreak{}e\allowbreak{}f\allowbreak{}\char32\relax\allowbreak{}\char95\relax\allowbreak{}e\allowbreak{}n\allowbreak{}c\allowbreak{}o\allowbreak{}d\allowbreak{}e\allowbreak{}\char95\relax\allowbreak{}c\allowbreak{}h\allowbreak{}u\allowbreak{}n\allowbreak{}k\allowbreak{}(\allowbreak{}s\allowbreak{}e\allowbreak{}l\allowbreak{}f\allowbreak{},\allowbreak{}\char32\relax\allowbreak{}t\allowbreak{}e\allowbreak{}x\allowbreak{}t\allowbreak{}\char95\relax\allowbreak{}b\allowbreak{}y\allowbreak{}t\allowbreak{}e\allowbreak{}s\allowbreak{})\allowbreak{}:} & 改写父类的单块编码方法，输入仍是一个 UTF-8 字节串。 \\[3pt]
\SourceLineNumber{83} & \BpeDocCodeCell{\char32\relax\allowbreak{}\char32\relax\allowbreak{}\char32\relax\allowbreak{}\char32\relax\allowbreak{}\char32\relax\allowbreak{}\char32\relax\allowbreak{}\char32\relax\allowbreak{}\char32\relax\allowbreak{}t\allowbreak{}e\allowbreak{}x\allowbreak{}t\allowbreak{}\char95\relax\allowbreak{}b\allowbreak{}y\allowbreak{}t\allowbreak{}e\allowbreak{}s\allowbreak{}\char32\relax\allowbreak{}=\allowbreak{}\char32\relax\allowbreak{}b\allowbreak{}y\allowbreak{}t\allowbreak{}e\allowbreak{}s\allowbreak{}(\allowbreak{}s\allowbreak{}e\allowbreak{}l\allowbreak{}f\allowbreak{}.\allowbreak{}b\allowbreak{}y\allowbreak{}t\allowbreak{}e\allowbreak{}\char95\relax\allowbreak{}s\allowbreak{}h\allowbreak{}u\allowbreak{}f\allowbreak{}f\allowbreak{}l\allowbreak{}e\allowbreak{}[\allowbreak{}b\allowbreak{}]\allowbreak{}\char32\relax\allowbreak{}f\allowbreak{}o\allowbreak{}r\allowbreak{}\char32\relax\allowbreak{}b\allowbreak{}\char32\relax\allowbreak{}i\allowbreak{}n\allowbreak{}\char32\relax\allowbreak{}t\allowbreak{}e\allowbreak{}x\allowbreak{}t\allowbreak{}\char95\relax\allowbreak{}b\allowbreak{}y\allowbreak{}t\allowbreak{}e\allowbreak{}s\allowbreak{})} & 逐字节查询 \BpeDocCodeInline{byte\char95\relax\allowbreak{}shuffle}，得到官方单字节编号，并用 \BpeDocCodeInline{bytes} 容器包装这些 0–255 的数。 \\[3pt]
\SourceLineNumber{84} & \BpeDocCodeCell{\char32\relax\allowbreak{}\char32\relax\allowbreak{}\char32\relax\allowbreak{}\char32\relax\allowbreak{}\char32\relax\allowbreak{}\char32\relax\allowbreak{}\char32\relax\allowbreak{}\char32\relax\allowbreak{}i\allowbreak{}d\allowbreak{}s\allowbreak{}\char32\relax\allowbreak{}=\allowbreak{}\char32\relax\allowbreak{}s\allowbreak{}u\allowbreak{}p\allowbreak{}e\allowbreak{}r\allowbreak{}(\allowbreak{})\allowbreak{}.\allowbreak{}\char95\relax\allowbreak{}e\allowbreak{}n\allowbreak{}c\allowbreak{}o\allowbreak{}d\allowbreak{}e\allowbreak{}\char95\relax\allowbreak{}c\allowbreak{}h\allowbreak{}u\allowbreak{}n\allowbreak{}k\allowbreak{}(\allowbreak{}t\allowbreak{}e\allowbreak{}x\allowbreak{}t\allowbreak{}\char95\relax\allowbreak{}b\allowbreak{}y\allowbreak{}t\allowbreak{}e\allowbreak{}s\allowbreak{})} & 调用 \BpeDocCodeInline{RegexTokenizer.\allowbreak{}\char95\relax\allowbreak{}encode\char95\relax\allowbreak{}chunk} 执行已学规则的 BPE。\BpeDocCodeInline{super}() 避免再次调用自己而无限递归。 \\[3pt]
\SourceLineNumber{85} & \BpeDocCodeCell{\char32\relax\allowbreak{}\char32\relax\allowbreak{}\char32\relax\allowbreak{}\char32\relax\allowbreak{}\char32\relax\allowbreak{}\char32\relax\allowbreak{}\char32\relax\allowbreak{}\char32\relax\allowbreak{}r\allowbreak{}e\allowbreak{}t\allowbreak{}u\allowbreak{}r\allowbreak{}n\allowbreak{}\char32\relax\allowbreak{}i\allowbreak{}d\allowbreak{}s} & 返回该块的最终编号列表。 \\[3pt]
\bottomrule
\end{tabular}
\end{minipage}\par\addvspace{0.65em}
\endgroup


为什么不用在 \BpeDocCodeInline{encode\char95\relax\allowbreak{}ordinary} 里再写一遍 GPT-4 专用流程？因为它调用的是 \BpeDocCodeInline{self.\allowbreak{}\char95\relax\allowbreak{}encode\char95\relax\allowbreak{}chunk}。当 \BpeDocCodeInline{self} 是 \BpeDocCodeInline{GPT4Tokenizer} 对象时，Python 会使用这里改写的版本。于是普通的正则分块流程得到复用，而特殊的换号步骤只写在这一小段里。


\Needspace{18\baselineskip}
\BpeDocSubsubsection{\BpeDocCodeInline{gpt4.\allowbreak{}py}：解码时换回原始字节，87–92 行}{gpt4.py：解码时换回原始字节，87–92 行}


\Needspace{4\baselineskip}
\begingroup\fontsize{8.9}{12.2}\selectfont\setstretch{1}
\setlength{\BpeDocTableWidth}{\dimexpr\linewidth-6\tabcolsep\relax}
\renewcommand{\arraystretch}{1.08}
\par\addvspace{0.45em}
\noindent\begin{minipage}{\linewidth}
{\small \BpeDocCodeInline{minbpe/\allowbreak{}gpt4.\allowbreak{}py}，第 87–92 行（仅列执行语句）}\par\vspace{0.35em}
\begin{tabular}{>{\raggedleft\arraybackslash}p{0.065\BpeDocTableWidth}>{\raggedright\arraybackslash}p{0.475\BpeDocTableWidth}>{\raggedright\arraybackslash}p{0.460\BpeDocTableWidth}}
\toprule
\textbf{行号} & \textbf{原代码} & \textbf{逐行解释} \\
\midrule
\SourceLineNumber{87} & \BpeDocCodeCell{\char32\relax\allowbreak{}\char32\relax\allowbreak{}\char32\relax\allowbreak{}\char32\relax\allowbreak{}d\allowbreak{}e\allowbreak{}f\allowbreak{}\char32\relax\allowbreak{}d\allowbreak{}e\allowbreak{}c\allowbreak{}o\allowbreak{}d\allowbreak{}e\allowbreak{}(\allowbreak{}s\allowbreak{}e\allowbreak{}l\allowbreak{}f\allowbreak{},\allowbreak{}\char32\relax\allowbreak{}i\allowbreak{}d\allowbreak{}s\allowbreak{})\allowbreak{}:} & 定义 \BpeDocCodeInline{GPT4Tokenizer} 自己的解码方法，覆盖父类版本。 \\[3pt]
\SourceLineNumber{89} & \BpeDocCodeCell{\char32\relax\allowbreak{}\char32\relax\allowbreak{}\char32\relax\allowbreak{}\char32\relax\allowbreak{}\char32\relax\allowbreak{}\char32\relax\allowbreak{}\char32\relax\allowbreak{}\char32\relax\allowbreak{}t\allowbreak{}e\allowbreak{}x\allowbreak{}t\allowbreak{}\char95\relax\allowbreak{}b\allowbreak{}y\allowbreak{}t\allowbreak{}e\allowbreak{}s\allowbreak{}\char32\relax\allowbreak{}=\allowbreak{}\char32\relax\allowbreak{}b\allowbreak{}"\allowbreak{}"\allowbreak{}.\allowbreak{}j\allowbreak{}o\allowbreak{}i\allowbreak{}n\allowbreak{}(\allowbreak{}s\allowbreak{}e\allowbreak{}l\allowbreak{}f\allowbreak{}.\allowbreak{}v\allowbreak{}o\allowbreak{}c\allowbreak{}a\allowbreak{}b\allowbreak{}[\allowbreak{}i\allowbreak{}d\allowbreak{}x\allowbreak{}]\allowbreak{}\char32\relax\allowbreak{}f\allowbreak{}o\allowbreak{}r\allowbreak{}\char32\relax\allowbreak{}i\allowbreak{}d\allowbreak{}x\allowbreak{}\char32\relax\allowbreak{}i\allowbreak{}n\allowbreak{}\char32\relax\allowbreak{}i\allowbreak{}d\allowbreak{}s\allowbreak{})} & 按普通词表查各编号，并把内部字节表示拼接起来。这里没有特殊 \BpeDocCodeInline{token} 分支。 \\[3pt]
\SourceLineNumber{90} & \BpeDocCodeCell{\char32\relax\allowbreak{}\char32\relax\allowbreak{}\char32\relax\allowbreak{}\char32\relax\allowbreak{}\char32\relax\allowbreak{}\char32\relax\allowbreak{}\char32\relax\allowbreak{}\char32\relax\allowbreak{}t\allowbreak{}e\allowbreak{}x\allowbreak{}t\allowbreak{}\char95\relax\allowbreak{}b\allowbreak{}y\allowbreak{}t\allowbreak{}e\allowbreak{}s\allowbreak{}\char32\relax\allowbreak{}=\allowbreak{}\char32\relax\allowbreak{}b\allowbreak{}y\allowbreak{}t\allowbreak{}e\allowbreak{}s\allowbreak{}(\allowbreak{}s\allowbreak{}e\allowbreak{}l\allowbreak{}f\allowbreak{}.\allowbreak{}i\allowbreak{}n\allowbreak{}v\allowbreak{}e\allowbreak{}r\allowbreak{}s\allowbreak{}e\allowbreak{}\char95\relax\allowbreak{}b\allowbreak{}y\allowbreak{}t\allowbreak{}e\allowbreak{}\char95\relax\allowbreak{}s\allowbreak{}h\allowbreak{}u\allowbreak{}f\allowbreak{}f\allowbreak{}l\allowbreak{}e\allowbreak{}[\allowbreak{}b\allowbreak{}]\allowbreak{}\char32\relax\allowbreak{}f\allowbreak{}o\allowbreak{}r\allowbreak{}\char32\relax\allowbreak{}b\allowbreak{}\char32\relax\allowbreak{}i\allowbreak{}n\allowbreak{}\char32\relax\allowbreak{}t\allowbreak{}e\allowbreak{}x\allowbreak{}t\allowbreak{}\char95\relax\allowbreak{}b\allowbreak{}y\allowbreak{}t\allowbreak{}e\allowbreak{}s\allowbreak{})} & 逐个查询 \BpeDocCodeInline{inverse\char95\relax\allowbreak{}byte\char95\relax\allowbreak{}shuffle}，将内部编号表示变回真实原始字节。 \\[3pt]
\SourceLineNumber{91} & \BpeDocCodeCell{\char32\relax\allowbreak{}\char32\relax\allowbreak{}\char32\relax\allowbreak{}\char32\relax\allowbreak{}\char32\relax\allowbreak{}\char32\relax\allowbreak{}\char32\relax\allowbreak{}\char32\relax\allowbreak{}t\allowbreak{}e\allowbreak{}x\allowbreak{}t\allowbreak{}\char32\relax\allowbreak{}=\allowbreak{}\char32\relax\allowbreak{}t\allowbreak{}e\allowbreak{}x\allowbreak{}t\allowbreak{}\char95\relax\allowbreak{}b\allowbreak{}y\allowbreak{}t\allowbreak{}e\allowbreak{}s\allowbreak{}.\allowbreak{}d\allowbreak{}e\allowbreak{}c\allowbreak{}o\allowbreak{}d\allowbreak{}e\allowbreak{}(\allowbreak{}"\allowbreak{}u\allowbreak{}t\allowbreak{}f\allowbreak{}-\allowbreak{}8\allowbreak{}"\allowbreak{},\allowbreak{}\char32\relax\allowbreak{}e\allowbreak{}r\allowbreak{}r\allowbreak{}o\allowbreak{}r\allowbreak{}s\allowbreak{}=\allowbreak{}"\allowbreak{}r\allowbreak{}e\allowbreak{}p\allowbreak{}l\allowbreak{}a\allowbreak{}c\allowbreak{}e\allowbreak{}"\allowbreak{})} & 按 UTF-8 解码为字符串；无效字节用替代字符处理。 \\[3pt]
\SourceLineNumber{92} & \BpeDocCodeCell{\char32\relax\allowbreak{}\char32\relax\allowbreak{}\char32\relax\allowbreak{}\char32\relax\allowbreak{}\char32\relax\allowbreak{}\char32\relax\allowbreak{}\char32\relax\allowbreak{}\char32\relax\allowbreak{}r\allowbreak{}e\allowbreak{}t\allowbreak{}u\allowbreak{}r\allowbreak{}n\allowbreak{}\char32\relax\allowbreak{}t\allowbreak{}e\allowbreak{}x\allowbreak{}t} & 返回文本。 \\[3pt]
\bottomrule
\end{tabular}
\end{minipage}\par\addvspace{0.65em}
\endgroup


这个固定提交有一个应如实写明的限制：\BpeDocCodeInline{GPT4Tokenizer.\allowbreak{}decode} 直接读取 \BpeDocCodeInline{self.\allowbreak{}vocab}，没有像 \BpeDocCodeInline{RegexTokenizer.\allowbreak{}decode} 那样处理 \BpeDocCodeInline{inverse\char95\relax\allowbreak{}special\char95\relax\allowbreak{}tokens}。其初始化也没有把特殊编号放进普通 \BpeDocCodeInline{self.\allowbreak{}vocab}。因此若将 \BpeDocCodeInline{encode}(..., \BpeDocCodeInline{allowed\char95\relax\allowbreak{}special}="\BpeDocCodeInline{all}") 得到的特殊 ID 交给这个 \BpeDocCodeInline{decode}，会在查表处遇到 \BpeDocCodeInline{KeyError}。普通文本的往返编码和包含特殊 ID 的往返编码，应分开说明；不能声称该版本在所有特殊 \BpeDocCodeInline{token} 情形下都完整等价于 tiktoken。


\Needspace{4\baselineskip}
\begingroup\fontsize{8.9}{12.2}\selectfont\setstretch{1}
\setlength{\BpeDocTableWidth}{\dimexpr\linewidth-6\tabcolsep\relax}
\renewcommand{\arraystretch}{1.08}
\par\addvspace{0.45em}
\noindent\begin{minipage}{\linewidth}
{\small \BpeDocCodeInline{minbpe/\allowbreak{}gpt4.\allowbreak{}py}，第 94–96 行（仅列执行语句）}\par\vspace{0.35em}
\begin{tabular}{>{\raggedleft\arraybackslash}p{0.065\BpeDocTableWidth}>{\raggedright\arraybackslash}p{0.475\BpeDocTableWidth}>{\raggedright\arraybackslash}p{0.460\BpeDocTableWidth}}
\toprule
\textbf{行号} & \textbf{原代码} & \textbf{逐行解释} \\
\midrule
\SourceLineNumber{95} & \BpeDocCodeCell{\char32\relax\allowbreak{}\char32\relax\allowbreak{}\char32\relax\allowbreak{}\char32\relax\allowbreak{}d\allowbreak{}e\allowbreak{}f\allowbreak{}\char32\relax\allowbreak{}t\allowbreak{}r\allowbreak{}a\allowbreak{}i\allowbreak{}n\allowbreak{}(\allowbreak{}s\allowbreak{}e\allowbreak{}l\allowbreak{}f\allowbreak{},\allowbreak{}\char32\relax\allowbreak{}t\allowbreak{}e\allowbreak{}x\allowbreak{}t\allowbreak{},\allowbreak{}\char32\relax\allowbreak{}v\allowbreak{}o\allowbreak{}c\allowbreak{}a\allowbreak{}b\allowbreak{}\char95\relax\allowbreak{}s\allowbreak{}i\allowbreak{}z\allowbreak{}e\allowbreak{},\allowbreak{}\char32\relax\allowbreak{}v\allowbreak{}e\allowbreak{}r\allowbreak{}b\allowbreak{}o\allowbreak{}s\allowbreak{}e\allowbreak{}=\allowbreak{}F\allowbreak{}a\allowbreak{}l\allowbreak{}s\allowbreak{}e\allowbreak{})\allowbreak{}:} & 显式定义 \BpeDocCodeInline{train} 方法以覆盖继承来的训练接口。 \\[3pt]
\SourceLineNumber{96} & \BpeDocCodeCell{\char32\relax\allowbreak{}\char32\relax\allowbreak{}\char32\relax\allowbreak{}\char32\relax\allowbreak{}\char32\relax\allowbreak{}\char32\relax\allowbreak{}\char32\relax\allowbreak{}\char32\relax\allowbreak{}r\allowbreak{}a\allowbreak{}i\allowbreak{}s\allowbreak{}e\allowbreak{}\char32\relax\allowbreak{}N\allowbreak{}o\allowbreak{}t\allowbreak{}I\allowbreak{}m\allowbreak{}p\allowbreak{}l\allowbreak{}e\allowbreak{}m\allowbreak{}e\allowbreak{}n\allowbreak{}t\allowbreak{}e\allowbreak{}d\allowbreak{}E\allowbreak{}r\allowbreak{}r\allowbreak{}o\allowbreak{}r} & 调用时直接抛出 \BpeDocCodeInline{NotImplementedError}，不会重新学习词表。 \\[3pt]
\bottomrule
\end{tabular}
\end{minipage}\par\addvspace{0.65em}
\endgroup


\Needspace{9\baselineskip}
\BpeDocSubsection{进阶附录：\BpeDocCodeInline{recover\char95\relax\allowbreak{}merges} 如何从现成词表找回合并关系}{进阶附录：recover\_merges 如何从现成词表找回合并关系}


前面已足以解释使用流程。下面两段用于回答更深入的“为什么 \BpeDocCodeInline{GPT4Tokenizer} 的初始化可以得到 \BpeDocCodeInline{merges}”。tiktoken 给出的是“一个完整字节串对应哪个 \BpeDocCodeInline{rank}”，MinBPE 想要的是“哪两个旧 \BpeDocCodeInline{token} 合成这个 \BpeDocCodeInline{token}”。恢复函数会对每个多字节 \BpeDocCodeInline{token}，只应用比它更早的规则，把它还原到即将完成最后一次合并的两个子块。这不是用新训练语料重新统计频次。源码注释把它称为一个小型 BPE 过程，报告应把其实际含义说清楚。


\Needspace{4\baselineskip}
\begingroup\fontsize{8.9}{12.2}\selectfont\setstretch{1}
\setlength{\BpeDocTableWidth}{\dimexpr\linewidth-6\tabcolsep\relax}
\renewcommand{\arraystretch}{1.08}
{\small \BpeDocCodeInline{minbpe/\allowbreak{}gpt4.\allowbreak{}py}，第 29–46 行（仅列执行语句）}

\begin{longtable}{>{\raggedleft\arraybackslash}p{0.065\BpeDocTableWidth}>{\raggedright\arraybackslash}p{0.475\BpeDocTableWidth}>{\raggedright\arraybackslash}p{0.460\BpeDocTableWidth}}
\toprule
\textbf{行号} & \textbf{原代码} & \textbf{逐行解释} \\
\midrule\endfirsthead
\toprule
\textbf{行号} & \textbf{原代码} & \textbf{逐行解释} \\
\midrule\endhead
\midrule\endfoot
\bottomrule\endlastfoot
\SourceLineNumber{29} & \BpeDocCodeCell{d\allowbreak{}e\allowbreak{}f\allowbreak{}\char32\relax\allowbreak{}r\allowbreak{}e\allowbreak{}c\allowbreak{}o\allowbreak{}v\allowbreak{}e\allowbreak{}r\allowbreak{}\char95\relax\allowbreak{}m\allowbreak{}e\allowbreak{}r\allowbreak{}g\allowbreak{}e\allowbreak{}s\allowbreak{}(\allowbreak{}m\allowbreak{}e\allowbreak{}r\allowbreak{}g\allowbreak{}e\allowbreak{}a\allowbreak{}b\allowbreak{}l\allowbreak{}e\allowbreak{}\char95\relax\allowbreak{}r\allowbreak{}a\allowbreak{}n\allowbreak{}k\allowbreak{}s\allowbreak{})\allowbreak{}:} & 定义恢复函数，输入为完整字节串到 \BpeDocCodeInline{rank} 的已有映射。 \\[3pt]
\SourceLineNumber{35} & \BpeDocCodeCell{\char32\relax\allowbreak{}\char32\relax\allowbreak{}\char32\relax\allowbreak{}\char32\relax\allowbreak{}m\allowbreak{}e\allowbreak{}r\allowbreak{}g\allowbreak{}e\allowbreak{}s\allowbreak{}\char32\relax\allowbreak{}=\allowbreak{}\char32\relax\allowbreak{}\char123\relax\allowbreak{}\char125\relax} & 准备保存恢复出的“两个编号 \(\rightarrow\) 新编号”映射。 \\[3pt]
\SourceLineNumber{36} & \BpeDocCodeCell{\char32\relax\allowbreak{}\char32\relax\allowbreak{}\char32\relax\allowbreak{}\char32\relax\allowbreak{}f\allowbreak{}o\allowbreak{}r\allowbreak{}\char32\relax\allowbreak{}t\allowbreak{}o\allowbreak{}k\allowbreak{}e\allowbreak{}n\allowbreak{},\allowbreak{}\char32\relax\allowbreak{}r\allowbreak{}a\allowbreak{}n\allowbreak{}k\allowbreak{}\char32\relax\allowbreak{}i\allowbreak{}n\allowbreak{}\char32\relax\allowbreak{}m\allowbreak{}e\allowbreak{}r\allowbreak{}g\allowbreak{}e\allowbreak{}a\allowbreak{}b\allowbreak{}l\allowbreak{}e\allowbreak{}\char95\relax\allowbreak{}r\allowbreak{}a\allowbreak{}n\allowbreak{}k\allowbreak{}s\allowbreak{}.\allowbreak{}i\allowbreak{}t\allowbreak{}e\allowbreak{}m\allowbreak{}s\allowbreak{}(\allowbreak{})\allowbreak{}:} & 遍历已有映射中的字节串 \BpeDocCodeInline{token} 及其 \BpeDocCodeInline{rank}。 \\[3pt]
\SourceLineNumber{37} & \BpeDocCodeCell{\char32\relax\allowbreak{}\char32\relax\allowbreak{}\char32\relax\allowbreak{}\char32\relax\allowbreak{}\char32\relax\allowbreak{}\char32\relax\allowbreak{}\char32\relax\allowbreak{}\char32\relax\allowbreak{}i\allowbreak{}f\allowbreak{}\char32\relax\allowbreak{}l\allowbreak{}e\allowbreak{}n\allowbreak{}(\allowbreak{}t\allowbreak{}o\allowbreak{}k\allowbreak{}e\allowbreak{}n\allowbreak{})\allowbreak{}\char32\relax\allowbreak{}=\allowbreak{}=\allowbreak{}\char32\relax\allowbreak{}1\allowbreak{}:} & 判断该 \BpeDocCodeInline{token} 是否只包含一个字节。 \\[3pt]
\SourceLineNumber{38} & \BpeDocCodeCell{\char32\relax\allowbreak{}\char32\relax\allowbreak{}\char32\relax\allowbreak{}\char32\relax\allowbreak{}\char32\relax\allowbreak{}\char32\relax\allowbreak{}\char32\relax\allowbreak{}\char32\relax\allowbreak{}\char32\relax\allowbreak{}\char32\relax\allowbreak{}\char32\relax\allowbreak{}\char32\relax\allowbreak{}c\allowbreak{}o\allowbreak{}n\allowbreak{}t\allowbreak{}i\allowbreak{}n\allowbreak{}u\allowbreak{}e\allowbreak{}\char32\relax\allowbreak{}\char35\relax\allowbreak{}\char32\relax\allowbreak{}s\allowbreak{}k\allowbreak{}i\allowbreak{}p\allowbreak{}\char32\relax\allowbreak{}r\allowbreak{}a\allowbreak{}w\allowbreak{}\char32\relax\allowbreak{}b\allowbreak{}y\allowbreak{}t\allowbreak{}e\allowbreak{}s} & 单字节是起点，没有两个子 \BpeDocCodeInline{token} 的合并关系，直接跳过。 \\[3pt]
\SourceLineNumber{39} & \BpeDocCodeCell{\char32\relax\allowbreak{}\char32\relax\allowbreak{}\char32\relax\allowbreak{}\char32\relax\allowbreak{}\char32\relax\allowbreak{}\char32\relax\allowbreak{}\char32\relax\allowbreak{}\char32\relax\allowbreak{}p\allowbreak{}a\allowbreak{}i\allowbreak{}r\allowbreak{}\char32\relax\allowbreak{}=\allowbreak{}\char32\relax\allowbreak{}t\allowbreak{}u\allowbreak{}p\allowbreak{}l\allowbreak{}e\allowbreak{}(\allowbreak{}b\allowbreak{}p\allowbreak{}e\allowbreak{}(\allowbreak{}m\allowbreak{}e\allowbreak{}r\allowbreak{}g\allowbreak{}e\allowbreak{}a\allowbreak{}b\allowbreak{}l\allowbreak{}e\allowbreak{}\char95\relax\allowbreak{}r\allowbreak{}a\allowbreak{}n\allowbreak{}k\allowbreak{}s\allowbreak{},\allowbreak{}\char32\relax\allowbreak{}t\allowbreak{}o\allowbreak{}k\allowbreak{}e\allowbreak{}n\allowbreak{},\allowbreak{}\char32\relax\allowbreak{}m\allowbreak{}a\allowbreak{}x\allowbreak{}\char95\relax\allowbreak{}r\allowbreak{}a\allowbreak{}n\allowbreak{}k\allowbreak{}=\allowbreak{}r\allowbreak{}a\allowbreak{}n\allowbreak{}k\allowbreak{})\allowbreak{})} & 调用辅助 \BpeDocCodeInline{bpe}，只允许使用 \BpeDocCodeInline{rank} 小于当前 \BpeDocCodeInline{token} 的规则；把返回列表转成元组。 \\[3pt]
\SourceLineNumber{40} & \BpeDocCodeCell{\char32\relax\allowbreak{}\char32\relax\allowbreak{}\char32\relax\allowbreak{}\char32\relax\allowbreak{}\char32\relax\allowbreak{}\char32\relax\allowbreak{}\char32\relax\allowbreak{}\char32\relax\allowbreak{}a\allowbreak{}s\allowbreak{}s\allowbreak{}e\allowbreak{}r\allowbreak{}t\allowbreak{}\char32\relax\allowbreak{}l\allowbreak{}e\allowbreak{}n\allowbreak{}(\allowbreak{}p\allowbreak{}a\allowbreak{}i\allowbreak{}r\allowbreak{})\allowbreak{}\char32\relax\allowbreak{}=\allowbreak{}=\allowbreak{}\char32\relax\allowbreak{}2} & 正常情况下应恰好剩两个子块；否则断言失败。 \\[3pt]
\SourceLineNumber{42} & \BpeDocCodeCell{\char32\relax\allowbreak{}\char32\relax\allowbreak{}\char32\relax\allowbreak{}\char32\relax\allowbreak{}\char32\relax\allowbreak{}\char32\relax\allowbreak{}\char32\relax\allowbreak{}\char32\relax\allowbreak{}i\allowbreak{}x\allowbreak{}0\allowbreak{}\char32\relax\allowbreak{}=\allowbreak{}\char32\relax\allowbreak{}m\allowbreak{}e\allowbreak{}r\allowbreak{}g\allowbreak{}e\allowbreak{}a\allowbreak{}b\allowbreak{}l\allowbreak{}e\allowbreak{}\char95\relax\allowbreak{}r\allowbreak{}a\allowbreak{}n\allowbreak{}k\allowbreak{}s\allowbreak{}[\allowbreak{}p\allowbreak{}a\allowbreak{}i\allowbreak{}r\allowbreak{}[\allowbreak{}0\allowbreak{}]\allowbreak{}]} & 查第一个子块的 \BpeDocCodeInline{rank}／编号。 \\[3pt]
\SourceLineNumber{43} & \BpeDocCodeCell{\char32\relax\allowbreak{}\char32\relax\allowbreak{}\char32\relax\allowbreak{}\char32\relax\allowbreak{}\char32\relax\allowbreak{}\char32\relax\allowbreak{}\char32\relax\allowbreak{}\char32\relax\allowbreak{}i\allowbreak{}x\allowbreak{}1\allowbreak{}\char32\relax\allowbreak{}=\allowbreak{}\char32\relax\allowbreak{}m\allowbreak{}e\allowbreak{}r\allowbreak{}g\allowbreak{}e\allowbreak{}a\allowbreak{}b\allowbreak{}l\allowbreak{}e\allowbreak{}\char95\relax\allowbreak{}r\allowbreak{}a\allowbreak{}n\allowbreak{}k\allowbreak{}s\allowbreak{}[\allowbreak{}p\allowbreak{}a\allowbreak{}i\allowbreak{}r\allowbreak{}[\allowbreak{}1\allowbreak{}]\allowbreak{}]} & 查第二个子块的 \BpeDocCodeInline{rank}／编号。 \\[3pt]
\SourceLineNumber{44} & \BpeDocCodeCell{\char32\relax\allowbreak{}\char32\relax\allowbreak{}\char32\relax\allowbreak{}\char32\relax\allowbreak{}\char32\relax\allowbreak{}\char32\relax\allowbreak{}\char32\relax\allowbreak{}\char32\relax\allowbreak{}m\allowbreak{}e\allowbreak{}r\allowbreak{}g\allowbreak{}e\allowbreak{}s\allowbreak{}[\allowbreak{}(\allowbreak{}i\allowbreak{}x\allowbreak{}0\allowbreak{},\allowbreak{}\char32\relax\allowbreak{}i\allowbreak{}x\allowbreak{}1\allowbreak{})\allowbreak{}]\allowbreak{}\char32\relax\allowbreak{}=\allowbreak{}\char32\relax\allowbreak{}r\allowbreak{}a\allowbreak{}n\allowbreak{}k} & 记录这两个子编号合并后得到当前 \BpeDocCodeInline{rank}。 \\[3pt]
\SourceLineNumber{46} & \BpeDocCodeCell{\char32\relax\allowbreak{}\char32\relax\allowbreak{}\char32\relax\allowbreak{}\char32\relax\allowbreak{}r\allowbreak{}e\allowbreak{}t\allowbreak{}u\allowbreak{}r\allowbreak{}n\allowbreak{}\char32\relax\allowbreak{}m\allowbreak{}e\allowbreak{}r\allowbreak{}g\allowbreak{}e\allowbreak{}s} & 返回恢复完成的合并表。 \\[3pt]
\end{longtable}
\endgroup


例如完整 \BpeDocCodeInline{token} 表里有 ab、abab，且 ab 的规则早于 abab。恢复 abab 时，先允许 a+b\(\rightarrow\)ab，但禁止使用 abab 自己的规则，于是停在 [ab,ab]，便知道 abab 是由两个 ab 合成的。这是理解 \BpeDocCodeInline{max\char95\relax\allowbreak{}rank} 参数最直接的办法。这里为了直观使用文本片段表示，实际代码处理的是 \BpeDocCodeInline{bytes}。


\Needspace{9\baselineskip}
\BpeDocSubsubsection{进阶附录：辅助 \BpeDocCodeInline{bpe} 的逐行过程，11–26 行}{进阶附录：辅助 bpe 的逐行过程，11–26 行}


\Needspace{4\baselineskip}
\begingroup\fontsize{8.9}{12.2}\selectfont\setstretch{1}
\setlength{\BpeDocTableWidth}{\dimexpr\linewidth-6\tabcolsep\relax}
\renewcommand{\arraystretch}{1.08}
{\small \BpeDocCodeInline{minbpe/\allowbreak{}gpt4.\allowbreak{}py}，第 11–26 行（仅列执行语句）}

\begin{longtable}{>{\raggedleft\arraybackslash}p{0.065\BpeDocTableWidth}>{\raggedright\arraybackslash}p{0.475\BpeDocTableWidth}>{\raggedright\arraybackslash}p{0.460\BpeDocTableWidth}}
\toprule
\textbf{行号} & \textbf{原代码} & \textbf{逐行解释} \\
\midrule\endfirsthead
\toprule
\textbf{行号} & \textbf{原代码} & \textbf{逐行解释} \\
\midrule\endhead
\midrule\endfoot
\bottomrule\endlastfoot
\SourceLineNumber{11} & \BpeDocCodeCell{d\allowbreak{}e\allowbreak{}f\allowbreak{}\char32\relax\allowbreak{}b\allowbreak{}p\allowbreak{}e\allowbreak{}(\allowbreak{}m\allowbreak{}e\allowbreak{}r\allowbreak{}g\allowbreak{}e\allowbreak{}a\allowbreak{}b\allowbreak{}l\allowbreak{}e\allowbreak{}\char95\relax\allowbreak{}r\allowbreak{}a\allowbreak{}n\allowbreak{}k\allowbreak{}s\allowbreak{},\allowbreak{}\char32\relax\allowbreak{}t\allowbreak{}o\allowbreak{}k\allowbreak{}e\allowbreak{}n\allowbreak{},\allowbreak{}\char32\relax\allowbreak{}m\allowbreak{}a\allowbreak{}x\allowbreak{}\char95\relax\allowbreak{}r\allowbreak{}a\allowbreak{}n\allowbreak{}k\allowbreak{})\allowbreak{}:} & 定义辅助函数，输入已有 \BpeDocCodeInline{rank} 表、一个待拆解字节串，以及允许合并等级的上界。 \\[3pt]
\SourceLineNumber{13} & \BpeDocCodeCell{\char32\relax\allowbreak{}\char32\relax\allowbreak{}\char32\relax\allowbreak{}\char32\relax\allowbreak{}p\allowbreak{}a\allowbreak{}r\allowbreak{}t\allowbreak{}s\allowbreak{}\char32\relax\allowbreak{}=\allowbreak{}\char32\relax\allowbreak{}[\allowbreak{}b\allowbreak{}y\allowbreak{}t\allowbreak{}e\allowbreak{}s\allowbreak{}(\allowbreak{}[\allowbreak{}b\allowbreak{}]\allowbreak{})\allowbreak{}\char32\relax\allowbreak{}f\allowbreak{}o\allowbreak{}r\allowbreak{}\char32\relax\allowbreak{}b\allowbreak{}\char32\relax\allowbreak{}i\allowbreak{}n\allowbreak{}\char32\relax\allowbreak{}t\allowbreak{}o\allowbreak{}k\allowbreak{}e\allowbreak{}n\allowbreak{}]} & 先把 \BpeDocCodeInline{token} 拆成多个单字节串；此处 \BpeDocCodeInline{parts} 是 \BpeDocCodeInline{bytes} 片段的列表，不是整数 ID 列表。 \\[3pt]
\SourceLineNumber{14} & \BpeDocCodeCell{\char32\relax\allowbreak{}\char32\relax\allowbreak{}\char32\relax\allowbreak{}\char32\relax\allowbreak{}w\allowbreak{}h\allowbreak{}i\allowbreak{}l\allowbreak{}e\allowbreak{}\char32\relax\allowbreak{}T\allowbreak{}r\allowbreak{}u\allowbreak{}e\allowbreak{}:} & 进入循环，直到没有合适的合并可以执行。 \\[3pt]
\SourceLineNumber{15} & \BpeDocCodeCell{\char32\relax\allowbreak{}\char32\relax\allowbreak{}\char32\relax\allowbreak{}\char32\relax\allowbreak{}\char32\relax\allowbreak{}\char32\relax\allowbreak{}\char32\relax\allowbreak{}\char32\relax\allowbreak{}m\allowbreak{}i\allowbreak{}n\allowbreak{}\char95\relax\allowbreak{}i\allowbreak{}d\allowbreak{}x\allowbreak{}\char32\relax\allowbreak{}=\allowbreak{}\char32\relax\allowbreak{}N\allowbreak{}o\allowbreak{}n\allowbreak{}e} & 本轮暂时还没选中相邻位置。 \\[3pt]
\SourceLineNumber{16} & \BpeDocCodeCell{\char32\relax\allowbreak{}\char32\relax\allowbreak{}\char32\relax\allowbreak{}\char32\relax\allowbreak{}\char32\relax\allowbreak{}\char32\relax\allowbreak{}\char32\relax\allowbreak{}\char32\relax\allowbreak{}m\allowbreak{}i\allowbreak{}n\allowbreak{}\char95\relax\allowbreak{}r\allowbreak{}a\allowbreak{}n\allowbreak{}k\allowbreak{}\char32\relax\allowbreak{}=\allowbreak{}\char32\relax\allowbreak{}N\allowbreak{}o\allowbreak{}n\allowbreak{}e} & 本轮暂时还没找到最小 \BpeDocCodeInline{rank}。 \\[3pt]
\SourceLineNumber{17} & \BpeDocCodeCell{\char32\relax\allowbreak{}\char32\relax\allowbreak{}\char32\relax\allowbreak{}\char32\relax\allowbreak{}\char32\relax\allowbreak{}\char32\relax\allowbreak{}\char32\relax\allowbreak{}\char32\relax\allowbreak{}f\allowbreak{}o\allowbreak{}r\allowbreak{}\char32\relax\allowbreak{}i\allowbreak{},\allowbreak{}\char32\relax\allowbreak{}p\allowbreak{}a\allowbreak{}i\allowbreak{}r\allowbreak{}\char32\relax\allowbreak{}i\allowbreak{}n\allowbreak{}\char32\relax\allowbreak{}e\allowbreak{}n\allowbreak{}u\allowbreak{}m\allowbreak{}e\allowbreak{}r\allowbreak{}a\allowbreak{}t\allowbreak{}e\allowbreak{}(\allowbreak{}z\allowbreak{}i\allowbreak{}p\allowbreak{}(\allowbreak{}p\allowbreak{}a\allowbreak{}r\allowbreak{}t\allowbreak{}s\allowbreak{}[\allowbreak{}:\allowbreak{}-\allowbreak{}1\allowbreak{}]\allowbreak{},\allowbreak{}\char32\relax\allowbreak{}p\allowbreak{}a\allowbreak{}r\allowbreak{}t\allowbreak{}s\allowbreak{}[\allowbreak{}1\allowbreak{}:\allowbreak{}]\allowbreak{})\allowbreak{})\allowbreak{}:} & \BpeDocCodeInline{zip} 产生相邻片段对；\BpeDocCodeInline{enumerate} 同时给出每一对左端的位置 i。 \\[3pt]
\SourceLineNumber{18} & \BpeDocCodeCell{\char32\relax\allowbreak{}\char32\relax\allowbreak{}\char32\relax\allowbreak{}\char32\relax\allowbreak{}\char32\relax\allowbreak{}\char32\relax\allowbreak{}\char32\relax\allowbreak{}\char32\relax\allowbreak{}\char32\relax\allowbreak{}\char32\relax\allowbreak{}\char32\relax\allowbreak{}\char32\relax\allowbreak{}r\allowbreak{}a\allowbreak{}n\allowbreak{}k\allowbreak{}\char32\relax\allowbreak{}=\allowbreak{}\char32\relax\allowbreak{}m\allowbreak{}e\allowbreak{}r\allowbreak{}g\allowbreak{}e\allowbreak{}a\allowbreak{}b\allowbreak{}l\allowbreak{}e\allowbreak{}\char95\relax\allowbreak{}r\allowbreak{}a\allowbreak{}n\allowbreak{}k\allowbreak{}s\allowbreak{}.\allowbreak{}g\allowbreak{}e\allowbreak{}t\allowbreak{}(\allowbreak{}p\allowbreak{}a\allowbreak{}i\allowbreak{}r\allowbreak{}[\allowbreak{}0\allowbreak{}]\allowbreak{}\char32\relax\allowbreak{}+\allowbreak{}\char32\relax\allowbreak{}p\allowbreak{}a\allowbreak{}i\allowbreak{}r\allowbreak{}[\allowbreak{}1\allowbreak{}]\allowbreak{})} & 把两个片段的 \BpeDocCodeInline{bytes} 拼接，在已有词表中查它的 \BpeDocCodeInline{rank}；没有则得到 \BpeDocCodeInline{None}。 \\[3pt]
\SourceLineNumber{19} & \BpeDocCodeCell{\char32\relax\allowbreak{}\char32\relax\allowbreak{}\char32\relax\allowbreak{}\char32\relax\allowbreak{}\char32\relax\allowbreak{}\char32\relax\allowbreak{}\char32\relax\allowbreak{}\char32\relax\allowbreak{}\char32\relax\allowbreak{}\char32\relax\allowbreak{}\char32\relax\allowbreak{}\char32\relax\allowbreak{}i\allowbreak{}f\allowbreak{}\char32\relax\allowbreak{}r\allowbreak{}a\allowbreak{}n\allowbreak{}k\allowbreak{}\char32\relax\allowbreak{}i\allowbreak{}s\allowbreak{}\char32\relax\allowbreak{}n\allowbreak{}o\allowbreak{}t\allowbreak{}\char32\relax\allowbreak{}N\allowbreak{}o\allowbreak{}n\allowbreak{}e\allowbreak{}\char32\relax\allowbreak{}a\allowbreak{}n\allowbreak{}d\allowbreak{}\char32\relax\allowbreak{}(\allowbreak{}m\allowbreak{}i\allowbreak{}n\allowbreak{}\char95\relax\allowbreak{}r\allowbreak{}a\allowbreak{}n\allowbreak{}k\allowbreak{}\char32\relax\allowbreak{}i\allowbreak{}s\allowbreak{}\char32\relax\allowbreak{}N\allowbreak{}o\allowbreak{}n\allowbreak{}e\allowbreak{}\char32\relax\allowbreak{}o\allowbreak{}r\allowbreak{}\char32\relax\allowbreak{}r\allowbreak{}a\allowbreak{}n\allowbreak{}k\allowbreak{}\char32\relax\allowbreak{}<\allowbreak{}\char32\relax\allowbreak{}m\allowbreak{}i\allowbreak{}n\allowbreak{}\char95\relax\allowbreak{}r\allowbreak{}a\allowbreak{}n\allowbreak{}k\allowbreak{})\allowbreak{}:} & 只有词表存在该合并，且它比目前候选更早或还没有候选，才更新最佳选择。 \\[3pt]
\SourceLineNumber{20} & \BpeDocCodeCell{\char32\relax\allowbreak{}\char32\relax\allowbreak{}\char32\relax\allowbreak{}\char32\relax\allowbreak{}\char32\relax\allowbreak{}\char32\relax\allowbreak{}\char32\relax\allowbreak{}\char32\relax\allowbreak{}\char32\relax\allowbreak{}\char32\relax\allowbreak{}\char32\relax\allowbreak{}\char32\relax\allowbreak{}\char32\relax\allowbreak{}\char32\relax\allowbreak{}\char32\relax\allowbreak{}\char32\relax\allowbreak{}m\allowbreak{}i\allowbreak{}n\allowbreak{}\char95\relax\allowbreak{}i\allowbreak{}d\allowbreak{}x\allowbreak{}\char32\relax\allowbreak{}=\allowbreak{}\char32\relax\allowbreak{}i} & 记录当前最佳相邻对的左端位置。 \\[3pt]
\SourceLineNumber{21} & \BpeDocCodeCell{\char32\relax\allowbreak{}\char32\relax\allowbreak{}\char32\relax\allowbreak{}\char32\relax\allowbreak{}\char32\relax\allowbreak{}\char32\relax\allowbreak{}\char32\relax\allowbreak{}\char32\relax\allowbreak{}\char32\relax\allowbreak{}\char32\relax\allowbreak{}\char32\relax\allowbreak{}\char32\relax\allowbreak{}\char32\relax\allowbreak{}\char32\relax\allowbreak{}\char32\relax\allowbreak{}\char32\relax\allowbreak{}m\allowbreak{}i\allowbreak{}n\allowbreak{}\char95\relax\allowbreak{}r\allowbreak{}a\allowbreak{}n\allowbreak{}k\allowbreak{}\char32\relax\allowbreak{}=\allowbreak{}\char32\relax\allowbreak{}r\allowbreak{}a\allowbreak{}n\allowbreak{}k} & 记录当前最小 \BpeDocCodeInline{rank}。 \\[3pt]
\SourceLineNumber{22} & \BpeDocCodeCell{\char32\relax\allowbreak{}\char32\relax\allowbreak{}\char32\relax\allowbreak{}\char32\relax\allowbreak{}\char32\relax\allowbreak{}\char32\relax\allowbreak{}\char32\relax\allowbreak{}\char32\relax\allowbreak{}i\allowbreak{}f\allowbreak{}\char32\relax\allowbreak{}m\allowbreak{}i\allowbreak{}n\allowbreak{}\char95\relax\allowbreak{}r\allowbreak{}a\allowbreak{}n\allowbreak{}k\allowbreak{}\char32\relax\allowbreak{}i\allowbreak{}s\allowbreak{}\char32\relax\allowbreak{}N\allowbreak{}o\allowbreak{}n\allowbreak{}e\allowbreak{}\char32\relax\allowbreak{}o\allowbreak{}r\allowbreak{}\char32\relax\allowbreak{}(\allowbreak{}m\allowbreak{}a\allowbreak{}x\allowbreak{}\char95\relax\allowbreak{}r\allowbreak{}a\allowbreak{}n\allowbreak{}k\allowbreak{}\char32\relax\allowbreak{}i\allowbreak{}s\allowbreak{}\char32\relax\allowbreak{}n\allowbreak{}o\allowbreak{}t\allowbreak{}\char32\relax\allowbreak{}N\allowbreak{}o\allowbreak{}n\allowbreak{}e\allowbreak{}\char32\relax\allowbreak{}a\allowbreak{}n\allowbreak{}d\allowbreak{}\char32\relax\allowbreak{}m\allowbreak{}i\allowbreak{}n\allowbreak{}\char95\relax\allowbreak{}r\allowbreak{}a\allowbreak{}n\allowbreak{}k\allowbreak{}\char32\relax\allowbreak{}>\allowbreak{}=\allowbreak{}\char32\relax\allowbreak{}m\allowbreak{}a\allowbreak{}x\allowbreak{}\char95\relax\allowbreak{}r\allowbreak{}a\allowbreak{}n\allowbreak{}k\allowbreak{})\allowbreak{}:} & 若没有候选，或候选等级已达到当前 \BpeDocCodeInline{token} 的 \BpeDocCodeInline{max\char95\relax\allowbreak{}rank} 上界，就停止。比较用 >=，因此只允许严格更小的 \BpeDocCodeInline{rank}。 \\[3pt]
\SourceLineNumber{23} & \BpeDocCodeCell{\char32\relax\allowbreak{}\char32\relax\allowbreak{}\char32\relax\allowbreak{}\char32\relax\allowbreak{}\char32\relax\allowbreak{}\char32\relax\allowbreak{}\char32\relax\allowbreak{}\char32\relax\allowbreak{}\char32\relax\allowbreak{}\char32\relax\allowbreak{}\char32\relax\allowbreak{}\char32\relax\allowbreak{}b\allowbreak{}r\allowbreak{}e\allowbreak{}a\allowbreak{}k} & 结束循环，保留尚未合并的片段。 \\[3pt]
\SourceLineNumber{24} & \BpeDocCodeCell{\char32\relax\allowbreak{}\char32\relax\allowbreak{}\char32\relax\allowbreak{}\char32\relax\allowbreak{}\char32\relax\allowbreak{}\char32\relax\allowbreak{}\char32\relax\allowbreak{}\char32\relax\allowbreak{}a\allowbreak{}s\allowbreak{}s\allowbreak{}e\allowbreak{}r\allowbreak{}t\allowbreak{}\char32\relax\allowbreak{}m\allowbreak{}i\allowbreak{}n\allowbreak{}\char95\relax\allowbreak{}i\allowbreak{}d\allowbreak{}x\allowbreak{}\char32\relax\allowbreak{}i\allowbreak{}s\allowbreak{}\char32\relax\allowbreak{}n\allowbreak{}o\allowbreak{}t\allowbreak{}\char32\relax\allowbreak{}N\allowbreak{}o\allowbreak{}n\allowbreak{}e} & 正常继续合并时必须已找到位置；该断言检查内部逻辑。 \\[3pt]
\SourceLineNumber{25} & \BpeDocCodeCell{\char32\relax\allowbreak{}\char32\relax\allowbreak{}\char32\relax\allowbreak{}\char32\relax\allowbreak{}\char32\relax\allowbreak{}\char32\relax\allowbreak{}\char32\relax\allowbreak{}\char32\relax\allowbreak{}p\allowbreak{}a\allowbreak{}r\allowbreak{}t\allowbreak{}s\allowbreak{}\char32\relax\allowbreak{}=\allowbreak{}\char32\relax\allowbreak{}p\allowbreak{}a\allowbreak{}r\allowbreak{}t\allowbreak{}s\allowbreak{}[\allowbreak{}:\allowbreak{}m\allowbreak{}i\allowbreak{}n\allowbreak{}\char95\relax\allowbreak{}i\allowbreak{}d\allowbreak{}x\allowbreak{}]\allowbreak{}\char32\relax\allowbreak{}+\allowbreak{}\char32\relax\allowbreak{}[\allowbreak{}p\allowbreak{}a\allowbreak{}r\allowbreak{}t\allowbreak{}s\allowbreak{}[\allowbreak{}m\allowbreak{}i\allowbreak{}n\allowbreak{}\char95\relax\allowbreak{}i\allowbreak{}d\allowbreak{}x\allowbreak{}]\allowbreak{}\char32\relax\allowbreak{}+\allowbreak{}\char32\relax\allowbreak{}p\allowbreak{}a\allowbreak{}r\allowbreak{}t\allowbreak{}s\allowbreak{}[\allowbreak{}m\allowbreak{}i\allowbreak{}n\allowbreak{}\char95\relax\allowbreak{}i\allowbreak{}d\allowbreak{}x\allowbreak{}\char32\relax\allowbreak{}+\allowbreak{}\char32\relax\allowbreak{}1\allowbreak{}]\allowbreak{}]\allowbreak{}\char32\relax\allowbreak{}+\allowbreak{}\char32\relax\allowbreak{}p\allowbreak{}a\allowbreak{}r\allowbreak{}t\allowbreak{}s\allowbreak{}[\allowbreak{}m\allowbreak{}i\allowbreak{}n\allowbreak{}\char95\relax\allowbreak{}i\allowbreak{}d\allowbreak{}x\allowbreak{}\char32\relax\allowbreak{}+\allowbreak{}\char32\relax\allowbreak{}2\allowbreak{}:\allowbreak{}]} & 保留选中位置之前的片段，把选中的两块拼成一块，再接上其后的片段；这一轮只合并这个位置。 \\[3pt]
\SourceLineNumber{26} & \BpeDocCodeCell{\char32\relax\allowbreak{}\char32\relax\allowbreak{}\char32\relax\allowbreak{}\char32\relax\allowbreak{}r\allowbreak{}e\allowbreak{}t\allowbreak{}u\allowbreak{}r\allowbreak{}n\allowbreak{}\char32\relax\allowbreak{}p\allowbreak{}a\allowbreak{}r\allowbreak{}t\allowbreak{}s} & 返回剩余字节片段列表。\BpeDocCodeInline{recover\char95\relax\allowbreak{}merges} 希望这里返回两个子块。 \\[3pt]
\end{longtable}
\endgroup


不要把两个辅助函数混淆：\BpeDocCodeInline{base.\allowbreak{}merge} 接收整数编号，针对已经选定的一种 \BpeDocCodeInline{pair} 从左向右替换多个不重叠出现；\BpeDocCodeInline{gpt4.\allowbreak{}bpe} 接收字节片段，在每轮重新扫描候选后只合并选中的一个位置，并受到 \BpeDocCodeInline{max\char95\relax\allowbreak{}rank} 限制。两者服务于不同阶段。


\begin{quote}\footnotesize\setlength{\parskip}{0pt} MinBPE 固定提交：\BpeDocCodeInline{base.\allowbreak{}py}，公共统计、合并和词表结构。
\par\nobreak
\url{https://github.com/karpathy/minbpe/blob/1acefe89412b20245db5a22d2a02001e547dc602/minbpe/base.py#L13-L95}
\end{quote}


\begin{quote}\footnotesize\setlength{\parskip}{0pt} MinBPE 固定提交：\BpeDocCodeInline{basic.\allowbreak{}py}，基础字节级 BPE 训练、编码与解码。
\par\nobreak
\url{https://github.com/karpathy/minbpe/blob/1acefe89412b20245db5a22d2a02001e547dc602/minbpe/basic.py#L15-L74}
\end{quote}


\begin{quote}\footnotesize\setlength{\parskip}{0pt} MinBPE 固定提交：\BpeDocCodeInline{regex.\allowbreak{}py}，预切分、普通编码和特殊 \BpeDocCodeInline{token} 处理。
\par\nobreak
\url{https://github.com/karpathy/minbpe/blob/1acefe89412b20245db5a22d2a02001e547dc602/minbpe/regex.py#L12-L164}
\end{quote}


\begin{quote}\footnotesize\setlength{\parskip}{0pt} MinBPE 固定提交：\BpeDocCodeInline{gpt4.\allowbreak{}py}，\BpeDocCodeInline{cl100k\char95\relax\allowbreak{}base} 包装、合并恢复、字节编号重排与接口限制。
\par\nobreak
\url{https://github.com/karpathy/minbpe/blob/1acefe89412b20245db5a22d2a02001e547dc602/minbpe/gpt4.py#L11-L96}
\end{quote}


\end{document}
